\subsection{Mazur's terminal error: retaining the square-root scales}
\label{subsec:Mazur-terminal}

Mazur's version--2 manuscript \cite{Mazur2026v2} gives another
natural-density argument, with an explicit logarithmic clock. Its
Theorems~1.1--1.2 state the odd-return constant
$501501/(5000\log2)$, the raw constant $1509503/(5000\log2)$,
and every fixed-target exponent $0<d<5/143$. These are the source's
claims; the terminal scalar calculation below is not a certificate
for the whole source package.

The two restrictions in its Lemmas~3.4--3.5 have different origins.
Phase averaging requires $d<1/(2\kappa)$ when
$\|q\log_2 3\|\ge c q^{1-\kappa}$. Terminal reconstruction is
stated with the additional restriction $d<1/20$. The latter is
introduced when square-root scales are replaced by larger pure
powers. Keeping the original expressions improves this part of
the argument without changing its reference level, tube, or measures.

\paragraph{Literal parameters and source locators.}
Let $B$ be a positive integer, $L=\log B$, and fix $C>0$. Put
\[
 \tau=\frac1{100000},\quad q=\frac1{10\log2},\quad
 f=\frac{340}{100000},\quad
 a=\frac1{200000},\quad b=\frac1{50000},
\]
\[
 n=\lfloor qL\rfloor,\quad m_0=\lfloor\tau L\rfloor,\quad
 m=\lceil m_0L^{-2/5}\rceil,\quad
 S=\sqrt{m\log m},\quad W=C\sqrt{n\log n},\quad F=fL.
\]
Here $n$ is the local passage horizon, $m_0$ the step-back
depth, and $m$ the smaller common reference level. They are not
interchangeable. These are the unnumbered schedule following (19)
in Section~3.3 and the tube definition preceding (26) of the paper,
and the definitions in
\texttt{A5ReferenceSchedule.lean}, \texttt{A5Physical.lean}
and \texttt{A5ReferenceEta.lean}, in its
\texttt{Erdos1135/ND/Band} directory.
The executable definition $n=\lfloor\lfloor\log_2 B\rfloor/10\rfloor$
equals $\lfloor qL\rfloor$ by integer division.
The exact nonnegative scalar majorant in the last file is
\begin{equation}
 \eta_B(C)=\frac{57SW}{F}+\frac{1021S^2}{F}+\frac8{m^3}.
 \label{eq:Mazur-eta}
\end{equation}
All following assertions concern large $B$, where $m,n\ge1$.

\begin{proposition}[a sharper terminal majorant]
\label{prop:Mazur-eta}
For $L\ge10^{10}$ the literal parameters satisfy
\begin{equation}
 \eta_B(C)\le
 \frac{57C\sqrt{bq}}{f}L^{-1/5}\log L+
 \frac{1021b}{f}L^{-2/5}\log L+
 8a^{-3}L^{-9/5}.
 \label{eq:Mazur-eta-envelope}
\end{equation}
Moreover, with $C$ fixed and positive,
\begin{equation}
 \frac{\eta_B(C)}{L^{-1/5}\log L}
 \longrightarrow
 A_C:=\frac{57C}{f}\sqrt{\frac35\tau q}>0.
 \label{eq:Mazur-eta-limit}
\end{equation}
In particular, for every $d<1/5$ and $\varepsilon>0$,
$\eta_B(C)\le\varepsilon L^{-d}$ for all sufficiently large $B$.
\end{proposition}
\begin{proof}
The floor and ceiling give
\[
 \tau L-1<m_0\le\tau L,\qquad
 m_0L^{-2/5}\le m<m_0L^{-2/5}+1.
\]
Since $L\ge10^{10}$, these imply
$aL^{3/5}\le m\le bL^{3/5}$: for the lower bound use
$\tau L-1\ge\tau L/2$; for the upper bound use
$1\le\tau L^{3/5}$.
Also $1\le m\le L$ and $1\le n\le qL<L$.
The last inequality uses $\log2>1/2$, obtained by integrating
$1/t>1/2$ on $(1,2)$.
Monotonicity of the logarithm now gives
\[
 S^2\le bL^{3/5}\log L,\qquad
 W^2\le C^2qL\log L.
\]
All factors are nonnegative. Taking square roots, multiplying,
and dividing by the original $F=fL$ gives the first term of
\eqref{eq:Mazur-eta-envelope}. The same estimate for $S^2$
gives its second term. The lower bound for $m$ gives its third.
No rescaling of either law or tube is involved.

For the asserted limit, the same floor and ceiling bounds give
\[
 \frac m{L^{3/5}}\to\tau,\quad \frac nL\to q,\quad
 \frac{\log m}{\log L}\to\frac35,\quad
 \frac{\log n}{\log L}\to1.
\]
Consequently
\[
 \frac{SW}{L^{4/5}\log L}
 =C\sqrt{\frac m{L^{3/5}}\frac nL
              \frac{\log m}{\log L}\frac{\log n}{\log L}}
 \longrightarrow C\sqrt{\frac35\tau q}.
\]
After division by $L^{-1/5}\log L$, the second term in
\eqref{eq:Mazur-eta} is $O(L^{-1/5})$ and its third is
$O(L^{-8/5}/\log L)$; both tend to zero.
This proves \eqref{eq:Mazur-eta-limit}.
Finally, if $\delta=1/5-d>0$, then
$L^{-\delta}\log L\to0$. For example,
$\log L\le (2/\delta)L^{\delta/2}$ for $L\ge1$,
by integrating $t^{-1}\le t^{\delta/2-1}$.
This proves the final assertion.
\end{proof}

\begin{figure}[ht]
\centering
\fbox{\begin{minipage}{0.92\linewidth}
\[
 \begin{array}{c@{\qquad}c}
 \text{Retained scales}&\text{Source's pure-power bounds}\\[2pt]
 S=O(L^{3/10}\sqrt{\log L})&S\le L^{7/20}\\
 W=O_C(L^{1/2}\sqrt{\log L})&W\le L^{3/5}\\
 F=fL&F=fL\\[3pt]
 SW/F=O_C(L^{-1/5}\log L)&SW/F=O(L^{-1/20})
 \end{array}
\]
The leading coefficient of the actual scalar majorant is
$A_C=(57C/f)\sqrt{(3/5)\tau q}$, not an unspecified change of
normalization. The lower two terms in \eqref{eq:Mazur-eta}
decay faster.
\end{minipage}}
\caption{Where the terminal exponent is lost, and how it is recovered.
This is a proof diagram for Proposition~\ref{prop:Mazur-eta},
not sampled orbit data. The original parameters and coefficients are
Mazur's; the retained-scale comparison is the present refinement.}
\label{fig:Mazur-scales}
\end{figure}




\begin{proposition}[the complete scalar terminal bound]
\label{prop:Mazur-terminal-rate}
Fix $C>0$ and $K\ge0$. Define the two nonproportional errors
and terminal scalar bounds by the source's exact expressions
\[
 p_H(m,K)=20K/m^{11}+80/m^3,\qquad p_F(m,K)=4p_H(m,K),
\]
\[
 {\cal R}_\sigma(B)=\frac{86}{25}\eta_B(C)
                  (1+11B^{-9/10})+14W p_\sigma(m,K),
 \qquad \sigma\in\{H,F\}.
\]
Then
\[
 0\le{\cal R}_H(B)\le{\cal R}_F(B),\qquad
 \frac{{\cal R}_\sigma(B)}{L^{-1/5}\log L}
       \longrightarrow\frac{86}{25}A_C .
\]
For every $0<d<1/5$ and $\varepsilon>0$ both bounds are
at most $\varepsilon(\log B)^{-d}$ for all sufficiently large
$B$. The threshold may depend on $C,K,d,\varepsilon$, but not
on a source branch, a band index, an endpoint, or an output event.
\end{proposition}
\begin{proof}
The factor $86/25$ is twice the original absorption coefficient
$43/25$; it has not been discarded. Positivity and $p_H\le p_F$
give the first assertion. With the unchanged schedules,
\[
 W/m^{11}=O_C(L^{-61/10}\sqrt{\log L}),\qquad
 W/m^3=O_C(L^{-13/10}\sqrt{\log L}).
\]
Indeed $W=O_C(L^{1/2}\sqrt{\log L})$ and
$m\asymp L^{3/5}$, so the exponents are
$1/2-33/5=-61/10$ and $1/2-9/5=-13/10$.
Both quantities are $o(L^{-1/5}\log L)$.
Since $B^{-9/10}\to0$, Proposition~\ref{prop:Mazur-eta}
therefore proves the two limits, and its final elementary
logarithmic comparison proves the rate assertion.
All four parameters were fixed before taking $B$ large.
The displayed scalar expressions contain none of the later
branch, band, endpoint or event parameters, which proves
the stated uniformity.
\end{proof}

\paragraph{Propagation through the source estimates.}
The source theorem
\texttt{eventually\_ndA5ReferenceTerminalAggregateRHS\_le\_rate}
bounds exactly $\mathcal R_H,\mathcal R_F$ above, but only for
$d<1/20$. Proposition~\ref{prop:Mazur-terminal-rate} proves
that scalar assertion for every $d<1/5$.
Its next theorem compares each nominal terminal sum with the
physical reference profile by these scalar bounds; the fixed
constant $K$ is chosen before $B$ and before the band and event.
Thus the improved estimate remains uniform when used there.

The subsequent genuine-band comparison is a four-term triangle
inequality. Its coverage error is at most $L^{-10}$,
its exact-to-nominal error at most $11B^{-9/10}$,
its terminal aggregation error is bounded by
Proposition~\ref{prop:Mazur-terminal-rate}, and its analytic
full-profile error is at most $6L^{-d}$.
The first two errors are respectively at most
$L^{-d}/4$ and $L^{-d}/2$ for large $B$ and $d<1/5$;
take $\varepsilon=1/4$ for the third.
The same displayed sum in the source is therefore
$(1/4+1/2+1/4+6)L^{-d}=7L^{-d}$.
The full-profile estimate still requires the separate
phase restriction $d<1/(2\kappa)$.
Consequently the restriction introduced by this terminal
calculation is $d<1/5$, not $d<1/20$.
This deduction uses the source's finite terminal comparison
and full-profile estimate at their stated scopes; those entire
dependency chains have not been independently replayed here.

For the source's $\kappa=143/10$, the phase ceiling
$5/143<1/20$ already dominates, so this refinement does not
change its printed headline exponent. For the published Rhin
input $\kappa=8.616$ used earlier in this note, the phase ceiling
is $1/17.232$: the improved terminal bound no longer prevents
that exponent from passing through this part of Mazur's
argument. This is a precise compatibility result between the
two arguments, not a new claim that a $599$-file Lean package
has been checked locally.

Equation~\eqref{eq:Mazur-eta-limit} also shows that this
particular positive scalar majorant is not $O(L^{-1/5})$.
That statement is about the retained upper-bound expression.
It is not a lower bound on the actual passage-law error and
does not exclude a sharper argument using cancellation.

\paragraph{Version and formal status.}
The manuscript is version~2, 21 July 2026. It names source
commit \texttt{f386357d453ac4dcf91242b76252d88a5a729906};
the downloaded source and its evidence records name
\texttt{ca3dd0d63920411213403092aecc6946619eb082}.
The latter's main file hash agrees with its advertised hash.
The source evidence reports the standard axioms
\texttt{propext}, \texttt{Classical.choice}, \texttt{Quot.sound}
and no \texttt{sorryAx}. These are the publisher's records,
not a build performed in this workbench. The archive contains
the first-party Lean closure but no bundled Mathlib, lake
manifest, or toolchain file. No formal verification status
is transferred from that evidence to the refinement above.

\subsubsection{Finite terminal records, boundary mass, and sharp absorption}
\label{sec:Mazur-finite}

The terminal scalar has a specific probabilistic meaning. It measures the
comparison between a finite sum of nominal affine weights and a reference
profile. Before this comparison can be used, one must retain the actual
integer sources and show why summing over many terminal records does not
multiply a probability mass by the number of records. We reconstruct that
finite mechanism here, then sharpen its algebraic absorption step.
The source locators are Mazur~\cite{Mazur2026v2}, Appendix~C.2--C.3,
pp.~21--23, and the Lean files named below; all refer to the downloaded
commit already identified in Section~\ref{subsec:Mazur-terminal}.

\paragraph{The exact inverse, with compatibility retained.}
Write $T=\operatorname{Syr}$ on positive odd integers. Fix an integer
threshold $B\ge1$, a positive integer $h$, a finite set $J$ of integer
times greater than $h$, and a finite set $\mathcal E$ of positive odd
integers whose first hitting time of $[1,B]$ is exactly $h$.
For $t\in J$ put $\nu=t-h$. A terminal record consists of
$(t,a,M)$, where $a=(a_1,\ldots,a_\nu)$ has positive integer entries
and $M\in\mathcal E$. Restrict these records to any finite collection;
Mazur uses the additional strict total-valuation tube
$|\ell-2\nu|<W$, where $\ell=\sum_i a_i$.

Put $s_i=a_1+\cdots+a_i$, $s_0=0$, and
\[
 c_a=\sum_{i=0}^{\nu-1}3^{\nu-1-i}2^{s_i},\qquad
 F_a=c_a/2^\ell,\qquad
 N_a(M)=\frac{2^\ell M-c_a}{3^\nu},\qquad
 N_a^0(M)=\frac{2^\ell M}{3^\nu}.
\]
Compatibility means $2^\ell M\equiv c_a\pmod{3^\nu}$.
Retain only compatible records with $M>F_a$. Then $N_a(M)$ is a
positive odd integer: $c_a$ is odd, the numerator is positive and
odd, and its odd divisor is $3^\nu$.
This is the inverse of the original affine equation
$T^\nu N=(3^\nu N+c_a)/2^\ell=M$ on the indicated word cylinder,
not a change of the map.

For clarity, terminal oddness proves the intermediate exponents too.
The one-letter assertion is immediate from $3N+1=2^{a_1}M$.
For a longer word, use $c_a=3^{\nu-1}+2^{a_1}c_{a'}$ to obtain
$2^{a_1}\mid3N+1$. The resulting integer $y=(3N+1)/2^{a_1}$
is positive and odd: the suffix equation
$3^{\nu-1}y+c_{a'}=2^{\ell-a_1}M$ has odd $c_{a'}$ and even
right side. Induction on the suffix proves all remaining exponents.
Conversely actual iteration gives the affine equation.
This gives both directions and the unique inverse.
It is the arithmetic content of
\texttt{A5TerminalAtomDenominatorComparison.lean},
\texttt{ndA5TerminalAtomCandidate\_odd\_and\_affine};
the source's endpoint-room estimate supplies $M>F_a$.

\begin{proposition}[one source per terminal record, across times]
\label{prop:Mazur-atoms}
Suppose $\max J-\min J\le h$ when $J$ is nonempty.
The map from compatible positive terminal records to $N_a(M)$ is
injective. In particular, restricting to any physical sampling band
preserves disjointness, and boundary candidates outside that band
also have multiplicity at most one.
\end{proposition}
\begin{proof}
Suppose the same $N$ has records at $t\le t'$.
The affine decoding gives $M=T^{t-h}N$ and
$M'=T^{t'-h}N=T^{t'-t}M$. Put $d=t'-t$.
Since $0\le d\le h$ and the first hit from $M$ is at $h$,
the first hit from $T^dM$ is exactly $h-d$:
its earlier iterates are the corresponding earlier iterates of $M$,
and its iterate at $h-d$ is $T^hM\le B$.
But $M'\in\mathcal E$ has first-hit time $h$, so $d=0$.
The times and word lengths coincide. Deterministic iteration from $N$
then gives the same ordered word and the same endpoint.
No band-membership assumption entered the proof.
\end{proof}

For the source's particular band, with $d=\log(4/3)$,
\[
 J=\left[\left\lceil\frac{\log(z/B)}d-3W\right\rceil,
          \left\lfloor\frac{\log(U/B)}d+3W\right\rfloor\right]
          \cap\mathbb N
 \quad\Longrightarrow\quad
 \operatorname{diam}J\le\frac{\beta}{d}+6W .
\]
The floor and ceiling can only reduce this diameter.
With fixed $C$, $W=o(\log B)$ and
$h=m_0\sim(\log B)/100000$, so the diameter is at most $h$
eventually. Also $z\ge B^{1001/1000}$ gives
$\min J\ge (\log B)/(1000d)-3W>h$ eventually:
$d<1/3$ implies $1/(1000d)>3/1000>1/100000$.
Thus the time hypotheses hold for the unchanged bands and schedules,
not for an arbitrary long interval of times.

The source proof is
\texttt{A5TerminalAtoms.lean},
\texttt{ndA5TerminalAtomIndex\_eq\_of\_common\_affineSource}.
This is a statement about terminal records, not yet identification of
their union with the genuine first-passage event.
For example, the actual orbit of $63$ first reaches $[1,50]$ at
time $35$, with value $23$; its values at times $36,37,38$ are
$35,53,5$. Thus an endpoint at time $37$ whose own first hit is one
step later does not imply that the source's first hit is time $38$.
Also the endpoints at times $34$ and $37$ are $61$ and $53$,
both with first-hit time one; their corresponding terminal times differ
by three and therefore violate the proposition's diameter hypothesis.
These exact examples explain why the source has a separate coverage
argument and a short time-window condition.

\paragraph{Actual and nominal weights.}
Take a nonempty band
$\mathcal B=\{N>0:N\text{ odd},\ z\le N<U\}$, where
$z>0$, $U=ze^\beta$, and $\beta>0$.
Write $D=\#\mathcal B$, $H=\sum_{N\in\mathcal B}N^{-1}$.
For a compatible record put
\[
 \alpha_a=F_a/M,\quad
 \delta_a=-\log(1-\alpha_a),\quad
 u_a=\log(N_a^0(M)/z).
\]
Then $0<\alpha_a<1$ and the exact identities are
\[
 N_a=(1-\alpha_a)N_a^0,\qquad
 \log(N_a/z)=u_a-\delta_a,\qquad
 \frac{3^\nu}{2^\ell M}=\frac1{N_a^0}
                  =\frac{1-\alpha_a}{N_a}.
 \tag{*}
\]
Let $P_a=\mathbf1_{\{z\le N_a<U\}}$ and
$Q_a=\mathbf1_{\{z<N_a^0<U\}}$; their endpoint conventions differ
deliberately. The actual masses of the union and the nominal sums are
\[
 U_H=\sum_a\frac{P_a}{HN_a},\quad U_F=\sum_a\frac{P_a}{D},
 \qquad I_H=\sum_a\frac{Q_a}{HN_a^0},\quad I_F=\sum_a\frac{Q_a}{D}.
\]
Here and below the index $a$ in sums abbreviates the entire terminal
record; different times have not been identified.
Proposition~\ref{prop:Mazur-atoms} gives $0\le U_H,U_F\le1$.

\begin{proposition}[boundary comparison with no record-count loss]
\label{prop:Mazur-boundaries}
If all retained records have $\delta_a\le\varepsilon$ and
$\alpha_a\le\rho$, define the odd-integer shells
\[
 S_-=[ze^{-\varepsilon},z)\cap(2\mathbb N+1),\qquad
 S_+=[Ue^{-\varepsilon},U)\cap(2\mathbb N+1).
\]
Then
\[
 |U_F-I_F|\le\frac{\#(S_-\cup S_+)}D,\qquad
 |U_H-I_H|\le\rho U_H+
       \frac1H\sum_{N\in S_-\cup S_+}\frac1N.
 \label{eq:Mazur-shell-bounds}
\]
There are also the one-sided bounds
\[
 I_F\le U_F+\frac{\#S_-}{D},\qquad
 I_H\le U_H+\frac1H\sum_{N\in S_-}\frac1N.
\]
\end{proposition}
\begin{proof}
If $Q_a=1$ and $P_a=0$, then $0<u_a<\beta$ while
$u_a-\delta_a<0$; consequently $N_a\in S_-$.
If $P_a=1$ and $Q_a=0$, positivity of $\delta_a$ excludes
$u_a\le0$, so $\beta\le u_a<\beta+\delta_a$ and
$N_a\in S_+$. Injectivity charges each mismatching record to
a distinct shell integer. This proves the flat bound.
For the harmonic bound insert the actual-support nominal sum
$I_H^{\rm phys}=\sum_a P_a/(HN_a^0)$. Identity (*) gives
\[
 0\le U_H-I_H^{\rm phys}
       =\sum_a\frac{P_a\alpha_a}{HN_a}\le\rho U_H.
\]
For support mismatch, $1/N_a^0\le1/N_a$, so injectivity
bounds $|I_H^{\rm phys}-I_H|$ by the stated shell sum.
The triangle inequality proves the assertion.
For the one-sided bounds only nominal-only records can increase
the nominal sum; these lie in $S_-$, while on the common support
the nominal harmonic weight is no greater than the actual one.
\end{proof}

Mazur's finite band estimates supply
$\varepsilon=B^{-9/10}$, $\rho=B^{-19749/20000}$, and the
respective shell ratios $<11B^{-9/10}$ and $\le5B^{-9/10}$.
The precise source declarations are in
\texttt{A5TerminalBoundaryShell.lean} and
\texttt{A5TerminalAtomNominalization.lean}.
Thus both nominal sums have the common cap $1+11B^{-9/10}$;
for the harmonic sum the triangle estimate alone already gives
$1+6B^{-9/10}$. The one-sided statement above retains still more
information if the individual lower shell is to be used.
These estimates explain, rather than assume, why the cap is for
the entire time-indexed sum.

\paragraph{The two integer rows and the flat exponential.}
Let $\theta=\log_2 3$, $A=\log_2(z/M)$, and
\[
 b_\nu=A+\nu\theta,\quad
 \ell_{\nu,r}=\lfloor b_\nu\rfloor+1+r,\quad
 u_{\nu,r}=(\ell_{\nu,r}-b_\nu)\log2
             =(r+1-\{b_\nu\})\log2,\quad r=0,1.
\]
The $+1$ is retained even when $b_\nu$ is integral.
The source bands have $0<\beta<2\log2$ for $\log B\ge300000$:
their logarithmic length before subdivision is at least $300$,
so $\beta<301/300<2\log2$.

\begin{proposition}[exact two-row reindexing]
\label{prop:Mazur-two-rows}
For $\nu\ge1$ and $0<W\le\nu/8$, let
$q_{\nu,\ell}=\binom{\ell-1}{\nu-1}2^{-\ell}$ for $\ell\ge\nu$
and zero otherwise. For every real-valued function $f$ on integers,
\[
 \sum_{\substack{\nu\le\ell\le3\nu\\|\ell-2\nu|<W\\
              0<(\ell-b_\nu)\log2<\beta}}
      q_{\nu,\ell}f(\ell)
 =
 \sum_{r=0}^1
 \mathbf1_{\{|\ell_{\nu,r}-2\nu|<W\}}
 \mathbf1_{\{0<u_{\nu,r}<\beta\}}
 q_{\nu,\ell_{\nu,r}}f(\ell_{\nu,r}).
\]
\end{proposition}
\begin{proof}
The integers strictly above $b_\nu$ begin at
$\lfloor b_\nu\rfloor+1$. Since $\beta/\log2<2$, only this
integer and its successor can satisfy the displayed band inequality.
The map $r\mapsto\lfloor b_\nu\rfloor+1+r$ is injective, with
inverse $\ell\mapsto\ell-\lfloor b_\nu\rfloor-1$.
The tube and $W\le\nu/8$ imply $\nu<\ell<3\nu$, so no
additional endpoint is lost. Each summand is identical after this
bijection; neither positivity of $f$ nor an approximation is used.
The expression for $q_{\nu,\ell}$ follows by counting the
$\binom{\ell-1}{\nu-1}$ positive compositions, each of geometric
probability $2^{-\ell}$.
\end{proof}

This is the finite identity in
\texttt{A5TerminalAtomReconstruction.lean},
\texttt{sum\_ndA5TerminalTotals\_endpointMass\_strict\_eq\_range\_two}.
Before any mixing estimate, let $\rho_{\nu,\ell}(M)$ be the fraction
of those compositions satisfying the original compatibility congruence.
Summing the harmonic nominal weights over one fixed $(\nu,\ell,M)$
cell gives exactly
\[
 \frac{3^\nu}{MH}\,q_{\nu,\ell}\rho_{\nu,\ell}(M).
\]
The flat expression is the same before the harmonic normalizer,
multiplied instead by $z e^{(\ell-b_\nu)\log2}/D$.
Indeed $(3^\nu/(2^\ell M))ze^{(\ell-b_\nu)\log2}=1$.
Thus the exponential appears exactly once; omitting it would change
the flat law to a different weighting.

Replacing the conditioned residue factor by its common coarse law
is the separate mixing step. Its resulting endpoint kernel is
$K_m(M)=3^m p_m(M\bmod3^m)/M$, with $p_m$ the source's
Syracuse offset law at the fixed reference depth $m$.
The two reference profiles retain this same kernel and multiply it by
\[
 \frac1H\sum_{\nu,r}Q_{\nu,r}
 \quad\hbox{and}\quad
 \frac zD\sum_{\nu,r}Q_{\nu,r}e^{u_{\nu,r}},
\]
respectively, where $Q_{\nu,r}$ is the product of the two indicators
and $q_{\nu,\ell_{\nu,r}}$ above, and $\nu=t-h$, $t\in J$.
The inverse time reindex is $t=\nu+h$. These are the literal
profiles in \texttt{A5ReferenceTerminalAggregate.lean}, not an
identification of the two sampling measures.

\begin{lemma}[sharp finite proportional absorption]
\label{lem:Mazur-sharp-absorption}
Let a finite collection have $Y_i,Z_i,E_i\ge0$ and
$0\le a_i\le\lambda<1$, with
$|Y_i-Z_i|\le E_i+a_iZ_i$ for each $i$. Then
\[
 \sum_i|Y_i-Z_i|
 \le\frac{\lambda\sum_iY_i+\sum_iE_i}{1-\lambda}.
 \label{eq:Mazur-sharp-absorption}
\]
In particular this bounds the absolute difference of the two sums.
The coefficients cannot be decreased using only these hypotheses.
\end{lemma}
\begin{proof}
Put $D_*=\sum_i|Y_i-Z_i|$. Summation and $Z_i\ge0$ give
$D_*\le\sum_iE_i+\lambda\sum_iZ_i$.
Also $\sum_iZ_i\le\sum_iY_i+D_*$.
Hence $(1-\lambda)D_*\le\sum_iE_i+\lambda\sum_iY_i$;
division is permitted since $1-\lambda>0$.
For one cell, choose any $Y,E\ge0$ and $Z=(Y+E)/(1-\lambda)$.
Then $Z\ge Y$ and $Z-Y=E+\lambda Z$, so equality holds.
This proves sharpness for this finite comparison problem, not
saturation by the Collatz distributions.
\end{proof}

In the source cell inequality,
$i=(t,\ell)$, $Z_i=q_{\nu,\ell}R_i$ and
$E_i=q_{\nu,\ell}p_\sigma(m,K)$, where $R_i$ is the nonnegative
coarse-reference expectation. The coefficient is
\[
 a_i=\frac{43}{25}
   \left(\operatorname{eps}_{m,\nu,\ell}+\frac8{m^3}\right),
 \qquad
 \lambda_B=\frac{43}{25}\eta_B(C).
\]
Here $\operatorname{eps}_{m,\nu,\ell}$ denotes exactly
\texttt{ndFixedTotalRatioEps} in the source, not a new choice.
The common terminal packet gives $0\le a_i\le\lambda_B\le1/2$.
Since $\sum_\ell q_{\nu,\ell}\le1$, summing the errors first over
totals and then over times costs at most $\#J\,p_\sigma$.
The nominal mass cap, on the other hand, is applied once after
both sums. With $\#J\le7W$, Lemma~\ref{lem:Mazur-sharp-absorption}
therefore replaces the source's bound by the smaller explicit scalar
\[
 \mathcal R_\sigma^*(B)=
 \frac{\lambda_B(1+11B^{-9/10})+7W p_\sigma(m,K)}
      {1-\lambda_B}.
 \label{eq:Mazur-sharp-terminal}
\]
This is an algebraic refinement of the local-cell estimate
used in \texttt{A5ReferenceFixedTimeAbsorption.lean}, followed
by the same finite aggregation. It does not purport to prove
the analytic conditioned-mixing estimate used to obtain those cells.

\begin{corollary}[leading constant of the refined terminal scalar]
\label{cor:Mazur-sharp-terminal}
For the unchanged schedules, fixed $C>0$ and $K\ge0$, eventually
$0\le\lambda_B<1$, and
\[
 \frac{\mathcal R_\sigma^*(B)}{(\log B)^{-1/5}\log\log B}
       \longrightarrow\frac{43}{25}A_C,\qquad
 \frac{\mathcal R_\sigma^*(B)}{\mathcal R_\sigma(B)}
       \longrightarrow\frac12,\qquad \sigma=H,F.
\]
For $0\le\lambda_B\le1/2$, $\mathcal R_\sigma^*\le\mathcal R_\sigma$.
\end{corollary}
\begin{proof}
Proposition~\ref{prop:Mazur-eta} gives $\lambda_B\to0$ and its
scaled limit $(43/25)A_C$. Proposition~\ref{prop:Mazur-terminal-rate}
proves that $Wp_\sigma$ is negligible on that scale.
Thus the denominator tends to one, proving the first limit.
The exact relation is
$\mathcal R_\sigma^*=\mathcal R_\sigma/[2(1-\lambda_B)]$;
the second limit and the finite comparison follow.
\end{proof}

The admissible terminal exponents remain every $d<1/5$.
In the four-term genuine-band estimate of the preceding section,
$\mathcal R_\sigma$ may be replaced by $\mathcal R_\sigma^*$;
all other terms and the separate phase restriction remain unchanged.
The finite inverse, boundary comparison and two-row identities above
have written proofs and exact arithmetic regression tests.
The coverage and conditioned-mixing dependencies are reconstructed
separately below; Section~\ref{sec:Mazur-genuine-assembly} now composes them
with this finite mechanism. No finite test certifies those analytic proofs
or supplies a formal replay of the source package.

\begin{figure}[ht]
\centering
\fbox{\begin{minipage}{0.92\linewidth}
\[
 \begin{array}{ccc}
 (t,a,M)&\longmapsto&N=(2^\ell M-c_a)/3^\nu\\
 \text{first hit from }M\text{ at }h
 &&\text{distinct records have distinct }N\\[2pt]
 N&=&(1-F_a/M)N^0<N^0\\
 Q=1,\ P=0&\Longrightarrow&ze^{-\varepsilon}\le N<z\\
 P=1,\ Q=0&\Longrightarrow&Ue^{-\varepsilon}\le N<U
 \end{array}
\]
The injectivity hypothesis is $\operatorname{diam}J\le h$.
Shell mass is counted once, not once per time or valuation total.
\[
 \sum_i|Y_i-Z_i|
 \le\frac{\lambda\sum_iY_i+\sum_iE_i}{1-\lambda},
 \quad
 \sum_iY_i\le1+11B^{-9/10},\quad
 \sum_iE_i\le7W p_\sigma .
\]
\end{minipage}}
\caption{Exact terminal records and the two different aggregation costs.
Propositions~\ref{prop:Mazur-atoms}--\ref{prop:Mazur-two-rows}
reconstruct Mazur's finite mechanism; Lemma~\ref{lem:Mazur-sharp-absorption}
retains the denominator in its final comparison. The diagram shows
formulas and set membership, not sampled orbit geometry.}
\label{fig:Mazur-finite}
\end{figure}

\subsection{Fixed-total descent with fibre cancellation}
\label{sec:Mazur-descent}

The terminal comparison needs more than equidistribution of a residue.
Its endpoint weights depend on the valuation total, several endpoints
may occupy the same residue, and the reference law has forgotten that
total. These are three distinct operations. The following reconstruction
keeps their maps separate and improves the error by using the signs of
the discarded masses. The source is Mazur~\cite{Mazur2026v2},
\texttt{ND/Fourier/FixedTotalFiniteResolvedDescent.lean},
\texttt{FixedTotalSourceLaw.lean}, \texttt{FixedTotalRatioWindow.lean},
\texttt{FixedTotalRatioNormalization.lean}, and
\texttt{ND/Band/A5FixedCellTerminalDescent.lean}, in the pinned source
release specified above.

\paragraph{Masses and tests have different projection maps.}
Write $G_j=\mathbb Z/3^j\mathbb Z$, let $0\le m\le n$, and put
$k=3^{n-m}$. The reduction $\pi:G_n\to G_m$ has the fibre
$\{x+3^mj:0\le j<k\}$ over the standard representative $x$.
For real masses $p$ on $G_n$, masses $q$ on $G_m$, and real tests
$c$ on $G_n$, define
\begin{equation}
 (\pi_*p)(x)=\sum_{\pi y=x}p(y),\qquad
 (Uq)(y)=k^{-1}q(\pi y),\qquad
 (Ac)(x)=k^{-1}\sum_{\pi y=x}c(y).
 \label{eq:Mazur-descent-maps}
\end{equation}
Here $\pi_*$ sums masses; $A$ averages tests. With the counting pairing,
\[
 \pi_*U=I,\qquad \langle Uq,c\rangle=\langle q,Ac\rangle.
\]
The map $U$ is an $\ell^1$ isometry, $\pi_*$ an $\ell^1$ contraction,
and $A$ a positive unital supremum-norm contraction.
All statements follow by summing over the displayed $k$-element fibres.
The pullback $\pi^*g=g\circ\pi$ satisfies $A\pi^*=I$.
The kernel of $A$ consists exactly of tests whose sum on every fibre
is zero. Thus $A$ is onto, and loses precisely these within-fibre
variations, not the total weight assigned to a fibre.

For a finite endpoint index set $I$, arbitrary addresses $a_i\in G_n$
and coefficients $d_i$, put $c(y)=\sum_{a_i=y}d_i$. Finite interchange gives
\[
 \langle p,c\rangle=\sum_{i\in I}p(a_i)d_i,\qquad
 (Ac)(x)=k^{-1}\sum_{\pi a_i=x}d_i.
\]
Coincident addresses add their weights. No injectivity or choice of one
endpoint per residue is permitted. In the terminal cell the harmonic
coefficient is $d_M=3^n/(MH)$ and the flat coefficient is
$d_M=z3^n e^{u_{n,\ell}(M)}/(MD)$, with the strict endpoint set
and $u_{n,\ell}(M)=(\ell-\log_2(z/M)-n\log_2 3)\log2$
as in the preceding section. Averaging changes $3^n$ to $3^m$;
it changes neither that endpoint set nor the flat exponential.

\begin{proposition}[cancellation on each residue fibre]
\label{prop:Mazur-centred-test}
For a probability mass $p$ on $G_n$, put $q=\pi_*p$, $r=p-Uq$ and
$\Omega=\|r\|_1$. For every real test $c$,
\begin{equation}
 |\langle p,c\rangle-\langle q,Ac\rangle|
 \le \frac12\sum_{x\in G_m}
       \operatorname{osc}_{\pi^{-1}x}(c)
       \sum_{\pi y=x}|r(y)|.
 \label{eq:Mazur-centred-test}
\end{equation}
In particular, $0\le c\le H$ gives the bound $H\Omega/2$.
The coefficient $1/2$ is best possible in this general statement.
\end{proposition}
\begin{proof}
The sum of $r$ on each fibre is zero, since $\pi_*Uq=q$.
On that fibre subtract the midpoint between the largest and smallest
values of $c$ without changing its pairing with $r$.
The absolute value of the resulting test is at most half its oscillation.
The triangle inequality, summed over fibres, proves the estimate.
For sharpness take $m=0,n=1$, $p=(1,0,0)$ and $c=(H,0,0)$,
with $H>0$. Then $Uq=(1/3,1/3,1/3)$, $\Omega=4/3$, and the
pairing difference is $2H/3$. Equality holds.
\end{proof}

This uses the full $\ell^1$ distance, as in the source, not a redefinition
of that distance as total variation. The source's $H\Omega$ bound is
valid; the zero fibre sums give the smaller displayed bound.

\paragraph{The exact law after conditioning.}
Let $a_i$ be independent positive integers with
$\mathbb P(a_i=j)=2^{-j}$, and let $S_j=a_1+\cdots+a_j$.
For $j\ge1$ set
\[
 b_j(u)=
 \begin{cases}\binom{u-1}{j-1}2^{-u}&u\ge j,\\0&u<j,\end{cases}
 \qquad
 X_j=\sum_{i=1}^j3^{i-1}2^{-S_i}\pmod{3^j}.
\]
The inverses of $2$ here exist in $G_j$.
Write $\mu_{j,u}=\operatorname{Law}(X_j\mid S_j=u)$ when $u\ge j$.

\begin{proposition}[conditioned prefix mixture and its likelihood]
\label{prop:Mazur-prefix-mixture}
For $1\le m<n$ and $\ell\ge n$, define probabilities on the integers by
\[
 P(u)=\frac{b_m(u)b_{n-m}(\ell-u)}{b_n(\ell)},\qquad Q(u)=b_m(u).
\]
Then
\begin{equation}
 \pi_*\mu_{n,\ell}=\sum_uP(u)\mu_{m,u},\qquad
 \operatorname{Law}(X_m)=\sum_uQ(u)\mu_{m,u}.
 \label{eq:Mazur-prefix-mixture}
\end{equation}
Zero-weight terms with $u<m$ can use any probability law.
On $u\ge m$, the likelihood is
$R(u)=P(u)/Q(u)=b_{n-m}(\ell-u)/b_n(\ell)$.
For $u+(n-m)<\ell$ its positive adjacent ratio is exactly
\begin{equation}
 \frac{R(u+1)}{R(u)}
 =\frac{2(\ell-u-(n-m))}{\ell-u-1},\qquad
 \frac{R(u+1)}{R(u)}-1
 =\frac{(\ell-2n)-(u-2m)+1}{\ell-u-1}.
 \label{eq:Mazur-adjacent-ratio}
\end{equation}
\end{proposition}
\begin{proof}
Every positive composition of $\ell$ into $n$ parts has probability
$2^{-\ell}$; there are $\binom{\ell-1}{n-1}$ such compositions.
Splitting after $m$ parts is a bijection onto the disjoint union over $u$
of a composition of $u$ into $m$ parts and a composition of $\ell-u$
into $n-m$ parts. Its inverse concatenates the two ordered lists.
Counting proves the formula for $P$ and $\sum_uP(u)=1$.
For fixed $u$, the prefix is uniform among its positive compositions.
Reduction modulo $3^m$ discards exactly the terms with $i>m$ in $X_n$,
so it equals $X_m$ of that prefix. This proves the first mixture.
The second is ordinary partition by $S_m$, with an absolutely
convergent sum when paired with a bounded test.
Finally
$b_j(v-1)/b_j(v)=2(v-j)/(v-1)$ for $v>j$.
Use $j=n-m$, $v=\ell-u$ and subtract one. The strict suffix room
ensures positive numerator and denominator; the $+1$ is not omitted.
\end{proof}

The affine endpoint offset uses the reversed word: if
$F_a=\sum_{i=1}^n3^{n-i}2^{-(a_i+\cdots+a_n)}$, then
$F_{a^{\rm rev}}=X_n(a)$. Word reversal is its own inverse and preserves
the geometric probability and total valuation. It therefore identifies
the two fixed-total offset laws. It does not preserve arbitrary prefix
events. This is exactly the separately stated reversal map in
\texttt{ndAffineFixedTotalSubmass\_eq\_section7}, not an identification
of an unreversed prefix with an affine tail.

\begin{lemma}[likelihood comparison retaining the actual masses]
\label{lem:Mazur-exact-likelihood}
Let $P,Q$ be probabilities on a countable set and $G$ a set on which
both are positive. Put $p_0=P(G^c)<1$, $q_0=Q(G^c)<1$.
Suppose the oscillation of $\log(P/Q)$ on $G$ is at most $\epsilon\ge0$.
For every $u\in G$,
\begin{equation}
 e^{-\epsilon}\frac{1-p_0}{1-q_0}
 \le \frac{P(u)}{Q(u)}
 \le e^\epsilon\frac{1-p_0}{1-q_0}.
 \label{eq:Mazur-exact-likelihood}
\end{equation}
If $p_0\le\theta_P<1$ and $q_0\le\theta_Q<1$, set
\[
 a=\max\left(\frac{e^\epsilon}{1-\theta_Q}-1,\
              1-e^{-\epsilon}(1-\theta_P)\right).
\]
For every $0\le\phi\le H$,
\begin{equation}
 \left|\sum_u(P(u)-Q(u))\phi(u)\right|
 \le a\,\mathbb E_Q\phi+H\max(\theta_P,\theta_Q).
 \label{eq:Mazur-signed-tail}
\end{equation}
When $\theta_P=\theta_Q=\theta$, one can take
$a=e^\epsilon/(1-\theta)-1$.
\end{lemma}
\begin{proof}
For fixed $u\in G$, exponentiation compares $R(v)=P(v)/Q(v)$
to $R(u)$ within the factors $e^{\pm\epsilon}$ at every $v\in G$.
Average over $Q(v)/Q(G)$. The exact average is $P(G)/Q(G)$.
Rearranging proves the first inequality and then $|R(u)-1|\le a$.
The central contribution is bounded by
$a\mathbb E_Q(\phi\mathbf1_G)\le a\mathbb E_Q\phi$.
The outside signed contribution lies in $[-Hq_0,Hp_0]$.
Its absolute value is at most $H\max(p_0,q_0)$, not the sum
of these two bounds. All sums converge absolutely because $\phi$
is bounded. For equal tails put $v=e^\epsilon/(1-\theta)\ge1$;
$v-1\ge1-v^{-1}$ since $(v-1)^2/v\ge0$.
\end{proof}

\paragraph{The source's literal central window.}
Let $m\ge150$, $8m\le n$, $D_\ell=|\ell-2n|\le n/8$, and put
\[
 S_m=\sqrt{m\log m},\quad R_m=\lceil16S_m\rceil,\quad
 G=\{u\ge m:|u-2m|\le R_m\},
\]
\[
 \epsilon_{m,n,\ell}=\frac{57S_mD_\ell+1021S_m^2}{n}.
\]
This is precisely \texttt{ndFixedTotalRatioEps}: its first term is
$57\sqrt{\log m}\,(mD_\ell/n)/\sqrt m=57S_mD_\ell/n$.

For completeness the window constants follow from
$\log m\le m/29$ and $S_m\ge2$. The function $\log x/x$ decreases
for $x\ge150$, and
$\log150=7\log2+\log(75/64)<49/10+11/64<51/10<150/29$.
Here $\log2<7/10$ follows from the first five terms of the exponential
series at $7/10$, whose sum exceeds $2$.
Thus $S_m\le m/\sqrt{29}$ and
$16m/\sqrt{29}\le3m-1$ for $m\ge150$ (square the positive sides).
Consequently
\[
 R_m\le3m,\qquad R_m+1\le18S_m,\qquad 2R_m+1\le34S_m.
\]
Every $u\in G$ satisfies $u\le5m$ and $u+(n-m)<\ell$.
The denominator in \eqref{eq:Mazur-adjacent-ratio} is at least
$15n/8-5m-1\ge5n/4-1\ge6n/5$.
Its numerator after subtracting one has absolute value at most
$D_\ell+R_m+1\le n/2+1\le27n/50$, since $n\ge1200$.
The ratio is therefore within $9/20$ of one.
Integration of $1/t$ on $[11/20,29/20]$ gives
$|\log t|\le2|t-1|$ there. Adjacent log differences are at most
$5(D_\ell+R_m+1)/(3n)$.
For $u<v$ in $G$, telescope over the half-open integer interval
$[u,v)$, all of whose indices remain in $G$. There are at most $34S_m$
terms, so
\[
 |\log R(u)-\log R(v)|
 \le\frac{170S_m(D_\ell+18S_m)}{3n}
 \le\epsilon_{m,n,\ell}.
\]
This reconstructs the source's pairwise likelihood estimate with its
literal signs, ceiling, and $57,1021$ constants.

The source tail estimates used at this point are
\begin{equation}
 P(G^c)\le2/m^3,\qquad Q(G^c)\le2/m^3
 \quad\text{when }mD_\ell/n\le S_m.
 \label{eq:Mazur-tail-input}
\end{equation}
Their exact declarations are
\texttt{ndFixedTotalHeadTotalPMF\_outside\_ratioGood\_le\_two\_div\_cube}
and \texttt{ndGeom2EndpointPMF\_outside\_ratioGood\_le\_two\_div\_cube}
in \texttt{FixedTotalRatioTails.lean}.
The conditioned proof applies its concentration estimate about $m\ell/n$
at radius $15S_m$, using the displacement bound $mD_\ell/n\le S_m$,
and obtains $2e^{-225S_m^2/(63m)}\le2/m^3$.
The unconditioned proof obtains
$2e^{-\min(R_m^2/(32m),R_m/8)}\le2/m^3$.
These are Mazur's source estimates. Complete written reconstructions,
including their concentration inputs, now appear in
Corollary~\ref{cor:Mazur-two-tails} below. The finite checks do not
certify the analytic inequalities or replay their Lean formalization.

\begin{corollary}[refined fixed-cell descent]
\label{cor:Mazur-centred-descent}
Use the preceding window and displacement hypotheses. For any test
$0\le c\le H$ on $G_n$, put
\[
 \Omega=\|\mu_{n,\ell}-U\pi_*\mu_{n,\ell}\|_1,\qquad
 V=\mathbb E_{\operatorname{Law}(X_m)}Ac,\qquad
 a^\dagger=\frac{\exp(\epsilon_{m,n,\ell})}{1-2/m^3}-1 .
\]
Using the source tail estimates \eqref{eq:Mazur-tail-input},
\begin{equation}
 |\mathbb E_{\mu_{n,\ell}}c-V|
 \le H\Omega/2+a^\dagger V+2H/m^3.
 \label{eq:Mazur-centred-descent}
\end{equation}
\end{corollary}
\begin{proof}
The first replacement, from $\mu_{n,\ell}$ to $U\pi_*\mu_{n,\ell}$,
costs at most $H\Omega/2$ by Proposition~\ref{prop:Mazur-centred-test}.
By Proposition~\ref{prop:Mazur-prefix-mixture}, the remaining difference
is $\sum_u(P(u)-Q(u))\phi(u)$, where
$\phi(u)=\mathbb E_{\mu_{m,u}}Ac$ lies in $[0,H]$.
Its $Q$-expectation is $V$. Apply
Lemma~\ref{lem:Mazur-exact-likelihood} with $\theta=2/m^3$.
\end{proof}

At a terminal cell the source bounds the assembled tests by $20$
(harmonic) and $80$ (flat). The harmonic bound uses the endpoint
coefficient bound $5$ and band harmonic mass at least $1/4$;
the flat bound uses pointwise domination by four times the harmonic
test. Their original guards include the source passage-window facts,
positive band count, $m_0+n\le n_0$, and, for the flat case,
$\log B\ge300000$. We retain these guards and do not substitute a
single-endpoint estimate for the assembled-test bound.
With $\Omega\le K/m^{11}$ the nonproportional errors become
\[
 p_H^\dagger=10K/m^{11}+40/m^3=p_H/2,\qquad
 p_F^\dagger=40K/m^{11}+160/m^3=p_F/2.
\]
The two halves have different reasons: zero fibre mass for the first
term, and the sign of the two outside contributions for the second.
This does not strengthen the Fourier theorem itself.

The mixing input on the terminal schedule is also covered by
Proposition~\ref{prop:Tao-joint-total-mixing}: the fine level $n$ lies between
$f\log B$ and $n_0=\lfloor q\log B\rfloor$, while
$m\sim\tau(\log B)^{3/5}$. Hence $m\ge n^{1/2}$ eventually,
and $|\ell-2n|\le W=C\sqrt{n_0\log n_0}\ll_C\sqrt{n\log n}$
uniformly over these cells. Taking $A=11$ there gives a single finite
$K$ for this range. This connects the previously reconstructed Tao-based
estimate to the present descent. Also $8m=o(F)$, $W=o(F)$, and
$mW/(FS)=O_C((\log B)^{-1/5})\to0$; hence the central-window
and displacement guards hold uniformly eventually. It makes no assertion about every
smaller reference scale in the source's formal package.

\begin{corollary}[retaining the exponential improves the terminal constant]
\label{cor:Mazur-exponential-terminal}
Keep the literal schedules of \eqref{eq:Mazur-eta}. Define
\begin{equation}
 \begin{gathered}
 E_B=\frac{57SW+1021S^2}{F},\quad
 \theta_B=\frac2{m^3},\quad
 \gamma_B=\frac{e^{E_B}}{1-\theta_B}-1,\\
 \mathcal R_\sigma^\dagger(B)=
 \frac{\gamma_B(1+11B^{-9/10})+(7/2)W p_\sigma(m,K)}
       {1-\gamma_B}.
 \end{gathered}
 \label{eq:Mazur-exponential-terminal}
\end{equation}
For all sufficiently large $B$, this is the terminal comparison bound
obtained from the preceding descent and the same nominal cap.
On the source range $0\le\lambda_B=(43/25)\eta_B\le1/2$,
\[
 0\le\gamma_B\le\lambda_B,\qquad
 \mathcal R_\sigma^\dagger\le\mathcal R_\sigma^*\le\mathcal R_\sigma.
\]
For fixed $C>0$, $K\ge0$, and $\sigma=H,F$,
\[
 \frac{\mathcal R_\sigma^\dagger(B)}
      {(\log B)^{-1/5}\log\log B}\longrightarrow A_C,\qquad
 \frac{\mathcal R_\sigma^\dagger}{\mathcal R_\sigma}
       \longrightarrow\frac{25}{86},\qquad
 \frac{\mathcal R_\sigma^\dagger}{\mathcal R_\sigma^*}
       \longrightarrow\frac{25}{43}.
\]
\end{corollary}
\begin{proof}
On each cell, $n\ge F$ and $D_\ell\le W$ give
$\epsilon_{m,n,\ell}\le E_B$. Therefore $a^\dagger\le\gamma_B$.
Multiplication by the single endpoint mass $q_{n,\ell}$
gives a cell error with absolute part $q_{n,\ell}p_\sigma/2$.
Sum first over $\ell$, using $\sum_\ell q_{n,\ell}\le1$, and
then over times, using $\#J\le7W$.
Apply Lemma~\ref{lem:Mazur-sharp-absorption} and the nominal mass
cap $1+11B^{-9/10}$ once. This gives the displayed quotient.

Since $\eta_B=E_B+4\theta_B$, the stated range implies
$0\le\theta_B<1/2$ and $E_B+2\theta_B\le1$.
Integration gives $-\log(1-\theta_B)\le2\theta_B$.
For $0\le x\le1$, convexity gives
$e^x-1\le(e-1)x<(43/25)x$. For instance the exponential series
through degree $6$ has remaining tail at most
$(1/7!)/(1-1/8)$, proving $e<68/25$.
Consequently
\[
 \gamma_B
 \le e^{E_B+2\theta_B}-1
 \le(43/25)(E_B+2\theta_B)\le\lambda_B.
\]
The function $(aM+e)/(1-a)$ increases in $a,e\ge0$ for $a<1$,
$M\ge0$, proving the finite comparison with $\mathcal R_\sigma^*$.

Write $s_B=(\log B)^{-1/5}\log\log B$. The earlier scalar proof gives
$E_B/s_B\to A_C$, $\theta_B/s_B\to0$ and $Wp_\sigma/s_B\to0$.
Since $E_B\to0$,
$\gamma_B=(e^{E_B}-1+\theta_B)/(1-\theta_B)$ has
$\gamma_B/s_B\to A_C$. The denominator and nominal cap tend to one.
The last two limits follow from the previously proved coefficients
$(86/25)A_C$ and $(43/25)A_C$, both positive.
\end{proof}

The four-term genuine-band estimate can therefore use
$\mathcal R_\sigma^\dagger$ in place of $\mathcal R_\sigma^*$.
Its boundary, exceptional-event and phase terms are unchanged.
The terminal exponent range is still $d<1/5$; the source's separate
phase restriction is not removed by this constant improvement.
The tail estimates and genuine first-passage identification are proved
separately in Section~\ref{sec:Mazur-concentration-coverage}; their
dependencies remain distinct from the finite terminal comparison.

\begin{figure}[ht]
\centering
\fbox{\begin{minipage}{0.92\linewidth}
\[
 \begin{array}{ccccc}
 \mu_{n,\ell}&\longrightarrow&U\pi_*\mu_{n,\ell}
 &\longrightarrow&U\operatorname{Law}(X_m)\\
 \substack{\text{fine conditioned}\\\text{law}}&&
 \substack{\text{coarse conditioned}\\\text{law lifted}}&&
 \substack{\text{unconditioned}\\\text{reference lifted}}\\
 &H\Omega/2&&a^\dagger V+2H/m^3&
 \end{array}
\]
Each arrow records its absolute pairing error against $0\le c\le H$.
The first uses zero mass on every residue fibre. The second uses the
exact prefix-mixture likelihood and retains the signed outside mass.
Endpoint collisions are summed before either comparison.
\end{minipage}}
\caption{Two different reasons for error in the fixed-cell comparison.
The arrows are comparisons of expectations, not deterministic maps
between sample outcomes. The actual mass and test maps are
\eqref{eq:Mazur-descent-maps}; the error is proved in
Corollary~\ref{cor:Mazur-centred-descent}, with the tail bounds proved in
Corollary~\ref{cor:Mazur-two-tails}.
Mazur's source construction is retained with finer error accounting.}
\label{fig:Mazur-descent}
\end{figure}

\subsection{Concentration from cuts and the actual passage event}
\label{sec:Mazur-concentration-coverage}

There are two remaining identifications behind the preceding comparison.
First, conditioning on a total valuation changes the distribution of an
initial block; its exceptional mass must be proved small.
Second, a terminal affine record need not by itself describe the
\emph{first} passage of its source. This section reconstructs both
mechanisms. The first uses finite counting and exponential moments.
The second uses the actual ordered valuation word and a lower bound on
the step-back endpoint. Neither is inferred from a numerical orbit survey.

The primary source locators are Mazur~\cite{Mazur2026v2},
\texttt{ND/Fourier/FixedTotalBars.lean},
\texttt{FixedTotalConcentration.lean}, \texttt{FixedTotalPGF.lean},
\texttt{FixedTotalMGF.lean}, \texttt{FixedTotalTail.lean},
\texttt{FixedTotalT3.lean}; and
\texttt{ND/Band/A5TerminalAtomCoverage.lean},
\texttt{A5GoodTimeLocalization.lean},
\texttt{Tao/Section5/ReversePrefix.lean},
\texttt{ReversePrefixScalar.lean}, and \texttt{PassLostWindow.lean}.
The written proofs below expose the numerical content of their guards.
They are source-level reconstructions, not a claim to have replayed the
Lean package or to have discovered sampling-without-replacement bounds.

\paragraph{Positive compositions as cuts.}
For $1\le n\le\ell$, let
\[
 \mathcal C_{n,\ell}=\{(a_1,\ldots,a_n)\in\mathbb N_{>0}^n:
                       a_1+\cdots+a_n=\ell\}.
\]
The map to subsets of $\{1,\ldots,\ell-1\}$ is
\begin{equation}
 (a_1,\ldots,a_n)\longmapsto
 \{a_1,a_1+a_2,\ldots,a_1+\cdots+a_{n-1}\}.
 \label{eq:Mazur-cuts}
\end{equation}
Its inverse takes consecutive differences after adjoining $0$ and $\ell$.
The cuts are strictly increasing, so all inverse differences are positive;
telescoping proves both compositions are the identity.
There are $\binom{\ell-1}{n-1}$ such words. Each has the same geometric
probability $2^{-\ell}$ before conditioning, hence the fixed-total law
is exactly the uniform cut-set law. The source stores a physical cut
$s$ as $s-1\in\operatorname{Fin}(\ell-1)$. Thus its strict zero-based
test $s-1<q$ is our inclusive physical test $s\le q$.
For $n=\ell=1$ both carriers consist of one empty-cut configuration.
No corresponding identification with a uniform cut set is asserted
for the infeasible positive-total, zero-length carrier.

\begin{proposition}[finite count and inclusive cut-count tails]
\label{prop:Mazur-cut-count}
Let $\mathcal B$ be a uniform $K$-subset of $\{1,\ldots,N\}$,
where $N\ge1$ and $0\le K\le N$. Put
$H_q=\#(\mathcal B\cap\{1,\ldots,q\})$, $0\le q\le N$,
and $\rho=K/N$. For $0\le r\le N$,
\begin{equation}
 \mathbb E\binom{H_q}{r}
   =\binom qr\frac{\binom Kr}{\binom Nr}
   \le\binom qr\rho^r.
 \label{eq:Mazur-cut-moments}
\end{equation}
For $q\ge1$ and $v\ge0$,
\begin{equation}
 \mathbb P(H_q-q\rho\ge v)\le e^{-2v^2/q},\qquad
 \mathbb P(q\rho-H_q\ge v)\le e^{-2v^2/q}.
 \label{eq:Mazur-cut-tails}
\end{equation}
Both events include equality. For $q=0$ the count is identically zero.
\end{proposition}
\begin{proof}
Count pairs $(R,\mathcal B)$ with $R\subseteq\{1,\ldots,q\}$,
$|R|=r$, and $R\subseteq\mathcal B$. For $r\le K$ their number is
$\binom qr\binom{N-r}{K-r}$; for $r>K$ it is zero.
Dividing by $\binom NK$ proves the identity, with binomial zeros retained.
For $r\le K$ the ratio equals
$\prod_{i=0}^{r-1}(K-i)/(N-i)\le(K/N)^r$;
for $r>K$ its numerator is zero.
More generally a prefix extending beyond $N$ uses $\min(q,N)$
in this count; it is not an additional population.

For $x\ge1$, expand $x^{H_q}=(1+(x-1))^{H_q}$ and extend by
the zero binomial terms to $r=q$. Positivity of $(x-1)^r$ gives
\[
 \mathbb E x^{H_q}\le(1+\rho(x-1))^q.
\]
For a Bernoulli variable of mean $\rho$, set
$g(t)=\log(1-\rho+\rho e^t)-\rho t$.
For $0<\rho<1$ its second derivative is $p_t(1-p_t)\le1/4$,
where $p_t=\rho e^t/(1-\rho+\rho e^t)$.
Since $g(0)=g'(0)=0$, integration twice gives $g(t)\le t^2/8$.
For $\rho=0,1$, $g=0$ and the same bound holds.
At $x=e^t$, $t\ge0$, it follows that
$\mathbb E e^{t(H_q-q\rho)}\le e^{qt^2/8}$.
On the event $H_q-q\rho\ge v$, the exponential is at least $e^{tv}$.
Markov's inequality with $t=4v/q$ proves the upper-tail bound,
including $v=0$. Complementation $\mathcal B\mapsto\mathcal B^c$
is a bijection of the uniform $K$- and $(N-K)$-subset carriers.
It changes $H_q$ to $q-H_q$ and $\rho$ to $1-\rho$,
so the same argument proves the lower tail. No independence of
the individual membership indicators was assumed.
\end{proof}

\begin{proposition}[the fixed-total prefix tail with its integer endpoints]
\label{prop:Mazur-order-statistic}
Let $n\ge3$, $1\le j<n$, $n\le\ell\le3n$, and let
$U_j=a_1+\cdots+a_j$ under the uniform law on $\mathcal C_{n,\ell}$.
For $8\le u\le3j$,
\begin{equation}
 \mathbb P\left(\left|U_j-\frac{j\ell}{n}\right|\ge u\right)
 \le 2\exp\left(-\frac{u^2}{63j}\right).
 \label{eq:Mazur-order-statistic}
\end{equation}
\end{proposition}
\begin{proof}
The variable $U_j$ is the $j$th physical cut. Its mean is
$c=j\ell/n$: permutation of the parts preserves the uniform
composition law, so every part has mean $\ell/n$.
Set $\rho=(n-1)/(\ell-1)$. Direct calculation gives
\[
 \rho\ge\frac14,\qquad j-1\le c\rho\le j,\qquad c\le3j.
\]
Indeed $4(n-1)\ge\ell-1$ follows from $n\ge3$, $\ell\le3n$;
also $j-c\rho=j(\ell-n)/(n(\ell-1))\in[0,1)$.

For the lower deviation put $q=\lfloor c-u\rfloor$.
If $q<j$ the event is empty, since $U_j\ge j$.
Otherwise $1\le q\le3j$ and $q<\ell$. The exact event is
$\{U_j\le q\}=\{H_q\ge j\}$.
As $q\rho\le c\rho-u\rho\le j-u/4$,
Proposition~\ref{prop:Mazur-cut-count} bounds its probability by
$e^{-u^2/(8q)}\le e^{-u^2/(24j)}\le e^{-u^2/(63j)}$.

For the upper deviation put $v=\lceil c+u\rceil$.
The support bound is $U_j\le\ell-(n-j)$; if $v$ exceeds it,
the event is empty. Otherwise put $q=v-1$.
Here $1\le q\le\ell-1$, $q<c+u\le6j$, and $q\ge c+u-1$.
The exact event is
$\{U_j\ge v\}=\{H_{v-1}\le j-1\}$.
In particular the ceiling is followed by a subtraction of one,
also when $c+u$ is integral. Since
$q\rho\ge j-1+(u-1)/4$, the lower cut-count tail gives
$e^{-(u-1)^2/(8q)}\le e^{-(u-1)^2/(48j)}$.
For $u=8+w$, $w\ge0$,
\[
 63(u-1)^2-48u^2=15w^2+114w+15\ge0.
\]
This proves the claimed upper-deviation bound.
The two one-sided events cover the absolute deviation; their sum
proves the assertion.
\end{proof}

\begin{corollary}[both exceptional masses in the fixed-cell descent]
\label{cor:Mazur-two-tails}
Keep $m,n,\ell,S_m,R_m,G,P,Q$ from
Section~\ref{sec:Mazur-descent}, with
$m\ge150$, $8m\le n$, $|\ell-2n|\le n/8$, and
$m|\ell-2n|/n\le S_m$. Then
\begin{equation}
 P(G^c)\le2m^{-25/7}\le2m^{-3},\qquad
 Q(G^c)\le2m^{-8}\le2m^{-3}.
 \label{eq:Mazur-two-tails}
\end{equation}
Thus both tail inputs in Corollary~\ref{cor:Mazur-centred-descent}
now have complete written proofs.
\end{corollary}
\begin{proof}
The law $P$ is the law of $U_m$ conditional on the total $\ell$.
Its mean differs from $2m$ by at most $S_m$.
Outside $G$ its support still has $u\ge m$, so $|u-2m|>R_m$.
Since $R_m\ge16S_m$, this implies
$|u-m\ell/n|\ge15S_m$.
The previously proved $S_m\ge2$ and $R_m\le3m$ imply
$8\le15S_m\le3m$. Also $n\le\ell\le3n$.
Proposition~\ref{prop:Mazur-order-statistic} gives
\[
 P(G^c)\le2e^{-225S_m^2/(63m)}=2m^{-25/7}.
\]

For the unconditioned law $Q$, the centered one-step moment is exactly
\[
 \mathbb E e^{t(a_i-2)}
       =\frac{e^{-t}}{2-e^t},\qquad t<\log2.
\]
Its logarithm $h(t)=-t-\log(2-e^t)$ satisfies
$h(0)=h'(0)=0$ and $h''(t)=2e^t/(2-e^t)^2$.
For $|t|\le1/4$, $e^t\le e^{1/4}<4/3$, because
$\log(4/3)=\int_1^{4/3}x^{-1}\,dx>1/4$.
Thus $h''(t)\le6$, and integration gives
$h(t)\le3t^2\le8t^2$ on this interval.
Independence now gives
$\mathbb E e^{t(\sum_{i=1}^m a_i-2m)}\le e^{8mt^2}$.
Using $t=\min(R/(16m),1/4)$ for the upper deviation and $-t$
for the lower one yields, for every $R>0$,
\[
 Q\left(\left|\sum_{i=1}^m a_i-2m\right|\ge R\right)
 \le2\exp\left(-\min\left(\frac{R^2}{32m},\frac R8\right)\right).
\]
In the first case $R\le4m$, the exponent is $-R^2/(32m)$;
in the second, $-R/4+m/2\le-R/8$.
At $R=R_m\le3m$ the quadratic term is the minimum, and
$R_m^2/(32m)\ge8\log m$. This proves the second bound.
All infinite geometric sums converge in the displayed range of $t$.
\end{proof}

The sharper, unequal tails can also be inserted directly into
Lemma~\ref{lem:Mazur-exact-likelihood}. For the same test as before, let
$\theta_P=2m^{-25/7}$, $\theta_Q=2m^{-8}$ and
\[
 a^\sharp=\max\left(
       \frac{e^{\epsilon_{m,n,\ell}}}{1-\theta_Q}-1,\
       1-e^{-\epsilon_{m,n,\ell}}(1-\theta_P)\right).
\]
Then the error is at most $H\Omega/2+a^\sharp V+H\theta_P$.
This is no greater than the previous bound because each tail is smaller.
In particular the harmonic and flat absolute cell errors may be
$10K/m^{11}+40m^{-25/7}$ and
$40K/m^{11}+160m^{-25/7}$, respectively.
The coefficient $a^\sharp$ is retained as a maximum; unequal tails
do not justify silently replacing it by its first entry.

\paragraph{What a terminal record has to prove about the orbit.}
We now use the positive odd-return map
$T(N)=(3N+1)/2^{v_2(3N+1)}$.
Write $x_k=T^kN$, $a_k=v_2(3x_{k-1}+1)$ and
$S_k=a_1+\cdots+a_k$. For any consecutive block of length $r$,
the exact affine identity gives
\begin{equation}
 x_{j+r}=\frac{3^r}{2^{S_{j+r}-S_j}}x_j+
 \sum_{i=1}^r3^{r-i}2^{-(a_{j+i}+\cdots+a_{j+r})}.
 \label{eq:Mazur-actual-block}
\end{equation}
The offset is nonnegative and at most $3^r$.
For $r=0$ it is zero. For $r\ge1$, bounding each inverse power
by one bounds the sum by $(3^r-1)/2<3^r$.
The \emph{upper} offset bound, not merely its positivity,
is needed to exclude an earlier passage.

Fix $C>0$. Use the unchanged schedules
\[
 L=\log B,\quad n_0=\lfloor L/(10\log2)\rfloor,\quad
 h=\lfloor L/100000\rfloor,
\]
\[
 W=C\sqrt{n_0\log n_0},\quad H_0=L^{3/5},\quad s_0=L^{7/10}.
\]
Put $d=\log(4/3)$, $\delta=d/\log2$, and
\[
 D_-=(4/3)^hB e^{-s_0},\qquad D_+=(4/3)^hB e^{s_0}.
\]
For a target set $E$ of positive odd integers, define the original
finite endpoint set
\[
 E'_B(E)=
 \{M\text{ odd}:\lceil D_-\rceil\le M\le\lfloor D_+\rfloor,\
       \tau_B(M)=h,\ T^hM\in E\},
\]
where $\tau_B$ is the first time at which the orbit is at most $B$,
and is $\infty$ if it never reaches that threshold.
These are exactly the closed integer endpoints of
\texttt{taoSection5EPrime}, not a changed window.

Take a band of positive odd integers $z\le N<U$ with
\[
 B^{1001/1000}\le z<U\le B^{(1001/1000)^3},
 \qquad 0<\log(U/z)<2\log2,
\]
which includes the source's two branches and their exact sub-bands.
Set
\[
 J=\left[\left\lceil\frac{\log(z/B)}d-3W\right\rceil,\
          \left\lfloor\frac{\log(U/B)}d+3W\right\rfloor\right]
        \cap\mathbb N,\qquad
 \mathcal G=\{N:|S_k-2k|<W\ (0\le k\le n_0)\}.
\]
Let $\mathcal A_E$ be the event of a finite first passage into $E$
from this band. Let $\mathcal U_E$ be its terminal-record union:
a member is an odd $N$ in the band for which some $t\in J$,
$r=t-h$, positive word $a$ of length $r$ and total $\ell$, and
$M\in E'_B(E)$ satisfy
\[
 r\le\ell\le3r,\qquad |\ell-2r|<W,\qquad
 \operatorname{Aff}_a(N)=M.
\]
The endpoint bound makes the family finite. The list coordinates are
equivalent to the source's extras multiset: the multiplicity of index
$i$ is $a_i-1$, with inverse $a_i=1+\text{multiplicity}(i)$.
This is the same unrestricted affine union used above;
no first-passage condition on $N$ has been inserted in its definition.

\begin{proposition}[actual first passage equals the terminal union on the good event]
\label{prop:Mazur-genuine-coverage}
For all sufficiently large $B$, depending on $C$ but uniformly in
the band and the target $E$,
\begin{equation}
 \mathcal A_E\cap\mathcal G=\mathcal U_E\cap\mathcal G.
 \label{eq:Mazur-genuine-coverage}
\end{equation}
The map from a good source to its terminal record is its actual first-hit
time $t$, actual first $t-h$ valuations, their total, and $M=T^{t-h}N$.
Its inverse is the integer affine inverse of
Proposition~\ref{prop:Mazur-atoms}.
\end{proposition}
\begin{proof}
We first verify the scalar guards; they are not assumptions about an
unknown orbit. Uniformly over the stated bands, as $L\to\infty$,
\[
 W=o(L),\quad W=o(H_0),\quad H_0=o(s_0),\quad s_0=o(h),\quad
 2\,3^{n_0}\le B.
\]
The last assertion follows from
$3^{n_0}\le B^{\log3/(10\log2)}<B^{1/5}$.
Also $h<\min J\le\max J<n_0$ whenever $J$ is nonempty, and
$W\le(t-h)/8$ for all $t\in J$ eventually.
For the lower edge use
$\log(z/B)\ge L/1000$ and $d<1/3$, giving
$\min J-h\ge(3/1000-1/100000)L-3W$.
For the upper edge use $(1001/1000)^3<101/100$, $d>1/4$ and
$n_0>L/10-1$, giving $\max J\le L/25+3W<n_0$.
Floors and ceilings have the displayed favourable directions.

If $N\in\mathcal G$ lies in the band, \eqref{eq:Mazur-actual-block} gives
\[
 T^{n_0}N\le(3/4)^{n_0}2^W U+3^{n_0}\le B
\]
eventually. Indeed the first term is at most
$(4/3)B^{(1001/1000)^3-d/(10\log2)}2^W$;
its exponent is less than $101/100-1/40<1$ before the $o(1)$
contribution of $2^W$. Thus a first hit exists with $1\le t\le n_0$.

The final inequality $T^tN\le B$, offset positivity, and
$S_t<2t+W$ imply
\[
 t>\frac{\log(N/B)}d-\frac W\delta.
\]
At its predecessor $T^{t-1}N>B$, the upper offset bound and
$2\,3^{n_0}\le B$ imply
$B/2<3^{t-1}N/2^{S_{t-1}}$.
Using $S_{t-1}>2(t-1)-W$ therefore gives
\[
 t<\frac{\log(N/B)}d+\frac{W+1}{\delta}+1.
\]
The exact integer inequality $3^{29}<2^{46}$ implies
$\delta>12/29$. For $W\ge6$,
$W/\delta<3W$ and $(W+1)/\delta+1<3W$.
The band bounds now place $t$ in the original $J$.

Put $r=t-h$ and $M=T^rN$.
Differences of the prefix inequalities give
$2k-2H_0\le S_{r+k}-S_r\le2k+2H_0$ for $k\le h$.
Since $T^hM\le B$,
\[
 M\le(4/3)^h2^{2H_0}B.
\]
Since $T^{h-1}M>B$ and its offset is at most $B/2$,
\[
 M>(3/8)(4/3)^h2^{-2H_0}B.
\]
For large $B$, $2(\log2)H_0+\log(8/3)\le s_0$.
Consequently $D_-\le M\le D_+$, with exactly the stated integer
rounding. First passage shifts to $\tau_B(M)=h$, and the landing
point is unchanged. The actual prefix word has
$|\ell-2r|<W\le r/8$, hence $r\le\ell\le3r$.
This constructs a terminal record and proves the forward inclusion.

Conversely take a record for $N\in\mathcal U_E\cap\mathcal G$.
Its odd affine endpoint decodes to the actual positive valuation word
and $T^rN=M$, by the affine inverse already proved.
Suppose some $j<r$ had $T^jN\le B$. The prefix deviations give
$S_r-S_j\ge2(r-j)-2H_0$, and \eqref{eq:Mazur-actual-block} then gives
\[
 M\le(3/4)^{r-j}2^{2H_0}B+3^{r-j}
      \le2^{2H_0}B+3^{n_0}.
\]
But $M\ge D_-$, whereas
\[
 2^{2H_0}B+3^{n_0}<D_-
\]
eventually: divide by $B$ and use
$hd-s_0-2(\log2)H_0\to+\infty$ and
$3^{n_0}/B\to0$.
This is a contradiction. There is no earlier hit before $r$;
the endpoint's first hit is exactly $h$ steps later, so $\tau_B(N)=t$
and the landing point belongs to $E$. This proves the reverse inclusion.

Finally $\operatorname{diam}J\le\log(U/z)/d+6W=o(h)$,
so Proposition~\ref{prop:Mazur-atoms} gives uniqueness of the record.
The construction and the affine inverse are inverse on these good events.
All eventual bounds used only $C$ and the fixed band envelopes,
not $E$ or the individual source.
\end{proof}

The distinction matters: without the good-prefix and endpoint bounds,
an orbit can hit $[1,B]$, leave it, and later supply a terminal record.
The earlier exact example with starting value $63$ illustrates precisely
that possibility. The reverse argument excludes it here by an explicit
upper bound on the later endpoint of any orbit that had already hit.

\begin{corollary}[one exceptional-event charge for both sampling laws]
\label{cor:Mazur-coverage-mass}
For every probability measure $p$ on the positive odd integers and
all sufficiently large $B$ as above,
\begin{equation}
 |p(\mathcal A_E)-p(\mathcal U_E)|\le p(\mathcal G^c).
 \label{eq:Mazur-coverage-mass}
\end{equation}
In particular this applies separately to the harmonic and uniform
probability laws on a nonempty band. The exceptional mass is charged
once after taking the full union, not once per time or valuation total.
\end{corollary}
\begin{proof}
The two event indicators agree on $\mathcal G$ and differ in absolute
value by at most one on its complement. Integrate this pointwise
inequality. For a measure supported on the band, the same proof can use
$p(\text{band}\cap\mathcal G^c)$, which is equal to $p(\mathcal G^c)$.
\end{proof}

For the original A5 sub-bands, with the band-mass and assembled-test
cap guards stated in Section~\ref{sec:Mazur-descent}, the genuine-band
comparison is now connected by proved maps and event identities:
\[
 |p_\sigma(\mathcal A_E)-\text{common reference profile}|
 \le p_\sigma(\mathcal G^c)
      +\text{the original boundary error}
      +\mathcal R_\sigma^\dagger .
\]
This inequality combines the actual-event identification here, the
nominal boundary comparison, and the refined terminal descent.
The band-measure estimates for $p_\sigma(\mathcal G^c)$ are proved
in Section~\ref{sec:Mazur-band-fullgood} below, through an exact residue
decoder and separate estimates for both sampling laws. Evaluation of
the common reference profile remains a separate mathematical step.

\subsection{Actual band laws and a maximal prefix bound}
\label{sec:Mazur-band-fullgood}

The good-prefix event is a property of an actual integer orbit. It is
therefore necessary to transport the geometric valuation law to the
integer sampling law, with a quantitative error. There is also a
separate improvement in the geometric estimate: summing a tail over
all times is unnecessary when the event asks whether a partial sum
ever crosses one common threshold.

The source comparison here is with Mazur~\cite{Mazur2026v2},
\texttt{ND/Band/A5HarmonicResidue.lean},
\texttt{A5HarmonicProp19.lean}, \texttt{A5HarmonicFullGood.lean},
\texttt{A5BandCoverageFullGood.lean}, and
\texttt{Tao/Probability/Geom2PrefixTypical.lean}, at the same pinned
commit. Its strict-truncation mechanism is in
\texttt{Tao/Syracuse/Prop19Laws.lean},
\texttt{Prop19Truncation.lean}, and
\texttt{TruncatedValuationTV.lean}.
The underlying valuation approximation is Tao's
Proposition~1.9~\cite{Tao2026v7}, source label \texttt{rach};
the proof is in his Section~4, source lines 597--647.
The following argument reconstructs the finite transport, then uses
the exact geometric moment to sharpen this particular prefix estimate.
It makes no priority claim for exponential stopping arguments.

\paragraph{The two sampling laws on the same integers.}
Let
\[
 I= \{N\in2\mathbb N+1:z\le N<U\},\qquad
 D=\#I>0,\qquad H=\sum_{N\in I}\frac1N,\qquad z\ge1.
\]
The upper endpoint remains excluded, even when it is an odd integer.
The uniform and harmonic laws are respectively
\[
 p_F(N)=D^{-1}\mathbf1_I(N),\qquad
 p_H(N)=\frac{\mathbf1_I(N)}{NH}.
\]
On their common support the exact density ratio is
$p_F(N)/p_H(N)=NH/D$, with inverse $D/(NH)$.
These are distinct probability measures on the same carrier;
the estimates below are proved separately for both.

\begin{proposition}[exact band residue errors]
\label{prop:Mazur-band-residues}
Let $M\ge1$, $q=2^M$, and let $\nu_M$ be the uniform law on the
$q/2$ odd residue classes modulo $q$. Let $\pi_M(N)=N\bmod q$.
Then, with full $\ell^1$ distance,
\begin{equation}
 \|\pi_{M*}p_F-\nu_M\|_1\le\frac{q}{2D},\qquad
 \|\pi_{M*}p_H-\nu_M\|_1\le\frac{q}{2\lceil z\rceil H}.
 \label{eq:Mazur-band-residues}
\end{equation}
The estimates retain the actual $D,H$ and real band endpoints.
\end{proposition}
\begin{proof}
Enumerate $I$ as $N_i=N_1+2(i-1)$, $1\le i\le D$.
For an odd residue $b$, put $a=q/2$,
$d_i=\mathbf1_{\{N_i\bmod q=b\}}-1/a$ and
$F_j=\sum_{i=1}^j d_i$. The residue repeats once every $a$ indices.
Writing $j=ka+r$, $0\le r<a$, proves $|F_j|\le1$,
including $a=1$. In particular the flat atom error is at most $1/D$.
Summing over the $a$ classes gives the first bound.

For $w_i=1/N_i$, finite summation by parts is the exact identity
\[
 \sum_{i=1}^D d_iw_i
 =F_Dw_D+\sum_{i=1}^{D-1}F_i(w_i-w_{i+1}).
\]
The weights decrease and are nonnegative. Hence the absolute value
is at most $w_D+\sum_{i<D}(w_i-w_{i+1})=w_1
\le1/\lceil z\rceil$. Divide by $H$ and sum over the same $a$
classes. No asymptotic count or endpoint replacement has been used.
\end{proof}

\begin{proposition}[strict decoding and the single ideal overflow mass]
\label{prop:Mazur-truncation-transport}
Let $n,M\ge1$, let $p$ be any probability law on positive odd
integers, and write $P_n$ for the law of its first $n$ actual valuations.
Let $Q_n=\operatorname{Geom}(2)^n$, and set
$\Delta_M=\|\pi_{M*}p-\nu_M\|_1$.
Then
\begin{equation}
 \|P_n-Q_n\|_1\le\Delta_M+2\Theta(n,M),\qquad
 \Theta(n,M)=2^{-(M-1)}
       \sum_{j=0}^{n-1}\binom{M-1}{j}.
 \label{eq:Mazur-truncation-transport}
\end{equation}
Consequently every set $A$ of length-$n$ positive words satisfies
\begin{equation}
 |P_n(A)-Q_n(A)|\le\frac{\Delta_M}{2}+\Theta(n,M).
 \label{eq:Mazur-event-half}
\end{equation}
Terms outside the binomial support are zero.
\end{proposition}
\begin{proof}
Let $V_n$ send an odd integer to its actual ordered valuation word.
Let $\mathcal W_{n,M}$ be the finite set of such positive words with
total $s<M$, and define $T_{n,M}$ on all positive length-$n$ words
by retaining the word in $\mathcal W_{n,M}$ and sending every other
word to one symbol $\bot$.

For an accepted word $a$, Lemma~\ref{lem:Tao-Prop19-residue-fibre} gives its
unique odd cylinder $u_a\bmod2^{s+1}$.
As $s+1\le M$, membership is determined by a residue modulo $2^M$.
Distinct accepted words have disjoint residue fibres: a common odd
residue has a positive lift, which cannot have two actual valuation
words. Define $d_{n,M}$ on the odd residues to output the accepted
word whose cylinder contains the residue, or $\bot$ if there is none.
This proves the commuting identity
\[
 d_{n,M}\circ\pi_M=T_{n,M}\circ V_n.
\]
The accepted fibre over $a$ has exactly $2^{M-s-1}$ residues,
namely the lifts of $u_a$ modulo $2^M$. Its uniform probability is
$2^{M-s-1}/2^{M-1}=2^{-s}$, exactly the geometric word probability.
All remaining residues form the one overflow fibre. Thus
$d_{n,M*}\nu_M=T_{n,M*}Q_n$.

Pushforward contracts full $\ell^1$: group the absolutely summable
differences in each fibre and apply the triangle inequality.
It follows that $\|T_{n,M*}P_n-T_{n,M*}Q_n\|_1\le\Delta_M$.
If $b_P,b_Q$ are the two overflow masses, the full untruncated
distance is at most the accepted-atom distance plus $b_P+b_Q$.
The truncated distance already contains $|b_P-b_Q|$, and
\[
 b_P+b_Q-|b_P-b_Q|=2\min(b_P,b_Q)\le2b_Q.
\]
This proves the first inequality with $b_Q=Q_n(s\ge M)$;
no separate bound for the actual orbit's overflow was assumed.

To evaluate $b_Q$, realize each geometric variable as the waiting
time to the next success in independent fair Bernoulli trials.
A run of $a_i-1$ failures followed by one success has probability
$2^{-a_i}$; disjoint runs multiply to $2^{-s}$.
The $n$th success occurs at time at least $M$ precisely when at most
$n-1$ successes occur among the first $M-1$ trials. Counting those
finite Bernoulli strings gives the displayed formula, also when
$M\le n$ and that probability is one.

Finally $P_n-Q_n$ has total mass zero. Its positive and negative
masses are equal, each half its full $\ell^1$ norm. The mass on
any event is between the negative of this number and this number.
This proves \eqref{eq:Mazur-event-half} without changing the
source's full-$\ell^1$ convention.
\end{proof}

\paragraph{One maximal estimate instead of a union over times.}
Let $A_i$ be independent with $\mathbb P(A_i=a)=2^{-a}$,
$a\ge1$, and write $X_k=\sum_{i=1}^k(A_i-2)$.
Its exact one-step logarithmic moment is
\[
 h(t)=\log\mathbb E e^{t(A_i-2)}
     =-t-\log(2-e^t),\qquad t<\log2.
\]
It has $h(0)=h'(0)=0$, $h''(t)>0$, and therefore $h(t)\ge0$.

\begin{proposition}[maximal geometric tail with the exact rate]
\label{prop:Mazur-maximal-tail}
For $n\ge2$ and real $0<w<n$, define
\[
 I(v)=(1+v)\log(1+v)-(2+v)\log(1+v/2)
       \quad(-1<v<\infty),
\]
\begin{equation}
 \Psi(n,w)=e^{-nI(w/n)}+e^{-nI(-w/n)}.
 \label{eq:Mazur-maximal-rate}
\end{equation}
Then
\begin{equation}
 \mathbb P\!\left(\max_{1\le k\le n}|X_k|\ge w\right)
       \le\Psi(n,w).
 \label{eq:Mazur-maximal-tail}
\end{equation}
There is no factor $n$ multiplying this estimate.
Also, for every $n\ge1$,
\begin{equation}
 \Theta(n,3n)\le(27/32)^n.
 \label{eq:Mazur-overflow-rate}
\end{equation}
\end{proposition}
\begin{proof}
For fixed $t<\log2$, $Z_k=\exp(tX_k-kh(t))$ has expectation
one and satisfies
$\mathbb E(Z_{k+1}\mid A_1,\ldots,A_k)=Z_k$.
This follows directly from independence and the displayed moment.
For $t>0$, stop at the first $k\le n$ with $X_k\ge w$, or at
$n$ if no crossing occurs. Call that bounded time $\tau$.
The identity
\[
 Z_{(k+1)\wedge\tau}-Z_{k\wedge\tau}
   =\mathbf1_{\{\tau>k\}}(Z_{k+1}-Z_k)
\]
has zero expectation: the indicator depends only on the first $k$
variables. Summing this finite identity proves $\mathbb E Z_\tau=1$.
All its terms are integrable; for example
$Z_\tau\le\sum_{k=0}^n Z_k$.
On a crossing, $Z_\tau\ge e^{tw-nh(t)}$, since $h(t)\ge0$.
Thus the upper maximum has probability at most $e^{nh(t)-tw}$.

For the lower maximum apply the same argument at a negative
parameter $t$, stopping when $X_k\le-w$. Its bound is
$e^{nh(t)+tw}$. Writing $r=w/n$, choose respectively
\[
 t_+=\log\frac{2(1+r)}{2+r},\qquad
 t_-=\log\frac{2(1-r)}{2-r}.
\]
Here $0<t_+<\log2$ and $t_-<0$, because $0<r<1$.
Since
\[
 h'(t)=\frac{2(e^t-1)}{2-e^t},
\]
these parameters solve $h'(t_+)=r$ and $h'(t_-)=-r$.
Substitution, with no discarded terms, gives
$rt_+-h(t_+)=I(r)$ and $-rt_--h(t_-)=I(-r)$.
The union of the two crossing events proves the result.

For the overflow only the final total is needed.
Exponential Markov at $t=\log(4/3)$ gives
\[
 Q_n(X_n\ge n)\le e^{-nt+nh(t)}
  =\left(\frac{27}{32}\right)^n,
 \qquad h(\log(4/3))=\log(9/8).
\]
The event is exactly the strict-decoder overflow $s\ge3n$.
\end{proof}

For comparison, using merely $h(\pm t)\le8t^2$ for $0\le t\le1/4$
in the same stopped argument gives
$2e^{-\min(w^2/(32n),w/8)}$, still without a union factor.
The optimized rate above retains more of the exact geometric law.
Its behavior near zero is explicit:
\[
 I(0)=I'(0)=0,\qquad I''(v)=\frac1{(1+v)(2+v)},\qquad
 I(v)=v^2\int_0^1\frac{1-u}{(1+uv)(2+uv)}\,du.
\]
On $|v|\le1/2$ the integrand and its $v$-derivative are bounded,
so $I(v)=v^2/4+O(|v|^3)$, with an absolute remainder constant.

\begin{corollary}[full-prefix failure under both actual band laws]
\label{cor:Mazur-band-fullgood}
Use the unchanged schedules of
Section~\ref{sec:Mazur-concentration-coverage}, and put $n=n_0$,
$W=C\sqrt{n\log n}$ and $w=\lceil W\rceil$.
For every fixed $C>0$, sufficiently large $B$ gives $1\le w<n$.
For every nonempty band $I$ as above, let $\mathcal G$ retain the
literal strict inequalities $|S_k-2k|<W$ for $0\le k\le n$.
Then
\begin{align}
 p_H(\mathcal G^c)
 &\le \Psi(n,w)+\frac{2^{3n-2}}{\lceil z\rceil H}
                  +\Theta(n,3n), \label{eq:Mazur-harmonic-fullgood}\\
 p_F(\mathcal G^c)
 &\le \Psi(n,w)+\frac{2^{3n-2}}{D}
                  +\Theta(n,3n). \label{eq:Mazur-flat-fullgood}
\end{align}
On the source's exact A5 sub-bands these imply, uniformly in the
branch and valid band index,
\begin{equation}
 p_\sigma(\mathcal G^c)
 \le\Psi(n,w)+B^{-701/1000}+(27/32)^n
 \le(2+o_C(1))n^{-C^2/4},\qquad \sigma=H,F.
 \label{eq:Mazur-band-fullgood-rate}
\end{equation}
The last expression is an asymptotic upper bound, not an asymptotic
formula for the actual exceptional probability.
\end{corollary}
\begin{proof}
Every deviation $S_k-2k$ is integral. Consequently
$|S_k-2k|\ge W$ is exactly $|S_k-2k|\ge w$, including integral $W$.
Equivalently it is strictly greater than the source radius
$\lceil W\rceil-1$. The zero prefix is never bad when $W>0$.
Thus $\mathcal G^c$ is the preimage under $V_n$ of the single
maximal event in Proposition~\ref{prop:Mazur-maximal-tail}.
Apply \eqref{eq:Mazur-event-half} at $M=3n$, with the two separate
residue estimates \eqref{eq:Mazur-band-residues}. This proves both
finite-band bounds.

Here are the band constants, with their original endpoints retained.
For an A5 branch $y=B^\alpha$ or $B^{\alpha^2}$,
$\alpha=1001/1000$, put
\[
 b=(\alpha-1)\log y,\quad K=\lfloor b\rfloor,\quad
 \beta=b/K,\quad z=y e^{j\beta},\quad U=z e^\beta,\quad 0\le j<K.
\]
Eventually $K\ge1$, $\beta\ge1$, $\beta\to1$, $z\ge B^\alpha$
and $z\ge6$. In particular $U\ge ez$.
The first odd point is less than $z+2$, and the successor of the
last odd point is at least $U$. Comparing the decreasing function
$1/t$ on each interval of length two gives
\[
 H\ge\frac12\log\frac{U}{z+2}
 \ge\frac12\left(1-\frac2z\right)\ge\frac13\ge\frac14.
\]
Direct lattice counting gives
$D>(U-z)/2-1\ge(e-1)z/2-1\ge z/2$ for $z\ge6$;
for the last inequality $e>8/3$ follows from its first five
Taylor terms. Since $n=\lfloor\log B/(10\log2)\rfloor$,
\[
 2^{3n}\le B^{3/10},\qquad
 \frac{2^{3n-2}}{\lceil z\rceil H}\le B^{-701/1000},\qquad
 \frac{2^{3n-2}}D\le\frac12 B^{-701/1000}.
\]
Keep the exact denominators in the finite bounds even though the
common upper bound uses the larger of these last two quantities.

Finally $w=C\sqrt{n\log n}+O(1)$. The integral expression for $I$ gives
\[
 nI(\pm w/n)=\frac{C^2}{4}\log n
   +O_C\!\left(\frac{(\log n)^{3/2}}{\sqrt n}\right)+o(1).
\]
Hence $\Psi(n,w)=(2+o_C(1))n^{-C^2/4}$.
The remaining two terms are negligible relative to every fixed
negative power of $n$: $n$ is comparable to $\log B$, and
\[
 (27/32)^n\le(32/27)B^{-\kappa},\qquad
 \kappa=\frac{\log(32/27)}{10\log2}>0.
\]
All constants used the band envelopes, not its index or a target set.
\end{proof}

This replaces the source's positive-prefix union bound
\[
 2n\exp\left[-\min\left(\frac{R^2}{32n},\frac R8\right)\right]
       +4\,2^{-n/128},\qquad R=\lceil W\rceil-1,
\]
whose first term is $(2+o_C(1))n^{\,1-C^2/32}$.
The improvement has three explicit causes: a maximum is controlled
without summing over times; the exact geometric moment is retained;
and an event costs half the full $\ell^1$ error.
The uniform band is treated directly, so no factor $e^\beta$ from
harmonic domination is needed. The new component estimate decays
for every fixed $C>0$; this is not a change to the source's separate
phase estimate or a claim of a stronger final natural-density theorem.

For the original A5 bands, Corollary~\ref{cor:Mazur-coverage-mass}
now bounds the genuine-passage/terminal-union discrepancy by the
appropriate right-hand side of
\eqref{eq:Mazur-harmonic-fullgood} or
\eqref{eq:Mazur-flat-fullgood}. Combining with the earlier boundary
comparison and $\mathcal R_\sigma^\dagger$ is therefore justified
with this quantitative exceptional term. A convex combination of
band laws retains their common bound because its weights sum to one;
it does not multiply the error by the number of bands.
The excluded upper endpoint of the full branch interval is not in
that combination. Its own singleton mass is retained in
Corollary~\ref{cor:Mazur-genuine-source}. The common profile is evaluated
in Section~\ref{sec:Mazur-exterior-assembly}; the resulting actual passage
comparison is Theorem~\ref{thm:Mazur-genuine-rate}.

\begin{corollary}[an explicit prefix constant in the Allikvere window]
\label{cor:Mazur-allikvere-prefix}
Let $x\to\infty$, $L=\log x$, $\alpha=1001/1000$,
$y=x^\alpha$ or $x^{\alpha^2}$, and
$n=\lfloor L/(10\log2)\rfloor$.
If $N$ is uniform on the positive odd integers in the
\emph{closed} interval $[y,y^\alpha]$, then
\begin{equation}
 \mathbb P\left(\max_{0\le k\le n}|S_k(N)-2k|\ge L^{3/5}\right)
 \le(2+o(1))\exp\left[-\frac{5\log2}{2}L^{1/5}\right].
 \label{eq:Mazur-allikvere-prefix}
\end{equation}
The error is uniform in the two branches.
\end{corollary}
\begin{proof}
Keep the closed upper endpoint by taking the equivalent half-open
integer carrier $y\le N<\lfloor y^\alpha\rfloor+1$.
Its odd count $D_y$ is comparable to $y^\alpha$.
Propositions~\ref{prop:Mazur-band-residues}
and~\ref{prop:Mazur-truncation-transport}, with $M=3n$, give event
error at most $2^{3n-2}/D_y+\Theta(n,3n)$.
Both terms decrease as a positive power of $x$.
Apply Proposition~\ref{prop:Mazur-maximal-tail} at
$w=\lceil L^{3/5}\rceil<n$. The integral rate identity gives
\[
 nI(\pm w/n)=\frac{w^2}{4n}+O(w^3/n^2)
            =\frac{5\log2}{2}L^{1/5}+o(1).
\]
Here $w^3/n^2=O(L^{-1/5})$ and the floor/ceiling changes also
contribute $o(1)$. Summing the two exponentials proves the result;
the power errors in $x$ are negligible relative to this exponential.
\end{proof}

This supplies an explicit constant for the earlier uniform-input
prefix bound. It does not replace the separate stronger power bound
for the event that no passage occurs by $n_0$: those are different
events, and their established estimates are both retained.

\subsection{The common endpoint profile with exact band corrections}
\label{sec:Mazur-common-profile}

The common center is not obtained by identifying uniform and harmonic
sampling. Their raw sums have different phase functions and different
means. This section reconstructs the finite comparison before estimating
the oscillatory sums. In particular, it keeps the shifted raw mass,
the two band denominators and the exterior center visible.

The source locators are Mazur~\cite{Mazur2026v2},
\texttt{A5TwoProfileRawQ.lean}, \texttt{A5TwoProfilePhysicalA6.lean},
\texttt{A5TwoProfileNormalizedA6.lean},
\texttt{A5ReferencePhysicalSplit.lean},
\texttt{A5ReferenceInteriorAggregate.lean}, and
\texttt{A5ReferenceFullCenterRate.lean}, at the pinned commit already
specified. The kernel is Tao's common term, represented in the package
by \texttt{Tao/Section5/CommonZ.lean}. The identities below are finite,
and do not presume the interior equidistribution or exterior tail estimates.

\paragraph{Exact lattice corrections.}
Keep the original odd band $I=[z,U)\cap(2\mathbb N+1)$,
$z\ge1$, $U>z$, $D=\#I>0$, $H=\sum_I1/N$,
and $\beta=\log(U/z)$. Define the signed discrepancies
\[
 \epsilon_D=D-\frac{U-z}{2},\qquad
 \epsilon_H=H-\frac{\beta}{2}.
\]
They are not discarded by replacing the band with a continuous interval.

\begin{proposition}[band discrepancies and retained denominator errors]
\label{prop:Mazur-band-denominators}
For the band just defined,
\begin{equation}
 |\epsilon_D|<1,\qquad |\epsilon_H|\le\frac1z.
 \label{eq:Mazur-band-denominators}
\end{equation}
For $d=\log(4/3)>0$, put $a=2/d$.
For any real raw values $R_H,R_F$, the following are exact:
\begin{align}
 \frac{R_H}{H}-a
 &=\frac{R_H-\beta/d}{H}-\frac{2\epsilon_H}{dH},
 \label{eq:Mazur-harmonic-denominator}\\
 \frac zD R_F-a
 &=\frac zD\left(R_F-\frac{e^\beta-1}{d}\right)
       -\frac{2\epsilon_D}{dD}.
 \label{eq:Mazur-flat-denominator}
\end{align}
Thus the absolute denominator charges are at most
$2/(zdH)$ and strictly less than $2/(dD)$, respectively.
If $z\ge6$ and $\beta\ge1$, then $H\ge1/3$ and $D>z/2$,
so these charges are at most $6/(zd)$ and less than $4/(zd)$.
\end{proposition}
\begin{proof}
For $z\le t\le U$, write
\[
 C(t)=\#\{N\text{ positive odd}:z\le N<t\},\qquad
 \Delta(t)=C(t)-(t-z)/2 .
\]
For $b\ge a$, the exact count of integers in $[a,b)$ is
$\lceil b\rceil-\lceil a\rceil$.
Apply this at $a=(z-1)/2$, $b=(t-1)/2$.
The two errors $\lceil u\rceil-u$ belong to $[0,1)$,
so $|\Delta(t)|<1$ and the claimed bound for $\epsilon_D=\Delta(U)$
follows, including integral odd endpoints.

For each $N\in I$ the identity
$1/N=1/U+\int_N^U t^{-2}\,dt$ gives, after a finite sum,
\[
 H=\frac D U+\int_z^U\frac{C(t)}{t^2}\,dt.
\]
The choices $N<t$ and $N\le t$ differ only at finitely many points
and give the same ordinary integral. The continuous expression
\[
 \frac{U-z}{2U}+
 \int_z^U\frac{t-z}{2t^2}\,dt=\frac12\log(U/z)
\]
is verified by integration. Subtracting proves the signed formula
\[
 \epsilon_H=\frac{\Delta(U)}U+
                      \int_z^U\frac{\Delta(t)}{t^2}\,dt .
\]
Its absolute value is at most
$1/U+\int_z^Ut^{-2}\,dt=1/z$.

Equations \eqref{eq:Mazur-harmonic-denominator} and
\eqref{eq:Mazur-flat-denominator} now follow by substitution and
$U-z=z(e^\beta-1)$. Finally
$H\ge\beta/2-1/z\ge1/3$ and
$D>(e-1)z/2-1>z/2$ for $z\ge6$; for the last strict inequality
$e>8/3$. All denominators are the original positive sums.
\end{proof}

Compared with the source estimates
$|H-\beta/2|\le3/(2z)$, $1/H\le4$ and $z/D<3$,
the displayed proof retains the exact signed errors and yields the
bounds $|H-\beta/2|\le1/z$, $1/H\le3$ and $z/D<2$
under the same $z\ge6,\beta\ge1$ conditions.
In particular an absolute raw harmonic error $E_H$ costs
$E_H/H+2/(zdH)\le3E_H+6/(zd)$, and an absolute raw flat
error $E_F$ costs $(z/D)E_F+2/(dD)\le2E_F+4/(zd)$.
The original coarse constants $4E_H+12/(zd)$ and
$3E_F+6/(zd)$ are preserved in the source comparison, not silently
substituted into its unreviewed consumers.

\paragraph{The literal two phase functions.}
Let $\gamma=\log_2 3$, $\delta=2-\gamma=d/\log2$, and suppose
$\log2<\beta<2\log2$. These inequalities hold on the source A5
bands for sufficiently large $B$.
For $W>0$, $A=\log_2(z/M)$ and $\nu\ge1$, define
\[
 s_\nu=\lfloor A+\nu\gamma\rfloor+1,\quad
 \theta_\nu=\{A+\nu\gamma\},\quad
 b=2-\beta/\log2\in(0,1),
\]
and, for every integer $s$,
\[
 p_\nu(s)=
 \begin{cases}\binom{s-1}{\nu-1}2^{-s},&s\ge\nu,\\0,&s<\nu.\end{cases}
\]
Retain the two distinct indicated atoms
\[
 q_0(\nu)=p_\nu(s_\nu)\mathbf1_{\{|s_\nu-2\nu|<W\}},\qquad
 q_1(\nu)=p_\nu(s_\nu+1)\mathbf1_{\{|s_\nu+1-2\nu|<W\}}.
\]
The strict integer anchor is used even at integral phase.
These formulas agree with the source's raw atoms at $A$ and $A+1$;
for $\nu>0$, clipping a negative signed level to zero leaves its
probability zero. The tube's finite upper cap removes no index:
$|\lfloor A-\delta\nu\rfloor+1|<W$ implies
$\nu\le(A+\lceil W\rceil)/\delta$.
Let $V$ be the unchanged positive physical time carrier
$\{\nu>0:\nu+h\in J\}$.

\begin{proposition}[two distinct raw profiles and their phase means]
\label{prop:Mazur-profile-phase}
For $\theta\in[0,1)$, set
\[
 f_H(\theta)=1+\mathbf1_{\{\theta>b\}},\qquad
 f_F(\theta)=2^{1-\theta}
                 +\mathbf1_{\{\theta>b\}}2^{2-\theta}.
\]
The two raw sums on $V$ are exactly
\begin{align}
 R_H(M)&=\sum_{\nu\in V}
        \big(q_0(\nu)+\mathbf1_{\{\theta_\nu>b\}}q_1(\nu)\big),\\
 R_F(M)&=\sum_{\nu\in V}
        \big(2^{1-\theta_\nu}q_0(\nu)
          +\mathbf1_{\{\theta_\nu>b\}}2^{2-\theta_\nu}q_1(\nu)\big).
 \label{eq:Mazur-literal-profiles}
\end{align}
Equivalently,
\begin{align}
 R_H(M)&=\sum_Vq_0(\nu)f_H(\theta_\nu)
       +\sum_V\mathbf1_{\{\theta_\nu>b\}}(q_1(\nu)-q_0(\nu)),\\
 R_F(M)&=\sum_Vq_0(\nu)f_F(\theta_\nu)
       +\sum_V\mathbf1_{\{\theta_\nu>b\}}2^{2-\theta_\nu}
                                      (q_1(\nu)-q_0(\nu)).
 \label{eq:Mazur-neighbour-residual}
\end{align}
Their ordinary integrals are
\begin{equation}
 \int_0^1 f_H(t)\,dt=\frac{\beta}{\log2},\qquad
 \int_0^1 f_F(t)\,dt=\frac{e^\beta-1}{\log2}.
 \label{eq:Mazur-phase-means}
\end{equation}
The gated upper-cell weight is strictly below $e^\beta$;
at $\theta=b$ the upper cell is absent.
\end{proposition}
\begin{proof}
The base position $(1-\theta)\log2$ belongs to $(0,\beta)$
throughout $[0,1)$. The successor position belongs to that
interval exactly when $(2-\theta)\log2<\beta$, or $\theta>b$.
This proves the first two formulas directly from
Proposition~\ref{prop:Mazur-two-rows}; the flat exponential is
attached once to its own cell. Add and subtract the gated $q_0$
term to obtain the next pair. No equality of the two tube supports
or their masses is asserted.

The harmonic integral is $1+(1-b)=\beta/\log2$.
The two flat integrals are
\[
 \int_0^1 2^{1-t}\,dt=\frac1{\log2},\qquad
 \int_b^1 2^{2-t}\,dt=\frac{2^{2-b}-2}{\log2}.
\]
Since $2^{2-b}=e^\beta$, their sum is as claimed.
The strict gate also gives $2^{2-\theta}<e^\beta$ when active.
Its endpoint exclusion changes neither integral but remains
essential to the finite sums.
\end{proof}

This explains the two raw centers: multiplying the respective means
by the source's raw-mass center $1/\delta$ gives
$\beta/d$ and $(e^\beta-1)/d$. It does not establish that an
individual raw sum equals its center. Both the weighted rotation error
and the signed neighbour correction in
\eqref{eq:Mazur-neighbour-residual} must still be estimated on their
actual supports.

\paragraph{The common kernel and the landing map.}
For a fixed depth $m$, retain
\[
 K_m(M)=\frac{3^m}{M}p_m(M\bmod3^m)\quad(M>0),
\]
where $p_m$ is the source's Syracuse offset probability law.
Let $\mathcal E_*$ be the original positive odd endpoint window
with $\tau_B(M)=h$, without imposing a landing target. Its lower
and upper integer endpoints remain
$\lceil(4/3)^hB e^{-(\log B)^{7/10}}\rceil$ and
$\lfloor(4/3)^hB e^{(\log B)^{7/10}}\rfloor$.
Define $\lambda(M)=T^hM\in[1,B]\cap(2\mathbb N+1)$.
Then the original $\mathcal E(E)$ is
$\mathcal E_*\cap\lambda^{-1}(E)$.

The weighted finite endpoint measure $\mu_m=\sum_{\mathcal E_*}
K_m(M)\delta_M$ therefore has the exact landing pushforward
\[
 Z(E)=\frac2d\,\lambda_*\mu_m(E)
     =\frac2d\sum_{M\in\mathcal E(E)}K_m(M).
\]
This is the source common profile. It has no band or sampling-law
parameter. It is not declared a probability measure: its total mass
has not been set to one. The landing map can have multiple endpoints
in a fibre, and that fibre has precisely the sum of their weights.
The map from endpoint measures to landing measures is linear and
positive; it loses the distribution within each landing fibre.

\begin{proposition}[signed common-center identity on the actual carrier]
\label{prop:Mazur-profile-residual}
For $S=\mathcal E(E)$, put $Z_0(S)=\sum_SK_m(M)$ and
\[
 r_H(M)=R_H(M)-\beta/d,\qquad
 r_F(M)=R_F(M)-(e^\beta-1)/d.
\]
The actual reference profiles
\[
 P_H(S)=\sum_SK_m(M)R_H(M)/H,\qquad
 P_F(S)=\sum_SK_m(M)(z/D)R_F(M)
\]
satisfy
\begin{align}
 P_H(S)-aZ_0(S)
 &=\frac1H\sum_SK_m(M)r_H(M)
                -\frac{2\epsilon_H}{dH}Z_0(S),\\
 P_F(S)-aZ_0(S)
 &=\frac zD\sum_SK_m(M)r_F(M)
                -\frac{2\epsilon_D}{dD}Z_0(S).
 \label{eq:Mazur-common-signed}
\end{align}
These formulas retain signs before the absolute-value estimate.
They hold for the empty target as well. In particular, the
denominator charge occurs once against $Z_0(S)$, not once per
endpoint or per time.
\end{proposition}
\begin{proof}
Substitute \eqref{eq:Mazur-harmonic-denominator} or
\eqref{eq:Mazur-flat-denominator} at every endpoint, multiply by
$K_m(M)$, and sum. The band discrepancies are independent of $M$.
The kernel is nonnegative because $M>0$ and $p_m$ is a probability
mass function, but positivity is not needed for the identities.
It gives the corresponding absolute error bounds by applying the
triangle inequality to these finite sums.
\end{proof}

\paragraph{An explicit all-endpoint envelope.}
For the same finite positive time set $V$, define
\[
 G_V=\sum_{\nu\in V}
         \frac{\binom{2\nu-2}{\nu-1}}{2^{2\nu-2}},\qquad
 L_H=\frac{G_V}{H},\qquad L_F=\frac{z e^\beta G_V}{D}.
\]
These quantities depend on the actual time carrier and band,
but not on the endpoint $M$.

\begin{proposition}[finite envelope and sharp exterior centering]
\label{prop:Mazur-profile-envelope}
For every positive endpoint $M$,
\[
 0\le R_H(M)/H\le L_H,\qquad
 0\le(z/D)R_F(M)\le L_F,\qquad
 G_V\le\sum_{\nu\in V}\nu^{-1/2}.
\]
Split the original endpoint set at its closed interior condition
\[
 S_{\rm in}=\{M\in S:
  |\log M-\log B-hd|\le4W/25\},\qquad
 S_{\rm out}=S\setminus S_{\rm in}.
\]
For either $\sigma=H,F$, write $P_{\sigma,\rm out}$ for its
profile restricted to $S_{\rm out}$ and
$Z_{\rm out}=\sum_{S_{\rm out}}K_m(M)$. Then
\begin{equation}
 |P_{\sigma,\rm out}-aZ_{\rm out}|
 \le\max\{a,L_\sigma-a\}\,Z_{\rm out}.
 \label{eq:Mazur-sharp-exterior}
\end{equation}
The constant is sharp given only $0\le P_{\sigma,\rm out}
\le L_\sigma Z_{\rm out}$.
For the full profile one may combine
\eqref{eq:Mazur-common-signed} on $S_{\rm in}$ with this
exterior term, retaining the exact denominators and
$Z_0(S_{\rm in})$ in the interior charge.
\end{proposition}
\begin{proof}
The adjacent endpoint masses obey, for $s\ge\nu$,
\[
 \frac{p_\nu(s+1)}{p_\nu(s)}
      =\frac{s}{2(s+1-\nu)}.
\]
For $\nu\ge2$ this ratio is at least one exactly when
$s\le2\nu-2$. The two maximal atoms are at $2\nu-2$ and
$2\nu-1$, with common mass
\[
 b_\nu=\frac{\binom{2\nu-2}{\nu-1}}{2^{2\nu-1}}.
\]
For $\nu=1$, $p_1(s)=2^{-s}$ has maximum $1/2=b_1$
at $s=1$. Thus each of the two indicated raw atoms is at most
$b_\nu$. Each active flat weight is less than $e^\beta$.
Summing gives $R_H\le2\sum_Vb_\nu=G_V$ and
$R_F\le e^\beta G_V$, without replacing an exterior physical
window by an intrinsic tube.

For completeness, let $c_k=\binom{2k}{k}/4^k$.
The recurrence $c_{k+1}/c_k=(2k+1)/(2k+2)$ and
\[
 4(k+1)^3-(2k+1)^2(k+2)=3k+2>0
\]
prove by induction from $c_0=1$ that
$c_k\le1/\sqrt{k+1}$. This proves the bound for $G_V$.

Nonnegative kernel weighting gives
$0\le P_{\sigma,\rm out}\le L_\sigma Z_{\rm out}$.
If $Z_{\rm out}=0$ the desired statement is immediate.
Otherwise $P_{\sigma,\rm out}/Z_{\rm out}$ lies in $[0,L_\sigma]$.
The largest distance of this interval from $a>0$ is
$\max\{a,|L_\sigma-a|\}=\max\{a,L_\sigma-a\}$.
Both endpoints of the interval realize the corresponding distances,
which proves sharpness at this information level.
Finally $S_{\rm in}$ and $S_{\rm out}$ are a disjoint partition,
including equality in the interior, so the two signed errors add
exactly before applying the triangle inequality.
\end{proof}

The source full-center argument instead uses a coarser common
envelope and the sum of that envelope with $a$. The positive
interval argument above retains the best coefficient supplied by
an envelope; it does not assert sharpness for actual Collatz orbits.
The finite identities explain the common center and carry the exact
weights into its next analytic stage. They do not yet replace the
source's interior weighted-rotation estimate or its exterior
coefficient-tail argument, and do not establish a new final
natural-density exponent.

\subsection{The shifted tube and the weighted phase sum}
\label{sec:Mazur-weighted-phase}

The remaining phase calculation has two distinct sources of error.
First, moving the integer level from $s_\nu$ to $s_\nu+1$ changes
both the mass and the set of permitted times. Second, the phases in
the original, unequally weighted sum need not have their integral
distribution. We give finite formulas for each effect. These formulas
retain the original masses; no probability normalization is imposed.

The source arguments are Mazur~\cite{Mazur2026v2},
\texttt{ND/Band/A5ShiftedRawComparison.lean},
\texttt{A5Regularity.lean}, and
\texttt{ND/Discrepancy/A5Padding.lean},
\texttt{A5PaddedProfile.lean}, \texttt{A5TwoProfileTail.lean}.
Their endpoint discrepancy majorant is defined in
\texttt{PhaseToDbar.lean}. All refer to the pinned source commit.
The recurrences and padding mechanism below reconstruct those
arguments, with exact finite terms retained.

\paragraph{The two supports.}
Keep $\gamma=\log_2 3$, $\delta=2-\gamma$, and $p_\nu,s_\nu,q_0,q_1$
from Section~\ref{sec:Mazur-common-profile}. Put
\[
 R=\lceil W\rceil-1,\qquad
 d_\nu=s_\nu-2\nu=\lfloor A-\delta\nu\rfloor+1 .
\]
Thus $R$ is a nonnegative integer. Define finite sets of positive integers
\[
 I_0=\{\nu:-R\le d_\nu\le R\},\qquad
 I_1=\{\nu:-R\le d_\nu+1\le R\}.
\]
These are the indicator supports, even where an endpoint mass vanishes.

\begin{proposition}[exact support change and signed shift]
\label{prop:Mazur-shift-support}
The common part, lost part and gained part are respectively
\[
 I_c=\{\nu:-R\le d_\nu\le R-1\},\quad
 I_-=\{\nu:d_\nu=R\},\quad I_+=\{\nu:d_\nu=-R-1\}.
\]
For every integer $r$,
\begin{equation}
 \{\nu\ge1:d_\nu=r\}
 =\mathbb Z_{\ge1}\cap
       \left(\frac{A-r}{\delta},\frac{A-r+1}{\delta}\right].
 \label{eq:Mazur-displacement-fibre}
\end{equation}
In particular $\#I_-,\#I_+\le3$. Suppose $I_0\cup I_1$ is nonempty
and its least element $\ell$ satisfies $\ell>R+1$. For every real
function $g$ on this union,
\begin{align}
 \sum_\nu g_\nu(q_1(\nu)-q_0(\nu))
 &=-\sum_{\nu\in I_c}g_\nu q_0(\nu)
       \frac{d_\nu+2}{2(\nu+d_\nu+1)}
       -\sum_{\nu\in I_-}g_\nu q_0(\nu)
       +\sum_{\nu\in I_+}g_\nu q_1(\nu).
 \label{eq:Mazur-shift-signed}
\end{align}
Define the exact nonnegative quantity
\[
 \mathcal E=
 \sum_{I_c}q_0(\nu)\frac{|d_\nu+2|}{2(\nu+d_\nu+1)}
       +\sum_{I_-}q_0(\nu)+\sum_{I_+}q_1(\nu).
\]
It equals $\sum_\nu|q_1(\nu)-q_0(\nu)|$. For $R\ge1$, put
\[
 c_R(\ell)=\max\left\{
 \frac{(R-2)_+}{2(\ell-R+1)},
 \frac{R+1}{2(\ell+R)}\right\};
 \qquad c_0(\ell)=0 .
\]
Writing $Z=\sum_{I_0}q_0(\nu)$, we have the finite bound
\begin{equation}
 \mathcal E\le c_R(\ell)Z+
       \sum_{I_-}q_0(\nu)+\sum_{I_+}q_1(\nu)
       \le c_R(\ell)Z+\frac3{\sqrt\ell}.
 \label{eq:Mazur-shift-envelope}
\end{equation}
\end{proposition}
\begin{proof}
For an integer $d$, $|d|<W$ is equivalent to $-R\le d\le R$.
Intersection and set subtraction give the three displayed supports.
The floor inequalities $r-1\le A-\delta\nu<r$ give
\eqref{eq:Mazur-displacement-fibre}, with its strict left endpoint.
The inequalities $1/3<\delta<1$ follow from
$(4/3)^3>2$ and $1<4/3<2$. Hence the interval length $1/\delta$
is less than three and cannot contain four integers.

On $I_c$, $s_\nu\ge\nu$ and the adjacent binomial ratio gives
\[
 \frac{q_1(\nu)}{q_0(\nu)}-1
 =\frac{2\nu+d_\nu}{2(\nu+d_\nu+1)}-1
 =-\frac{d_\nu+2}{2(\nu+d_\nu+1)} .
\]
Our hypothesis on $\ell$ makes these denominators and both masses
positive. On the other two parts, precisely one indicator is zero.
This proves the signed identity and the exact expression for
$\mathcal E$. If $R=0$, the common part is empty.
Otherwise the function $|d+2|/[2(\nu+d+1)]$, for
$-R\le d\le R-1$, decreases up to $d=-2$ and increases thereafter:
the derivative of $(d+2)/(\nu+d+1)$ is
$(\nu-1)/(\nu+d+1)^2>0$. Its maximum is therefore bounded by the
two endpoint expressions in $c_R(\nu)$, each nonincreasing in $\nu$.
Finally Proposition~\ref{prop:Mazur-profile-envelope} proves
$p_\nu(s)\le1/(2\sqrt\nu)$ for every integer $s$.
There are at most six boundary indices, proving the last bound.
The exact boundary sums, rather than this ceiling, remain available.
\end{proof}

For the harmonic gate take $g_\nu=\mathbf1_{\{\theta_\nu>b\}}$;
for the flat gate take
$g_\nu=\mathbf1_{\{\theta_\nu>b\}}2^{2-\theta_\nu}$.
Their absolute corrections are at most $\mathcal E$ and
$e^\beta\mathcal E$, respectively. The signed expression
\eqref{eq:Mazur-shift-signed} retains cancellation, including the
zero common-part correction at $d_\nu=-2$.

\paragraph{Variation along the original time coordinate.}
Assume $I_0$ is nonempty, write $l=\min I_0$, $h=\max I_0$,
and suppose $l>R$. The floor $d_\nu$ decreases by zero or one at
each step, since $0<\delta<1$. Thus $I_0=[l,h]\cap\mathbb Z$.
Write $q_\nu=q_0(\nu)$, $Z=\sum_{l}^{h}q_\nu$ and
\[
 \mathcal V=q_l+q_h+\sum_{\nu=l}^{h-1}|q_{\nu+1}-q_\nu|.
\]
Both boundary terms occur, also when $l=h$.

\begin{proposition}[raw recurrences and full variation]
\label{prop:Mazur-raw-variation}
For $l\le\nu<h$, $s_{\nu+1}-s_\nu$ is one or two, and
\begin{equation}
 \frac{q_{\nu+1}}{q_\nu}-1=
 \begin{cases}
  d_\nu/(2\nu),&s_{\nu+1}=s_\nu+1,\\[2mm]
  (d_\nu^2+d_\nu-2\nu)/(4\nu(\nu+d_\nu+1)),
       &s_{\nu+1}=s_\nu+2 .
 \end{cases}
 \label{eq:Mazur-time-ratios}
\end{equation}
In particular, with
\[
 \eta(l,R)=\max\left\{\frac R{2l},
       \frac{R^2+R+2l}{4l(l-R+1)}\right\},
\]
one has
\begin{equation}
 \mathcal V\le q_l+q_h+\eta(l,R)Z
 \le\frac1{2\sqrt l}+\frac1{2\sqrt h}+\eta(l,R)Z .
 \label{eq:Mazur-raw-variation}
\end{equation}
Under the explicit additional scalar guards
\[
 W>1,\quad W\le l/8,\quad Z\ge1/32,\quad
 W\ge(37/20)\sqrt l ,
\]
the same raw variation satisfies $\mathcal V\le20ZW/l$.
\end{proposition}
\begin{proof}
Consecutive levels rise by one or two because $1<\gamma<2$.
Cancelling factorials in the two cases gives ratios
$s_\nu/(2\nu)$ and
$s_\nu(s_\nu+1)/[4\nu(s_\nu+1-\nu)]$, respectively.
Substitution of $s_\nu=2\nu+d_\nu$ proves
\eqref{eq:Mazur-time-ratios}.
The first absolute defect is at most $R/(2l)$.
The second is at most
\[
 \frac{R^2+R}{4\nu(\nu-R+1)}
       +\frac1{2(\nu-R+1)}\le
 \frac{R^2+R+2l}{4l(l-R+1)}.
\]
Multiplication by $q_\nu$ and summation proves the first bound;
the atom ceiling proves the second.

For the last statement it is enough to retain the source's
coarser bound $2W/\nu$ on each relative defect.
Here is a direct check. The first defect is less than $W/(2\nu)$.
For the second, $|d_\nu|<W\le\nu/8$, $\nu>8$, and $W>1$
give
\[
 |d_\nu^2+d_\nu-2\nu|
 \le W^2+W+2\nu\le(9/4)W\nu,\qquad
 4\nu(\nu+d_\nu+1)\ge(7/2)\nu^2 .
\]
The ratio is therefore at most $(9/14)W/\nu<2W/\nu$.
The internal variation is at most $2WZ/l$.
The two endpoint masses sum to at most $1/\sqrt l$,
which is at most $32Z/\sqrt l\le(640/37)ZW/l<18ZW/l$.
Adding proves the constant 20, without replacing $Z$ by one.
\end{proof}

\paragraph{An exact padding map and its error formula.}
Put $s=h-l+1$ and $u_k=q_{l+k}$ for $0\le k<s$.
For any integer $P\ge1$, pad on the left by $P$ zero entries:
\[
 \omega_j=0\ (0\le j<P),\qquad
 \omega_{P+k}=u_k\ (0\le k<s),\qquad
 \phi'=A+(l-P)\gamma\pmod1 .
\]
The map $\mathbb R^s\to\mathbb R^{P+s}$ so defined is linear,
injective, and preserves the sum and the $\ell^1$ norm.
Its inverse on its image discards the zero head.
The circle translation is invertible as well. In particular,
\[
 \{\phi'+(P+k)\gamma\}=\{A+(l+k)\gamma\}.
\]
Negative $l-P$ causes no loss: this is an integer circle translation,
not truncated subtraction of natural numbers.

\begin{proposition}[weighted phase error with both endpoints retained]
\label{prop:Mazur-raw-abel}
For $\sigma=H,F$, let $\mu_\sigma=\int_0^1f_\sigma$ and define
\[
 S_m^\sigma=\sum_{j=0}^{m-1}
       \big(f_\sigma(\{\phi'+j\gamma\})-\mu_\sigma\big).
\]
The exact weighted error is
\begin{align}
 \sum_{\nu=l}^h q_\nu f_\sigma(\theta_\nu)-\mu_\sigma Z
 =u_{s-1}S_{P+s}^\sigma-u_0S_P^\sigma
       +\sum_{k=0}^{s-2}(u_k-u_{k+1})S_{P+k+1}^\sigma .
 \label{eq:Mazur-raw-abel}
\end{align}
Define the actual weak endpoint discrepancies
\[
 D_m=\sup_{0\le t\le1}
 \left|\frac1m\#\{0\le j<m:\{\phi'+j\gamma\}\le t\}-t\right|,
 \qquad D_*=\max_{P\le m\le P+s}D_m ,
\]
and the exact Abel coefficient
\[
 \mathcal B_P=P|u_0|+(P+s)|u_{s-1}|
       +\sum_{k=0}^{s-2}(P+k+1)|u_k-u_{k+1}|.
\]
Then
\begin{equation}
 \left|\sum_{\nu=l}^h q_\nu f_\sigma(\theta_\nu)-\mu_\sigma Z\right|
 \le C_\sigma D_*\mathcal B_P,\qquad
 C_H=1,\quad C_F=2e^\beta-1,\qquad
 \mathcal B_P\le(P+s)\mathcal V .
 \label{eq:Mazur-raw-abel-bound}
\end{equation}
All quantities are finite and defined from the original phase orbit.
\end{proposition}
\begin{proof}
Write each centered summand as $S_{j+1}^\sigma-S_j^\sigma$,
multiply by $\omega_j$, and telescope. Only the jump at $P$,
the internal differences and the terminal term survive, giving
\eqref{eq:Mazur-raw-abel}. This also proves the identity for $s=1$.

Let $E_m(t)=\#\{j<m:\{\phi'+j\gamma\}\le t\}-mt$.
For $f_H$, the strict upper gate gives $S_m^H=-E_m(b)$.
For $f_F$, the jump immediately to the right of $b$ has size $e^\beta$.
On the two open pieces its derivative is $-\log2\,f_F(t)$.
Integrating the identity
$f_F(x)=f_F(1)-\int_x^1 f_F'(t)\,dt
-e^\beta\mathbf1_{\{x\le b\}}$ for $0\le x<1$
and then summing it at the finite sample points yields
\[
 S_m^F=-\int_0^1 E_m(t)f_F'(t)\,dt-e^\beta E_m(b).
\]
Here $f_F(1)$ denotes its left limit, namely $3$.
The total integral of $|f_F'|$ is
$(2-e^\beta/2)+(3e^\beta/2-3)=e^\beta-1$.
Thus $|S_m^\sigma|\le C_\sigma mD_m$, with the weak
endpoint convention accounting exactly for a sample at $b$.
Inserting these bounds into the exact identity proves
\eqref{eq:Mazur-raw-abel-bound}.
Each length coefficient in $\mathcal B_P$ is at most $P+s$,
which proves its bound by the full variation.
\end{proof}

This calculation explains a constant already available inside the
source. For its equal padding $P=s$, a nonincreasing uniform
endpoint majorant $D(m)$ gives $D_*\le D(s)$.
Its proved constant-20 variation estimate therefore gives
$40\,ZsW D(s)/l$ for the harmonic error, and
$(2e^\beta-1)$ times this for the flat error.
The later source adapter instead relaxes 20 to 32 and consequently
uses 64. This is a direct consequence of its stronger lemma, not
a new equidistribution theorem. The source uses
$D(m)=\min\{1,(6+2/(\pi c))m^{-1/\kappa}\}$ under its
specified Diophantine input. The finite proof above uses the
actual $D_m$ and does not certify that analytic input by computation.

For the full intrinsic two-cell sums $\mathcal R_H,\mathcal R_F$
over $I_0\cup I_1$, the exact decomposition now reads
\[
 \mathcal R_\sigma-\mu_\sigma/\delta
 =\left(\sum_{I_0}q_0f_\sigma-\mu_\sigma Z\right)
       +\sum_{I_0\cup I_1}g_\nu^\sigma(q_1-q_0)
       +\mu_\sigma(Z-1/\delta).
\]
Consequently the three respective absolute charges are
$C_\sigma D_*\mathcal B_P$,
$\mathcal E$ or $e^\beta\mathcal E$,
and $\mu_\sigma|Z-1/\delta|$.
The raw-mass error is still present. These intrinsic sums equal
the original physical sums only when their two supports lie in
the physical time window. For an arbitrary physical window, insert
its indicator into every displayed finite sum; neither missing
times nor their weights can be dropped. The source's interior
containment, raw-mass asymptotic, and exterior coefficient tail
remain separate parts of the analytic reconstruction.

\subsection{Raw mass, Gaussian coordinates, and physical interior profiles}
\label{sec:Mazur-mass-interior}

The total weight in the tube is not one. Its limit is $1/\delta$,
and this factor is essential to both raw phase centers. We now prove
the mass estimate with explicit finite errors and then evaluate the
original physical sums. This connects Mazur's A6 calculation to the
Gaussian envelope in the Gaussian-envelope argument in Section~\ref{subsec:Allikvere-phase-average}; the envelope here
is exactly half of that earlier function. The truncation scales differ,
so their tail estimates must not be identified.

The source comparison is with Mazur~\cite{Mazur2026v2},
\texttt{ND/Band/A6PhysicalRate.lean},
\texttt{A6PhysicalInteriorRate.lean}, \texttt{A5Physical.lean},
and \texttt{ND/Discrepancy/A5TwoProfilePhysicalA6.lean}, at the pinned
commit. The factorial estimate used below is Robbins's original
formula~\cite[(1)--(2), p.~26]{Robbins1955}, not an asymptotic formula
without its remainder.

\begin{proposition}[an exact atom formula and a finite logarithmic error]
\label{prop:Mazur-atom-gaussian}
For integers $\nu\ge1$, $k$ with $|k|\le\nu/2$, put $t=k/\nu$ and
\[
 I(t)=(1+t)\log(1+t)-(2+t)\log(1+t/2).
\]
Let $r_n$ be defined by
$n!=\sqrt{2\pi n}(n/e)^n e^{r_n}$. Then
\begin{equation}
 p_\nu(2\nu+k)=\frac{e^{-\nu I(t)+r_{2\nu+k}-r_\nu-r_{\nu+k}}}
 {2\sqrt{\pi\nu}\sqrt{(1+t/2)(1+t)}}.
 \label{eq:Mazur-exact-atom}
\end{equation}
For $G_\nu(k)=(2\sqrt{\pi\nu})^{-1}e^{-k^2/(4\nu)}$,
\begin{equation}
 \left|\log\frac{p_\nu(2\nu+k)}{G_\nu(k)}\right|
 \le\frac{4|k|}{3\nu}+\frac{16|k|^3}{27\nu^2}+\frac1{4\nu}.
 \label{eq:Mazur-atom-log-error}
\end{equation}
Now retain the strict tube $I_0$, $R=\lceil W\rceil-1$,
$d_\nu=\lfloor A-\delta\nu\rfloor+1$, and $l=\min I_0$.
Suppose $I_0$ is nonempty and $W\le l/2$. Put $\xi=A/\delta$ and
\[
 G_\xi(u)=\frac1{2\sqrt{\pi u}}
    \exp\!\left(-\frac{\delta^2(u-\xi)^2}{4u}\right)\quad(u>0),
 \qquad
 \varepsilon=\frac{11W}{6l}+\frac{16W^3}{27l^2}+\frac1{2l}.
\]
For every $\nu\in I_0$, $q_0(\nu)>0$ and
\begin{equation}
 e^{-\varepsilon}G_\xi(\nu)\le q_0(\nu)
                  \le e^\varepsilon G_\xi(\nu).
 \label{eq:Mazur-rounded-gaussian}
\end{equation}
The exact rounding term in its logarithm is
$-(2d_\nu e_\nu-e_\nu^2)/(4\nu)$, where
$e_\nu=1-\{A-\delta\nu\}\in(0,1]$.
\end{proposition}
\begin{proof}
Set $s=2\nu+k$ and $m=\nu+k\ge\nu/2>0$. The exact identity
$p_\nu(s)=(\nu/s)\binom{s}{\nu}2^{-s}$, followed by Robbins's
factorial formula at $s,\nu,m$, gives
\eqref{eq:Mazur-exact-atom}. In particular the factor $\nu/s$
is not removed. Robbins proves
$1/(12n+1)<r_n<1/(12n)$.
Since $s\ge\nu+1$,
$r_s<1/(12s)<1/(12\nu+1)<r_\nu$. Thus
\[
 -\frac1{4\nu}<-(r_\nu+r_m)<r_s-r_\nu-r_m<0.
\]
On $[-1/2,1/2]$ the derivative of
$-\tfrac12\log(1+t/2)-\tfrac12\log(1+t)$ has absolute value
at most $1/3+1=4/3$.
Also $I(0)=I'(0)=0$, $I''(0)=1/2$, and
\[
 I''(t)=\frac1{(1+t)(2+t)},\qquad
 |I'''(t)|=\frac1{(1+t)^2}-\frac1{(2+t)^2}\le\frac{32}{9}.
\]
The last expression decreases with $t$. Taylor's integral remainder
therefore gives $|I(t)-t^2/4|\le16|t|^3/27$.
Taking logarithms in the exact identity proves the first bound.
This is the same rate function $I$ used in
Proposition~\ref{prop:Mazur-maximal-tail}, now in a local atom formula.

For the tube, $|d_\nu|<W\le\nu/2$ and
$d_\nu=A-\delta\nu+e_\nu$. Hence
$d_\nu^2-(A-\delta\nu)^2=2d_\nu e_\nu-e_\nu^2$,
whose absolute value is at most $2W+1$.
Adding this divided by $4\nu$ to
\eqref{eq:Mazur-atom-log-error}, and replacing $\nu$ by its lower
bound $l$, gives precisely $\varepsilon$. At integral phase
$e_\nu=1$, not zero; the same identity and estimate apply.
\end{proof}

\begin{proposition}[finite tube mass and its uniform asymptotic]
\label{prop:Mazur-tube-mass}
Assume $R\ge1$, $\xi>R/\delta$, and $W\le l/2$. Set
\[
 a=\xi-\frac R\delta,\quad b=\xi+\frac{R+1}{\delta},\quad
 t_*=\sqrt{\xi^2+\delta^{-4}}-\delta^{-2},\quad
 M_\xi=G_\xi(t_*),
\]
\[
 u_-=\frac{R}{\delta\sqrt a},\qquad
 u_+=\frac{R+1}{\delta\sqrt b},\qquad
 T=\frac{2}{\delta^2\sqrt\pi}
   \left(\frac{e^{-\delta^2u_-^2/4}}{u_-}
              +\frac{e^{-\delta^2u_+^2/4}}{u_+}\right),\quad
 Q=2M_\xi+T.
\]
The strict tube is exactly $(a,b]\cap\mathbb Z$, and its original
mass $Z=\sum_{I_0}q_0(\nu)$ satisfies
\begin{equation}
 \left|Z-\frac1\delta\right|
 \le (e^\varepsilon-1)\left(\frac1\delta+Q\right)+Q.
 \label{eq:Mazur-finite-mass}
\end{equation}
In particular, fix $C\ge1/2$. As $N\to\infty$, for
\[
 W=C\sqrt{N\log N},\qquad \frac N{43}\le\xi\le\frac N{13},
\]
all the preceding guards hold eventually and, uniformly over every
real center in this interval,
\begin{equation}
 Z=\frac1\delta+
 O_C\!\left(\frac{(\log N)^{3/2}}{\sqrt N}\right).
 \label{eq:Mazur-mass-asymptotic}
\end{equation}
Here $C$ is fixed before $N$ tends to infinity.
\end{proposition}
\begin{proof}
The floor inequalities in
\eqref{eq:Mazur-displacement-fibre} give the exact interval $(a,b]$.
Its length $(2R+1)/\delta>1$ ensures nonemptiness; $a>0$ ensures
positive indices. the Gaussian-envelope argument in Section~\ref{subsec:Allikvere-phase-average}, with $z=\xi$,
gives $\int_0^\infty G_\xi=1/\delta$ and the unique maximum at
$t_*$. The zero-extended function $G_\xi\mathbf1_{(a,b]}$ is
unimodal, with total variation at most $2M_\xi$, including its two
endpoint jumps. For each integer cell, the difference between its
left-endpoint value and its integral is bounded by that cell's
variation. Summing over cells gives
\[
 \left|\sum_{\nu\in I_0}G_\xi(\nu)-\int_a^bG_\xi(u)\,du\right|
 \le 2M_\xi.
\]
This estimate respects an integral $a$ or $b$ with the displayed
strict/weak conventions.

For the continuous tails use the explicit increasing bijection
\[
 v=\sqrt u-\xi/\sqrt u,\qquad
 u=\left(\frac{v+\sqrt{v^2+4\xi}}2\right)^2.
\]
It sends $(0,\infty)$ onto $\mathbb R$ and $a,b$ to $-u_-,u_+$.
Its exact density transformation is
\[
 G_\xi(u)\,du=\frac{e^{-\delta^2v^2/4}}{2\sqrt\pi}
       \left(1+\frac v{\sqrt{v^2+4\xi}}\right)dv.
\]
The parenthesis lies strictly between zero and two.
For $u_0>0$, integration of
$(v/u_0)e^{-\delta^2v^2/4}$ on $[u_0,\infty)$ proves
\[
 \int_{u_0}^\infty e^{-\delta^2v^2/4}dv
 \le\frac{2e^{-\delta^2u_0^2/4}}{\delta^2u_0}.
\]
Thus the two discarded continuous tails sum to at most $T$.
Writing $Z_G=\sum_{I_0}G_\xi(\nu)$, we have
$|Z_G-1/\delta|\le Q$. By
\eqref{eq:Mazur-rounded-gaussian},
$|Z-Z_G|\le(e^\varepsilon-1)Z_G$.
These two bounds prove \eqref{eq:Mazur-finite-mass}, including
its product error.

For the asymptotic, $W=o_C(N)$, $R=W+O(1)$ and
$l=\xi+O_C(W+1)\asymp N$, uniformly in the stated center interval.
In particular $a>0$, $l>2W$ and $R\ge1$ eventually.
The explicit expression for $\varepsilon$ is
$O_C((\log N)^{3/2}/\sqrt N)=o_C(1)$.
Also $t_*\ge\xi/2$ eventually, so
$M_\xi\le(2\pi\xi)^{-1/2}=O(N^{-1/2})$.
Both tail exponents satisfy, uniformly,
\[
 \frac{\delta^2u_\pm^2}{4}
 \ge\left(\frac{13C^2}{4}+o_C(1)\right)\log N
 \ge\frac23\log N
\]
for sufficiently large $N$, because $13C^2/4\ge13/16>2/3$.
Moreover $u_\pm\gg_C\sqrt{\log N}$. Consequently
$T=O_C(N^{-2/3}/\sqrt{\log N})$ and $Q=O_C(N^{-1/2})$.
Substitution in the finite bound proves the asserted rate.
These tails are polynomially small; the different cutoff in
the Gaussian-envelope argument in Section~\ref{subsec:Allikvere-phase-average} had smaller tails and was not used
to justify this truncation.
\end{proof}

\begin{corollary}[the two intrinsic centers from the actual raw weights]
\label{cor:Mazur-intrinsic-centers}
Under the asymptotic schedule of
Proposition~\ref{prop:Mazur-tube-mass}, let
$\log2<\beta<2\log2$, $\mu=8.616$ and $d=\log(4/3)$.
The literal two-cell sums over $I_0\cup I_1$ satisfy
\begin{align}
 \mathcal R_H&=\frac{\beta}{d}
       +O_C\!\left(N^{-1/(2\mu)}
                         (\log N)^{1/2-1/(2\mu)}\right),\\
 \mathcal R_F&=\frac{e^\beta-1}{d}
       +O_C\!\left(N^{-1/(2\mu)}
                         (\log N)^{1/2-1/(2\mu)}\right).
 \label{eq:Mazur-intrinsic-centers}
\end{align}
Both errors are uniform in the real center and in $\beta$.
In particular each is $O_{C,c}(N^{-c})$ for every
$0<c<1/17.232$.
\end{corollary}
\begin{proof}
The base support has length $s=O_C(W)$ and endpoints comparable to
$N$. Proposition~\ref{prop:Mazur-tube-mass} bounds $Z$ uniformly.
Put $w_\nu=G_\xi(\nu)\mathbf1_{I_0}(\nu)$. Its full variation
is at most $2M_\xi=O(N^{-1/2})$. Write the exact difference
\[
 \sum_{I_0}q_0(f_\sigma-\mu_\sigma)
 =\sum_{I_0}w(f_\sigma-\mu_\sigma)
       +\sum_{I_0}(q_0-w)(f_\sigma-\mu_\sigma).
\]
The second sum is bounded in absolute value by
$\|f_\sigma-\mu_\sigma\|_\infty(e^\varepsilon-1)(1/\delta+Q)$.
Thus every original atom is retained with an explicit error;
no equality between it and the Gaussian is assumed.
Use the Abel estimate in Section~\ref{subsec:Allikvere-phase-average}, whose Diophantine input is
Rhin~\cite[p.~160, (8)]{Rhin1987} and whose discrepancy inequalities
are those of Kuipers--Niederreiter~\cite[Ch.~2]{KN1974}.
As periodic functions, $f_H,f_F$ have variations $2$ and
$2e^\beta\le8$. These include the wrap-around jump at zero;
the interval coefficient $2e^\beta-1$ in the exact weak-discrepancy
formula is a different quantity. The finite bound on the first sum
is $C_\gamma\operatorname{Var}(f_\sigma)s^{1-1/\mu}2M_\xi$.
The total centered weighted error is consequently at most
\[
 O_C\!\left(s^{1-1/\mu}N^{-1/2}
                     +(\log N)^{3/2}N^{-1/2}\right)
 =O_C\!\left(N^{-1/(2\mu)}
                   (\log N)^{1/2-1/(2\mu)}\right).
\]
Both centered functions have uniformly bounded sup norm.
No division by $Z$ occurs. Using only the raw-variation ceiling
would instead give the larger logarithmic power $1-1/(2\mu)$;
the explicit Gaussian comparison saves a factor $\sqrt{\log N}$.

For the shift correction, $\ell=\min(I_0\cup I_1)\asymp N$ and
$R=O_C(W)$, so \eqref{eq:Mazur-shift-envelope} gives
$\mathcal E=O_C(\sqrt{\log N/N})$. Its flat multiplier $e^\beta$
is bounded by four. Finally both phase means are bounded, and
\eqref{eq:Mazur-mass-asymptotic} bounds the raw-mass correction
by $O_C((\log N)^{3/2}/\sqrt N)$.
Insert these three estimates in the exact decomposition at the end
of Section~\ref{sec:Mazur-weighted-phase}. The last two errors are
smaller than the first. The means in
\eqref{eq:Mazur-phase-means}, together with $\delta=d/\log2$,
give precisely the two displayed centers. The stated weaker powers
follow since a fixed power of $\log N$ is $o(N^\eta)$ for every
$\eta>0$.
\end{proof}

\begin{proposition}[transport to the original physical window]
\label{prop:Mazur-physical-centers}
Let $B,z,M>0$, $m_0\in\mathbb Z_{\ge0}$,
$\tau=\log M-\log B-m_0d$, $A=\log(z/M)/\log2$,
$t_0=\log(z/B)/d$, and $\log2<\beta<2\log2$.
For $W>0$, suppose
\begin{equation}
 \frac{|\tau|}{d}+\frac{W+1}{\delta}\le3W.
 \label{eq:Mazur-directional-margin}
\end{equation}
The map $\nu\mapsto n=\nu+m_0$ embeds $I_0\cup I_1$ in
\[
 J=\mathbb Z_{\ge0}\cap
       [\,t_0-3W,\ t_0+\beta/d+3W\,].
\]
It is bijective onto its image with inverse $n\mapsto n-m_0$;
all image points have $n>m_0$. Zero extension gives exact equality
between each intrinsic two-cell sum and its original physical sum
on $\{\nu>0:\nu+m_0\in J\}$.

For the original source schedule
\[
 \alpha=1001/1000,\quad L=\log B,\quad
 N=\lfloor L/(10\log2)\rfloor,\quad m_0=\lfloor L/100000\rfloor,
 \quad W=C\sqrt{N\log N},
\]
fix $C\ge1/2$, let $y=B^\alpha$ or $B^{\alpha^2}$, and put
\[
 K=\lfloor(\alpha-1)\log y\rfloor,\quad
 \beta=(\alpha-1)\log y/K,\quad z=y e^{j\beta}\quad(0\le j<K).
\]
Uniformly in the branch, $j$, and every $M>0$ with
$|\tau|\le4W/25$, for sufficiently large $B$,
\begin{align}
 R_H(M)&=\beta/d+
 O_C\!\left(N^{-1/17.232}(\log N)^{1/2-1/17.232}\right),\\
 R_F(M)&=(e^\beta-1)/d+
 O_C\!\left(N^{-1/17.232}(\log N)^{1/2-1/17.232}\right).
 \label{eq:Mazur-physical-centers}
\end{align}
These are interior estimates on the original physical carrier.
\end{proposition}
\begin{proof}
The exact affine identity is
\[
 \xi=A/\delta=t_0-m_0-\tau/d.
\]
The union of the two indicator supports consists exactly of the
positive integers in
$((A-R)/\delta,(A+R+2)/\delta]$. Thus every image point obeys
\[
 n>t_0-\tau/d-R/\delta,\qquad
 n\le t_0-\tau/d+(R+2)/\delta.
\]
The first inequality and the margin put $n$ above $t_0-3W$.
For the upper endpoint, $\beta/d>1/\delta$ and $R<W$ give
\[
 |\tau|/d+(R+2)/\delta
 <\beta/d+|\tau|/d+(W+1)/\delta
 \le\beta/d+3W.
\]
This proves containment of both supports without asserting that they
are equal. Translation is injective, its displayed inverse is exact
on its image, and all summands outside the union vanish.
This proves the equality of the finite sums, not just an estimate
of discarded endpoints. In particular the one-unit shift is absorbed
by the upper band width; a symmetric enlargement by that unit is
unnecessary for this containment argument.

For the source schedule, $K\to\infty$ and $\beta\to1$ uniformly in
the two branches; thus the required interval for $\beta$ holds
eventually. Since $z\in[y,y^\alpha)$, the exact identity for $\xi$
gives uniformly
\[
 \frac{\alpha-1}{d}-\frac1{100000}+o_C(1)
 \le\frac{\xi}{L}
 \le\frac{\alpha^3-1}{d}+o_C(1).
\]
Here the floor error is at most $1/L$, and
$|\tau|/L\le4W/(25L)=o_C(1)$. Since $N/L\to1/(10\log2)$,
these bounds put $\xi$ in $[N/43,N/13]$ eventually. For an explicit
check of the strict room, use
\[
 0.6931<\log2<0.6932,\qquad
 0.2876<d<0.2877,
\]
which give
\[
 10(0.6931)\left(\frac{0.001}{0.2877}-0.00001\right)>1/43,
 \qquad
 10(0.6932)\frac{(1001/1000)^3-1}{0.2876}<1/13.
\]
The logarithm intervals follow from
$\log x=2\sum_{k=0}^{J-1}u^{2k+1}/(2k+1)+E_J$,
$u=(x-1)/(x+1)$, with
$0<E_J<2u^{2J+1}/((2J+1)(1-u^2))$:
take $J=12$ for $x=2$ and $J=6$ for $x=4/3$.
This also gives
\[
 \frac{4}{25d}+\frac1\delta
 =\frac{4/25+\log2}{d}
 <\frac{0.16+0.6932}{0.2876}<3.
\]
Consequently the margin holds for all sufficiently large $W$,
uniformly under $|\tau|\le4W/25$. All required guards of the mass
and intrinsic-profile results therefore hold with $C$ fixed before
$B\to\infty$. Exact transport proves the physical assertions.
\end{proof}

The same band denominators from
Proposition~\ref{prop:Mazur-band-denominators} now give
$R_H/H=2/d+O_{C,c}(N^{-c})$ and
$(z/D)R_F=2/d+O_{C,c}(N^{-c})$ throughout the interior,
for $0<c<1/17.232$. The signed denominator errors and the exact
kernel residual in Proposition~\ref{prop:Mazur-profile-residual}
are unchanged. The exterior coefficient estimate and full kernel
summation are proved in Section~\ref{sec:Mazur-exterior-assembly} below.
Mazur's source uses its stated $\kappa=14.3$
phase input; the rate here uses the independently cited Rhin bound,
already used in the Allikvere reconstruction. No improved final
Mazur natural-density theorem or local Lean replay is asserted.

\subsection{Exterior endpoint weights and the full reference profile}
\label{sec:Mazur-exterior-assembly}

The interior estimate must be summed against the endpoint kernel, not
against a probability measure obtained by dividing that kernel by its
mass. The remaining issue is therefore quantitative: how much of the
original weighted endpoint set lies outside the interior? The native
reciprocal-fibre estimate contains a linear factor in the last valuation.
Summing that factor exactly, rather than bounding it by a multiple of
$\log B$, gives a substantially smaller exterior error.

The source route is Mazur~\cite{Mazur2026v2},
\texttt{Tao/Section5/EndpointValuationFiber.lean},
\texttt{EndpointRatio.lean}, \texttt{EndpointLowRoom.lean},
\texttt{EndpointFiberEnvelope.lean}, \texttt{EndpointCoefficientBound.lean},
and \texttt{Geom2HighWeightScalar.lean};
then \texttt{ND/Band/A5ReferenceFilteredCoefficient.lean},
\texttt{A5ReferenceExteriorTupleTail.lean},
\texttt{A5ReferenceTerminalCap.lean}, and
\texttt{ND/Probability/Geom2FiniteTwoEventTail.lean}.
The final comparison is with
\texttt{ND/Discrepancy/A5ReferenceFullCenterRate.lean}.
These are the pinned v2 sources, not a replay of their formal build.
The geometric law and first-passage mechanism originate in
Tao~\cite{Tao2026v7}, the definition of ``Geometric random variable''
(v7 source label \texttt{geom}) and the Section 5 coefficient argument
(source label \texttt{loam}).
We reconstruct the needed finite estimates before taking limits.

\paragraph{Original endpoints and their exact fibres.}
Write $h=\lfloor L/100000\rfloor$, $L=\log B$,
$N=\lfloor L/(10\log2)\rfloor$, and $d=\log(4/3)$.
For a finite set $S$ of positive odd integers with first passage at time
$h$ under the odd Syracuse map $T$, define
\[
 c_q(S,x)=3^q\sum_{\substack{M\in S\\M\equiv x\pmod{3^q}}}\frac1M .
\]
In the application $S$ is the original $\mathcal E(E)$ or its
interior/exterior subset; its rounded lost-window endpoints remain
unchanged. For the chronological valuation word
$b=(b_1,\ldots,b_h)$, put
\[
 s=\sum_{i=1}^h b_i,\quad s_-=\sum_{i=1}^{h-1}b_i,\quad a=b_h,
 \quad w(b)=2^{-s},\quad b_0(b)=(1+a/2)2^{-s}.
\]
The fibre $F_{b,x}$ consists of precisely the $M\in S$ with word $b$
and residue $x$. Only nonempty fibres enter the finite word carrier.
Thus $c_q(S,x)=\sum_b3^q\sum_{F_{b,x}}1/M$ exactly.
The map $M\mapsto(b(M),M)$ is a bijection from the residue-restricted
endpoint set to the disjoint union of these fibres, with inverse
$(b,M)\mapsto M$. Projection to the word alone is generally many-to-one;
its pushforward retains the reciprocal sum in each fibre.

\begin{proposition}[native endpoint coefficient with its decaying remainder]
\label{prop:Mazur-native-coefficient}
Suppose $B>1$, $h\ge1$, $2\,3^{h-1}\le B$, and $h+q\le N$.
For every such finite $S$ and every $x\in\mathbb Z/3^q\mathbb Z$,
\begin{equation}
 c_q(S,x)\le\sum_{b:F_{b,x}\ne\varnothing}b_0(b)
             +\epsilon_B,\qquad
 \epsilon_B=3B^{-4997/600000}.
 \label{eq:Mazur-native-coefficient}
\end{equation}
In particular $0\le c_q(S,x)\le2+\epsilon_B$.
The assertion holds for filtered original endpoint sets without
requiring the enlarged same-word fibres to obey the filter.
\end{proposition}
\begin{proof}
The exact affine iteration has positive offset
\[
 T^jM=3^j2^{-s_j}M+F_j,\qquad
 F_j=\sum_{i=1}^j\frac{3^{j-i}}{2^{b_i+\cdots+b_j}}
 \le(3/2)^j-1<3^j .
\]
This includes $j=0$ with $F_0=0$.
Since $T^{h-1}M>B$ and $B\ge2\,3^{h-1}$, the affine main term
at time $h-1$ exceeds $B/2$. At time $h$, positivity and
$T^hM\le B$ give
\[
 B2^{s_-}<2\,3^{h-1}M,\qquad 3^hM\le B2^s .
\]
Moreover $T^{h-1}M<2\,3^{h-1}2^{-s_-}M$ and
$2^a\le3T^{h-1}M+1\le4T^{h-1}M$, so
$2^s<8\,3^{h-1}M$. Consequently each fibre lies in the inclusive
integer interval $[A_0,U]$, and also in $[A_1,U]$, where
\[
 A_0=\left\lfloor\frac{B2^{s_-}}{2\,3^{h-1}}\right\rfloor+1,\quad
 U=\left\lfloor\frac{B2^s}{3^h}\right\rfloor,\quad
 A_1=\max\left\{A_0,\left\lfloor\frac{2^s}{8\,3^{h-1}}\right\rfloor+1\right\}.
\]
For a nonempty fibre these lower endpoints are at most $U$.
The strict quotient inequalities give $U<2^aA_0$, hence
$\log(U/A_i)\le a\log2\le a$.

The exact valuation cylinder is one odd class modulo $2^{s+1}$.
Indeed the one-step condition is $3M+1\equiv2^{b_1}\pmod{2^{b_1+1}}$.
Inductively the next prescribed odd residue lifts its unique previous
class to a unique class modulo the next power of two: solving
$3M+1=2^{b_1}M_1$ modulo $2^{s+1}$ is possible and unique since
$3$ is invertible modulo that power. This also proves the converse
by dividing each congruence by its prescribed power of two.
Combining this class with $M\equiv x\pmod{3^q}$ gives one class
modulo $Q=2^{s+1}3^q$, since the two moduli are coprime.
No endpoints or word multiplicities are identified in this step.

For any nonempty portion of a $Q$-progression in $[A,U]$, $A\ge1$,
decreasingness of $1/t$, integrated between successive points, gives
\[
 \sum\frac1M\le\frac1A+\frac1Q\log(U/A)
 \le\frac1A+\frac1Q\{1/A+\log(U/A)\}.
\]
To verify the first inequality, retain the first point separately;
for each later point $M$ its contribution is at most
$Q^{-1}\int_{M-Q}^M dt/t$. The intervals concatenate, and their
union lies in $[A,U]$. This also bounds every subset of the progression.

Call a word high when $2^s>B^{1/2}$, strictly. On the low branch,
$4^s\le B$. Also $2^{4(h+q)}\le B$, because $h+q\le N$.
Thus $2^s2^{2(h+q)}\le B$ by multiplying the two square bounds.
As $a\le s$ and $h\ge1$, this implies
$2^{a+2}3^{h+q-1}\le B$. Multiplication by $2^{s_-}$ shows
$Q\le\lfloor B2^{s_-}/(2\,3^{h-1})\rfloor<A_0$.
The retained progression estimate is therefore at most
\[
 3^q\{Q^{-1}+Q^{-1}(1+a)\}=(1+a/2)2^{-s}.
\]
On the high branch, use $A_1$. Its outer term satisfies
$3^q/A_1<(8/3)3^{h+q}2^{-s}\le3\,3^{h+q}2^{-s}$.
The remaining term is at most $(1+a)2^{-s-1}\le b_0(b)$.
This proves the pointwise envelope
$b_0(b)+3\,3^{h+q}2^{-s}\mathbf1_{\{2^s>B^{1/2}\}}$.

Here and below enlarging a finite word carrier to all positive words
of length $h$ is legitimate by nonnegativity, without duplicating words.
For independent $X_i$ with $\Pr(X_i=r)=2^{-r}$, $r\ge1$, the centered
one-coordinate moment generating function is
\[
 g(t)=\mathbb E e^{t(X_i-2)}
     =\frac{e^{-t}}{2-e^t},\qquad t<\log2.
\]
Since $h\le L/100000$ and $\log2<7/10$, the high event implies
$\sum_i(X_i-2)>(7/10)L$.
For $0\le t\le1/4$, $e^t<4/3$ and
$(\log g)''=2e^t/(2-e^t)^2\le6$.
Here $\log(4/3)>1/4$ follows directly by integrating $1/x$.
As $(\log g)(0)=(\log g)'(0)=0$, $\log g(1/4)\le3/16<1/2$.
Markov's inequality now bounds the high mass by
$e^{-7L/40+h/2}\le B^{-34999/200000}$.
Finally $3^{6(h+q)}\le2^{10(h+q)}\le B$, so the total high charge
is at most $\epsilon_B$. This proves the first assertion.

The elementary series $\sum_{a\ge1}2^{-a}=1$ and
$\sum_{a\ge1}a2^{-a}=2$ show
$\sum_{b\in\mathbb N_{>0}^h}b_0(b)=2$.
This proves the last bound, retaining the smaller $2+\epsilon_B$
rather than replacing the decaying remainder by $3$.
The proof never enlarged the original endpoint filter before retaining
its nonempty word carrier.
\end{proof}

\begin{proposition}[exact weighted two-event estimate]
\label{prop:Mazur-weighted-chernoff}
Let $h\ge2$, $u>0$, $0<v<h-1$, and let $\mathcal C$ be any
finite set of positive words of length $h$ such that every $b\in\mathcal C$
satisfies at least one of
\[
 s-2h>u,\qquad s_--2(h-1)<-v .
\]
With the rate function from Section~\ref{sec:Mazur-band-fullgood},
\[
 I(t)=(1+t)\log(1+t)-(2+t)\log(1+t/2)\quad(t>-1),
\]
one has
\begin{equation}
 \sum_{b\in\mathcal C}(1+b_h/2)2^{-|b|}
 \le \left(2+\frac{u}{2h}\right)e^{-hI(u/h)}
       +2e^{-(h-1)I(-v/(h-1))}.
 \label{eq:Mazur-weighted-chernoff}
\end{equation}
\end{proposition}
\begin{proof}
For $t<\log2$, summing the terminal factor exactly gives
\[
 \mathbb E[(1+X_h/2)e^{t\sum_{i=1}^h(X_i-2)}]
 =g(t)^h\left(1+\frac1{2-e^t}\right).
\]
This follows by differentiating the absolutely convergent geometric
series: $\mathbb E[X_he^{t(X_h-2)}]=2e^{-t}/(2-e^t)^2$.
The upper event is bounded, for $0<t<\log2$, by
$e^{-tu}g(t)^h(1+(2-e^t)^{-1})$.
Take
\[
 e^t=\frac{2(1+u/h)}{2+u/h}.
\]
This lies strictly between $1$ and $2$; direct substitution gives
$e^{-tu}g(t)^h=e^{-hI(u/h)}$ and
$1+(2-e^t)^{-1}=2+u/(2h)$.

For the lower event, the terminal coordinate is independent of the
first $h-1$ coordinates and its weighted expectation equals $2$.
Taking $e^t=2(1-v/(h-1))/(2-v/(h-1))\in(0,1)$ in the ordinary
lower-tail Markov estimate gives its displayed contribution.
Apply the union inequality to these two nonnegative weighted events.
It remains valid when they overlap; no independence of the two events
is asserted. Absolute convergence follows from $e^t<2$ and the
linear terminal factor. No terminal cutoff or factor $\log B+1$
has been introduced.
\end{proof}

\begin{corollary}[exterior coefficient on the original carrier]
\label{cor:Mazur-exterior-coefficient}
Retain Proposition~\ref{prop:Mazur-native-coefficient}'s schedule.
For $W>0$ put
\[
 u=\frac{4W}{25\log2},\qquad
 \kappa_8=\frac{\log(8/3)}{\log2},\qquad v=u-\kappa_8.
\]
Suppose $h\ge2$ and $0<v<h-1$. Let
$S_{\rm out}=\{M\in S:|\log M-L-hd|>4W/25\}$.
Then, uniformly in $x$,
\begin{equation}
 c_q(S_{\rm out},x)\le
 \left(2+\frac{u}{2h}\right)e^{-hI(u/h)}
 +2e^{-(h-1)I(-v/(h-1))}+\epsilon_B
 =:\mathcal T_B(W).
 \label{eq:Mazur-exterior-finite}
\end{equation}
For fixed $C>0$ and $W=C\sqrt{N\log N}$, all the guards hold eventually,
and the displayed majorant satisfies
\begin{equation}
 \mathcal T_B(W)=(4+o_C(1))N^{-A_C},\qquad
 A_C=\frac{64C^2}{(\log2)^3}.
 \label{eq:Mazur-exterior-rate}
\end{equation}
This is an asymptotic for the explicit majorant, not an equality for
the actual exterior coefficient.
\end{corollary}
\begin{proof}
The cleared first-passage inequalities in the preceding proof give
the exact corridor
\[
 (s_--2(h-1))\log2-\log(8/3)
 <\log M-L-hd\le(s-2h)\log2 .
\]
An exterior endpoint thus supplies either the strict upper event
$s-2h>u$ or the strict predecessor event $s_--2(h-1)<-v$.
Each retained word has such an endpoint witness, so the word satisfies
the event even after its reciprocal fibre is enlarged.
Apply \eqref{eq:Mazur-native-coefficient} and
\eqref{eq:Mazur-weighted-chernoff}. The exact buffer $\kappa_8$
has not been relaxed to the source's $3$.

For the limit, $h=L/100000+O(1)$ and
$N=L/(10\log2)+O(1)$, with $C$ fixed before $B$ grows.
In particular $u/h\to0$, $v/(h-1)\to0$, and $v>0$ eventually.
The proved derivative identity $I''(t)=1/((1+t)(2+t))$ gives
$I(t)=t^2/4+O(t^3)$ near zero. Therefore
\[
 hI(u/h)=\frac{u^2}{4h}+O(u^3/h^2)
        =A_C\log N+o_C(1),
\]
since $u^3/h^2=O_C((\log L)^{3/2}/\sqrt L)\to0$ and the floor
errors in $N/h$ contribute $O_C(\log L/L)$.
Replacing $u$ by $u-\kappa_8$ and $h$ by $h-1$ changes the
exponent by $O_C(\sqrt{\log L/L}+\log L/L)=o_C(1)$.
Both prefactors tend to $2$. Also
$\epsilon_B=o_C(N^{-A_C})$. This proves the stated asymptotic.
The room condition follows from
$(h-1)\log3+\log2=L\log3/100000+O(1)<L$.
\end{proof}

\begin{proposition}[kernel assembly without changing its mass]
\label{prop:Mazur-full-profile}
For the original schedule of Proposition~\ref{prop:Mazur-physical-centers},
fix $C\ge1/2$, put
$q=\lceil hL^{-2/5}\rceil$, and retain the original offset probability
$p_q$ and $K_q(M)=3^qp_q(M\bmod3^q)/M$.
For the original endpoint set $S=\mathcal E(E)$, write
$Z_q(S)=\sum_SK_q(M)$ and $a_*=2/d$.
Then
\[
 0\le Z_q(S)\le2+\epsilon_B,\qquad
 0\le Z_q(S_{\rm out})\le\mathcal T_B(W).
\]
If $e_\sigma$ is any valid uniform bound for the interior error
$|P_\sigma(M)-a_*|$, where $P_H(M)=R_H(M)/H$ and
$P_F(M)=(z/D)R_F(M)$, the exact finite estimate is
\begin{equation}
 \left|\sum_SK_q(M)P_\sigma(M)-a_*Z_q(S)\right|
 \le e_\sigma Z_q(S_{\rm in})
    +\max\{a_*,L_\sigma-a_*\}Z_q(S_{\rm out}).
 \label{eq:Mazur-full-finite}
\end{equation}
Here $L_H,L_F$ are the proved original-carrier envelopes of
Proposition~\ref{prop:Mazur-profile-envelope}.
For the two source branches, all bands, and all landing targets $E$,
uniformly as $B\to\infty$,
\begin{equation}
 \sum_SK_q(M)P_\sigma(M)
 =\frac2d Z_q(S)+
 O_C\!\left(N^{-1/17.232}(\log N)^{1/2-1/17.232}\right),
 \qquad \sigma=H,F.
 \label{eq:Mazur-full-profile-rate}
\end{equation}
\end{proposition}
\begin{proof}
Grouping endpoints by residue gives the exact identity
\[
 Z_q(S)=\sum_{x\in\mathbb Z/3^q\mathbb Z}p_q(x)c_q(S,x).
\]
The summation weights are nonnegative with sum one. Thus the
residue-uniform coefficient bounds pass to $Z_q$ without a factor
$3^q$ and without declaring $Z_q$ to have mass one.
The reference level obeys $q=O(L^{3/5})$, hence $h+q\le N$
eventually. The preceding propositions apply.

The closed interior and strict exterior partition $S$ exactly.
On the former, the triangle inequality gives $e_\sigma Z_q(S_{\rm in})$.
On the latter, the proved positive interval bound gives the second
term in \eqref{eq:Mazur-full-finite}. This proves that estimate
before any upper bound on the two masses is substituted.

The actual positive time carrier is
$V=\{\nu>0:\nu+h\in J\}$, of cardinality at most
$6W+\beta/d+1=O_C(W)$.
Uniformly over the two branches and bands, its minimum is at least
$N/43$ for large $B$: use
$\nu\ge t_0-h-3W$ and the same strict lower constant calculation
in Proposition~\ref{prop:Mazur-physical-centers}, with $3W/L=o_C(1)$.
This assertion is independent of $M$.
Hence $G_V\le\#V/\sqrt{\min V}=O_C(\sqrt{\log N})$.
The retained denominator estimates give $H\ge1/3$, $z/D\le2$,
and $e^\beta<4$, so $L_H,L_F=O_C(\sqrt{\log N})$.
For $C\ge1/2$, $A_C\ge16/(\log2)^3>48$; the last strict bound
follows already from $\log2<0.6932$.
The entire exterior charge is therefore
$O_C(N^{-A_C}\sqrt{\log N})$.

On the interior, Proposition~\ref{prop:Mazur-physical-centers} and
the exact signed denominator identities give
\[
 e_\sigma=O_C\!\left(N^{-1/17.232}
              (\log N)^{1/2-1/17.232}\right).
\]
The denominator contribution $O(1/z)$ is smaller uniformly.
Since $Z_q(S_{\rm in})\le2+\epsilon_B$, combining these estimates
proves \eqref{eq:Mazur-full-profile-rate}.
\end{proof}

The common center here is the same landing measure
$(2/d)\lambda_*\mu_q$ already constructed, not a new probability law.
This completes the reference-profile input using Rhin's cited
Diophantine estimate. The source's terminal-coordinate cap is unnecessary
for the reconstructed exterior estimate, and its coarse total coefficient
$5$ can be retained more accurately as $2+\epsilon_B$.
Section~\ref{sec:Mazur-genuine-assembly} now combines this estimate
with actual first passage and the whole-source mixtures. The original
clock and threshold-iteration consumers remain to be reconstructed.
Neither a full-package formal replay nor a new final Mazur
natural-density exponent is asserted here.

\subsection{From reference profiles to genuine passage laws}
\label{sec:Mazur-genuine-assembly}

The reference profile is useful because it approximates an actual question:
where does a Collatz orbit first enter the integers at most $B$? This section
completes that passage from the profile to the orbit. Both ways of sampling
the starting integer, harmonic and uniform, are kept. The finite boundary
charges and the possible upper endpoint of the source interval are kept too.

The source comparison is Mazur~\cite{Mazur2026v2},
\texttt{ND/Band/A5ExactAffineCorrection.lean},
\texttt{A5TerminalBoundaryShell.lean},
\texttt{A5FixedCellTerminalDescent.lean},
\texttt{A5SourceMixture.lean}, and
\texttt{ND/Discrepancy/A5ReferenceGenuineBandRate.lean},
\texttt{A5ReferenceGenuineSourceRate.lean}, at the pinned downloaded commit.
The last two files use $C=40$, a factor-seven band estimate and a factor-eight
whole-source estimate, with exponent restrictions $c<1/20$ and
$c<1/(2\kappa)$. We reconstruct their four comparisons using the sharper
estimates proved above. The analytic inputs remain Tao's cited mixing
theorem and Rhin's cited Diophantine estimate; there is no new Fourier
decay theorem here.

Throughout this section put
\[
 L=\log B,\quad N=\lfloor L/(10\log2)\rfloor,\quad
 h=\lfloor L/100000\rfloor,\quad m=\lceil hL^{-2/5}\rceil,
\]
\[
 W=C\sqrt{N\log N},\quad d=\log(4/3).
\]
Fix $C\ge1/2$. Write $\alpha=1001/1000$, and for either
$y=B^\alpha$ or $y=B^{\alpha^2}$ set
\[
 b_y=(\alpha-1)\log y,\quad k_y=\lfloor b_y\rfloor,\quad
 \beta=b_y/k_y,\quad z_j=ye^{j\beta},\quad U_j=z_{j+1},
 \qquad 0\le j<k_y.
\]
All statements below concern sufficiently large integer $B$; the threshold
may depend on $C$ but not on the branch, band index or target set.
Write $D_j,H_j$ for the exact odd-band count and reciprocal sum.
The original target set is $E'_B(E)$ from
Section~\ref{sec:Mazur-concentration-coverage}. In particular its endpoints
have actual first-passage time $h$. Its lower bound
$(4/3)^hBe^{-L^{7/10}}$ exceeds $B$ eventually.

\begin{proposition}[explicit boundary charges and their constants]
\label{prop:Mazur-genuine-boundaries}
For a retained band suppress $j$ and put
\[
 \varepsilon=B^{-9/10},\qquad \rho=B^{-19749/20000}.
\]
Every retained compatible terminal record satisfies
$0<F_a/M\le\rho$ and $0<-\log(1-F_a/M)\le\varepsilon$.
Define the following positive \emph{majorants}, using the original
$z,U,D,H$:
\begin{align*}
 b_{H,-}&=\frac{\varepsilon/2+e^\varepsilon/z}{H},&
 b_H&=\frac{\varepsilon+e^\varepsilon(1/z+1/U)}H,\\
 b_{F,-}&=\frac{z(1-e^{-\varepsilon})/2+1}{D},&
 b_F&=\frac{(z+U)(1-e^{-\varepsilon})/2+2}{D}.
\end{align*}
For the terminal union masses $U_\sigma$ and nominal sums $I_\sigma$,
\begin{equation}
 |U_H-I_H|\le\rho U_H+b_H,\qquad
 |U_F-I_F|\le b_F,\qquad
 0\le I_\sigma\le1+b_{\sigma,-}.
 \label{eq:Mazur-explicit-boundaries}
\end{equation}
Uniformly in the branches and bands,
\[
 \frac{b_{H,-}}{\varepsilon}\longrightarrow1,\quad
 \frac{b_H}{\varepsilon}\longrightarrow2,\quad
 \frac{b_{F,-}}{\varepsilon}\longrightarrow\frac1{e-1},\quad
 \frac{b_F}{\varepsilon}\longrightarrow\frac{e+1}{e-1}.
\]
These are limits of the displayed bounds, not of the actual errors.
\end{proposition}
\begin{proof}
The original padded time interval gives
\[
 \frac{L}{1000d}-3W-h\le \nu=t-h
 \le\frac{(\alpha^3-1)L}{d}+3W .
\]
The directed bounds $0.2876<d<0.2877$ established above imply
\[
 \frac{340}{100000}<\frac1{1000d}-\frac1{100000},
 \qquad \frac{\alpha^3-1}{d}<\frac{105}{10000}.
\]
Since $W=o(L)$, eventually
$F=(340/100000)L\le\nu\le(105/10000)L$,
uniformly over all retained times. The elementary bound $F_a\le3^\nu$
and $\log3<11/10$ give $F_a\le B^{231/20000}$.
The latter logarithmic inequality follows, for example, from the
exponential series through degree five at $11/10$, whose sum exceeds $3$.
Since $M>B\ge B^{999/1000}$, division gives $F_a/M\le\rho$.
Here $\rho/\varepsilon=B^{-1749/20000}\to0$. Eventually $\rho\le1/2$
and $2\rho\le\varepsilon$, so
\[
 0<-\log(1-F_a/M)\le\frac{F_a/M}{1-F_a/M}\le2\rho\le\varepsilon.
\]
Strict positivity uses $\nu\ge1$.

For either endpoint $x=z,U$, the odd shell $[xe^{-\varepsilon},x)$
has fewer than $x(1-e^{-\varepsilon})/2+1$ points.
Proposition~\ref{prop:Mazur-band-denominators}, applied with lower endpoint
$xe^{-\varepsilon}\ge1$, bounds its reciprocal sum by
$\varepsilon/2+e^\varepsilon/x$; the empty case satisfies the same
inequality directly. Apply Proposition~\ref{prop:Mazur-boundaries}.
Using the sum of both shell bounds overcovers their union; using just the
lower shell gives the one-sided cap. Injectivity of the compatible
record-to-source map is essential here and was proved in
Proposition~\ref{prop:Mazur-atoms}. It holds on the original time interval
because $\beta/d+6W\le h$ eventually.

Finally $\beta\to1$ uniformly, $z\ge B^\alpha$, and the exact discrepancies
give $H=\beta/2+O(1/z)$ and
$D=z(e^\beta-1)/2+O(1)$. Thus $1/(z\varepsilon)\to0$ uniformly.
Also $(1-e^{-\varepsilon})/\varepsilon\to1$ and
$e^\varepsilon\to1$. Substitution proves all four limits.
\end{proof}

\begin{proposition}[assembled cell tests with the endpoint mass retained]
\label{prop:Mazur-genuine-caps}
Put $\epsilon_B=3B^{-4997/600000}$ as in
Proposition~\ref{prop:Mazur-native-coefficient}. At any retained
$(\nu,\ell)$, let
\[
 S_{\nu,\ell}=\{M\in E'_B(E):z<2^\ell M/3^\nu<U\}.
\]
On the fine residue group $G_\nu=\mathbb Z/3^\nu\mathbb Z$ define
\begin{align*}
 c_H(x)&=\sum_{\substack{M\in S_{\nu,\ell}\\M\equiv x\ (3^\nu)}}
                       \frac{3^\nu}{MH},\\
 c_F(x)&=\sum_{\substack{M\in S_{\nu,\ell}\\M\equiv x\ (3^\nu)}}
                       \frac{3^\nu}{M}\frac{2^\ell M/3^\nu}{D}.
\end{align*}
Then
\begin{equation}
 0\le c_\sigma\le A_\sigma,\qquad
 A_H=\frac{2+\epsilon_B}{H},\qquad
 A_F=(2+\epsilon_B)\frac UD .
 \label{eq:Mazur-genuine-caps}
\end{equation}
The exact projection $\pi:G_\nu\to G_m$, $m\le\nu$, and its fibre average
$Ac(x)=3^{m-\nu}\sum_{\pi v=x}c(v)$ give respectively
\[
 Ac_H(x)=\frac{3^m}H
       \sum_{\substack{M\in S_{\nu,\ell}\\M\equiv x\ (3^m)}}\frac1M,\qquad
 Ac_F(x)=\frac{3^m}D
       \sum_{\substack{M\in S_{\nu,\ell}\\M\equiv x\ (3^m)}}
                       \frac{2^\ell M/3^\nu}{M}.
\]
All endpoint collisions are retained. Uniformly,
$A_H\to4$ and $A_F\to4e/(e-1)$.

Let $K_C$ be the single fixed-total mixing constant furnished by
Corollary~\ref{cor:fixed-total} with exponent $11$ on this schedule.
Keep $E_B,\gamma_B$ from \eqref{eq:Mazur-exponential-terminal}, and
the original finite time set $J$. Then the nominal sums and full physical
reference profiles $\mathcal P_\sigma$ obey
\begin{equation}
 |I_\sigma-\mathcal P_\sigma|
 \le\widehat{\mathcal R}_\sigma:=
 \frac{\gamma_B(1+b_{\sigma,-})+
   \#J\,A_\sigma\{K_C/(2m^{11})+2/m^3\}}{1-\gamma_B}.
 \label{eq:Mazur-genuine-terminal}
\end{equation}
In particular $\widehat{\mathcal R}_\sigma=O_C(L^{-1/5}\log L)$,
and its ratio to $L^{-1/5}\log L$ tends to the positive terminal
constant $A_C$ of \eqref{eq:Mazur-eta-limit}.
\end{proposition}
\begin{proof}
The retained times have $h+\nu\le N$ eventually, by the bounds in the
preceding proof. The room inequality $2\,3^{h-1}\le B$ also holds.
Every subset $S_{\nu,\ell}$ is therefore covered by
Proposition~\ref{prop:Mazur-native-coefficient} at depth $\nu$.
Division by $H$ proves the harmonic cap. On strict support
$2^\ell M/3^\nu<U$; keeping this factor inside the original sum and then
bounding it by $U$ proves the flat cap. Summing first over residue fibres
proves the two projection identities: every labelled $M$ appears once,
and $3^{m-\nu}3^\nu=3^m$. This is an averaging map on tests, not
replacement of the endpoint measure by a probability measure.
The denominator limits in the preceding proof prove the limits of $A_\sigma$.

To obtain the terminal estimate, condition a positive geometric word of
length $\nu$ on its total $\ell$. Its total probability is
$q_{\nu,\ell}=\binom{\ell-1}{\nu-1}2^{-\ell}$. The fine conditional
offset law paired with $c_\sigma$ gives exactly the nominal cell mass
divided by $q_{\nu,\ell}$, by the compatibility calculation in
Section~\ref{sec:Mazur-finite}. Its projected unconditioned pairing gives
the corresponding coarse reference cell. Apply
Corollary~\ref{cor:Mazur-centred-descent} with test cap $A_\sigma$,
$\Omega\le K_C/m^{11}$ and the uniform likelihood coefficient
$\gamma_B$. Multiplication by $q_{\nu,\ell}$ gives
\[
 |Y_{\nu,\ell}-Z_{\nu,\ell}|
 \le q_{\nu,\ell} A_\sigma
       \{K_C/(2m^{11})+2/m^3\}+\gamma_B Z_{\nu,\ell}.
\]
The hypotheses are supplied by the original scales:
$F\le\nu\le N$, $m\asymp L^{3/5}\ge\sqrt{\nu}$ eventually,
$|\ell-2\nu|<W=O_C(\sqrt{\nu\log\nu})$,
$8m=o(F)$, $W=o(F)$ and $mW/(F\sqrt{m\log m})\to0$.
Thus the mixing constant is chosen before $B$, uniformly in all cells;
the two tail estimates used in descent were proved in
Corollary~\ref{cor:Mazur-two-tails}.

Now sum over totals using $\sum_\ell q_{\nu,\ell}\le1$, then over times.
Lemma~\ref{lem:Mazur-sharp-absorption} and
$I_\sigma\le1+b_{\sigma,-}$ give
\eqref{eq:Mazur-genuine-terminal}. The exact two-row reindexing and
time translation $t=\nu+h$ identify the reference sum with
$\mathcal P_\sigma$, without changing either band weight.
Since $\#J=O_C(W)$ and $A_\sigma=O(1)$,
the nonproportional numerator is
$O_C(L^{-13/10}\sqrt{\log L})=o(L^{-1/5}\log L)$.
The limit of $\gamma_B/(L^{-1/5}\log L)$ was proved in
Corollary~\ref{cor:Mazur-exponential-terminal}; both remaining factors
tend to one. This proves the limit. Eventually
$A_H\le20$, $A_F\le80$, $b_{\sigma,-}\le11\varepsilon$ and $\#J\le7W$,
so this bound also refines $\mathcal R_\sigma^\dagger$.
\end{proof}

\begin{theorem}[the unchanged genuine passage measures]
\label{thm:Mazur-genuine-rate}
For each band let $p_{\sigma,j}$ be its actual harmonic or uniform
probability law, and define the finite-passage subprobability
\[
 G_{\sigma,j}(E)=
 p_{\sigma,j}\{n:\tau_B(n)<\infty,\ T^{\tau_B(n)}n\in E\}.
\]
No assertion that every orbit reaches $B$ is implicit in this definition.
Keep the original, unnormalized common measure
\[
 \mathcal Z_B(E)=\frac2d\sum_{M\in E'_B(E)}
                         \frac{3^m p_m(M\bmod3^m)}M.
\]
Let $\chi_{\sigma,j}$ denote the corresponding exact full-prefix
majorant in \eqref{eq:Mazur-harmonic-fullgood} or
\eqref{eq:Mazur-flat-fullgood}, and let $Q_{\sigma,j}(E)$ denote the
right side of \eqref{eq:Mazur-full-finite}. The explicit finite conclusion is
\begin{align}
 |G_{H,j}(E)-\mathcal Z_B(E)|
 &\le\chi_{H,j}+\rho U_H+b_H+
                  \widehat{\mathcal R}_H+Q_{H,j}(E),\nonumber\\
 |G_{F,j}(E)-\mathcal Z_B(E)|
 &\le\chi_{F,j}+b_F+
                  \widehat{\mathcal R}_F+Q_{F,j}(E).
 \label{eq:Mazur-genuine-four}
\end{align}
In particular, with $b_*=1/17.232$,
\begin{equation}
 \sup_{\text{branches},\,j,\,E,\,\sigma}
 |G_{\sigma,j}(E)-\mathcal Z_B(E)|
 =O_C\!\left(L^{-b_*}(\log L)^{1/2-b_*}\right).
 \label{eq:Mazur-genuine-rate}
\end{equation}
This holds for every fixed $C\ge1/2$, including $C=1/2$.
\end{theorem}
\begin{proof}
Insert, in order, the actual terminal union mass $U_\sigma$, nominal
sum $I_\sigma$, and physical reference profile $\mathcal P_\sigma$.
The resulting four differences are bounded by
Corollary~\ref{cor:Mazur-coverage-mass} with
Corollary~\ref{cor:Mazur-band-fullgood}, the preceding two propositions,
and Proposition~\ref{prop:Mazur-full-profile}. This proves the finite
inequalities and explains each term; it is not a substitution of the
reference event for actual first passage.

The estimates already proved are uniform in $E$, including targets
depending on $B$. Coverage contributes
$O_C(N^{-C^2/4})$; the boundary and affine terms are $O(B^{-9/10})$;
the terminal term is $O_C(L^{-1/5}\log L)$.
The full reference error is
$O_C(N^{-b_*}(\log N)^{1/2-b_*})$.
Since $C^2/4\ge1/16>b_*$ and $N\sim L/(10\log2)$, the first three
terms are of smaller order than the last. For example their quotient
for the terminal term is
$O_C(L^{-(1/5-b_*)}(\log L)^{1/2+b_*})\to0$.
This proves \eqref{eq:Mazur-genuine-rate} and the stated order of
quantifiers: one threshold and constant for all branches, bands and
landing targets, once $C$ is fixed.
\end{proof}

\begin{corollary}[whole-source laws and exact completion at nonpassage]
\label{cor:Mazur-genuine-source}
On each original closed source block
\[
 \mathcal I_y=[y,y^\alpha]\cap(2\mathbb N+1),
 \qquad y=B^\alpha\ \hbox{or}\ B^{\alpha^2},
\]
let $G_{\sigma,y}$ be its genuine finite-passage subprobability under
the harmonic or uniform starting law. Then
\[
 \sup_{\sigma,y,E}|G_{\sigma,y}(E)-\mathcal Z_B(E)|
 =O\!\left(L^{-b_*}(\log L)^{1/2-b_*}\right).
\]
For every $0<c<b_*$, the source's factor-seven and factor-eight
conclusions therefore hold eventually with $L^{-c}$ in place of
its rate, without the additional restriction $c<1/20$.
The implicit constants above may be fixed by taking $C=1/2$.

Define the actual total landing map
\[
 \Lambda_B:2\mathbb N+1\longrightarrow
 \mathcal Y_B:=\{r>0:r\text{ odd},\ r\le B\}\sqcup\{\partial\},
 \qquad
 \Lambda_B(n)=
 \begin{cases}T^{\tau_B(n)}n,&\tau_B(n)<\infty,\\
 \partial,&\tau_B(n)=\infty .
 \end{cases}
\]
Here $\partial$ is one extra point, not a numerical landing value.
For any two of the four whole-source starting laws, their pushforwards
$\widehat G_1,\widehat G_2$ satisfy
\begin{equation}
 \|\widehat G_1-\widehat G_2\|_{\rm TV}
 :=\sup_{A\subseteq\mathcal Y_B}|\widehat G_1(A)-\widehat G_2(A)|
 =\sup_{E\subseteq\mathcal Y_B\setminus\{\partial\}}
                      |G_1(E)-G_2(E)|
 =O\!\left(L^{-b_*}(\log L)^{1/2-b_*}\right).
 \label{eq:Mazur-completed-tv}
\end{equation}
In particular it is at most $16L^{-c}$ eventually for every $0<c<b_*$.
For every real function $f$ on $\mathcal Y_B$, the difference of its
expectations is bounded by
$(\max f-\min f)\|\widehat G_1-\widehat G_2\|_{\rm TV}$.

The original common mass itself satisfies
\[
 \mathcal Z_B(\mathcal Y_B\setminus\{\partial\})
 =1+O\!\left(L^{-b_*}(\log L)^{1/2-b_*}\right),
\]
or, equivalently, the unscaled original kernel mass on $E'_B([1,B])$
is $d/2+O(L^{-b_*}(\log L)^{1/2-b_*})$. No division by that mass is used.
\end{corollary}
\begin{proof}
The half-open bands partition $[y,y^\alpha)$, with exactly one additional
point if $y^\alpha$ is an odd integer. Write $H_y,D_y$ for the reciprocal
sum and count of the closed block. The respective band weights are
$H_j/H_y$ and $D_j/D_y$; the possible top weights are
$p_H=(1/y^\alpha)/H_y$ and $p_F=1/D_y$ when that point exists, and zero
otherwise. In either case $\sum_jw_{\sigma,j}+p_\sigma=1$ exactly.
For any event its whole-source mass is
$\sum_jw_{\sigma,j}G_{\sigma,j}(E)+t_\sigma(E)$, where
$0\le t_\sigma(E)\le p_\sigma$.

For any nonnegative center $Z$ and valid band errors $r_j$ this gives
the finite identity and bound
\[
 G_{\sigma,y}(E)-Z
 =\sum_jw_{\sigma,j}(G_{\sigma,j}(E)-Z)+(t_\sigma(E)-p_\sigma Z),
\]
\[
 |G_{\sigma,y}(E)-Z|
 \le\sum_jw_{\sigma,j}r_j+
       p_\sigma\max\{Z,|1-Z|\}.
\]
The last factor is exact for a variable $t_\sigma/p_\sigma\in[0,1]$;
when $p_\sigma=0$ its term is zero. The first band gives
$H_y\ge1/3$ and $D_y>y/2$, hence $p_H\le3/B$ and $p_F<2/B$.
The endpoint coefficient bound gives
$0\le\mathcal Z_B(E)\le(2/d)(2+\epsilon_B)$.
Thus the top costs $O(B^{-1})$ uniformly, not the number of bands.
This proves the whole-source rate. Because
$L^{-(b_*-c)}(\log L)^{1/2-b_*}\to0$, each band and source error
is $o(L^{-c})$; the particular source constants $7$ and $8$ follow.

For the completion, restriction to the finite landing points gives $G_i$
exactly, and the mass at $\partial$ is $1-G_i([1,B])$. If a set $A$
contains $\partial$, its complementary set does not, and equality of the
total probability masses gives
$|\widehat G_1(A)-\widehat G_2(A)|
 =|G_1(\mathcal Y_B\setminus A)-G_2(\mathcal Y_B\setminus A)|$.
This proves the equality of the two suprema. The triangle inequality
through the \emph{unchanged} common measure proves the rate and constant
$16$; there is no separate factor for the nonpassage atom.
For the expectation bound, subtract $\min f$ and write
$f-\min f=\int_0^{\max f-\min f}\mathbf1_{\{f-\min f>s\}}\,ds$.
The domain is finite, so summation and integration interchange directly.

Finally the good-prefix event ensures passage by time $N$, as proved in
Proposition~\ref{prop:Mazur-genuine-coverage}. Band failure mass is at most
$\chi_{\sigma,j}$; whole-source failure mass is at most its mixture plus
$p_\sigma$. At $C=1/2$ this is $O(N^{-1/16})+O(B^{-1})$.
Therefore $G_{\sigma,y}([1,B])=1+O(N^{-1/16})$.
Its comparison with $\mathcal Z_B([1,B])$ proves the asserted mass limit.
\end{proof}

These results strengthen the genuine passage-law stage of Mazur's
argument and connect it to the previously reconstructed Allikvere
passage comparison. Section~\ref{sec:Mazur-real-timed} carries them
through real thresholds and the original scheduled iteration to a
timed natural-density theorem. No Lean replay is inferred from the
written proofs.

\begin{figure}[ht]
\centering
\fbox{\begin{minipage}{0.95\linewidth}
\[
 \begin{array}{ccccc}
 G_{\sigma,j}&\longleftrightarrow&U_\sigma&
                    \longleftrightarrow&I_\sigma\\
 &\chi_{\sigma,j}&&b_\sigma+\mathbf1_{\sigma=H}\rho U_H&\\[5pt]
 I_\sigma&\longleftrightarrow&\mathcal P_\sigma&
                    \longleftrightarrow&\mathcal Z_B\\
 &\widehat{\mathcal R}_\sigma&&Q_{\sigma,j}(E)&
 \end{array}
\]
The first node is actual finite first passage; the last is the original
common measure. Each arrow is an error estimate, not equality of maps.
The inequalities are \eqref{eq:Mazur-genuine-four}. Convex band weights
then give the whole-source law, with its separate upper-endpoint charge.
The total map $\Lambda_B$ adds the actual nonpassage point without
altering any finite landing mass; \eqref{eq:Mazur-completed-tv} gives
its exact comparison.
\end{minipage}}
\caption{How the four estimates reach actual Collatz passage laws.
This retains Mazur's original comparisons with the quantitative
refinements proved here from the cited Tao and Rhin inputs.
The endpoints of the diagram and all intermediate quantities are defined
in this section; the arrows do not denote deterministic orbit maps.}
\label{fig:Mazur-genuine-chain}
\end{figure}

\subsection{Real thresholds and the timed density conclusion}
\label{sec:Mazur-real-timed}

The preceding rate concerns genuine landing probabilities. To use it in
the scale iteration, two changes must be made explicitly: the source
formalization sends nonpassage to the number \(1\), and its iteration
uses real thresholds and real source windows. Neither change requires
a loss in the logarithmic exponent. Keeping the exact geometric sum of
the original scheduled times also improves the source's clock constants.

The source route is Mazur~\cite{Mazur2026v2},
\texttt{ND/Probability/PassTotalizer.lean},
\texttt{ND/Discrepancy/A5ReferenceTotalizedSourceRate.lean},
\texttt{A5ReferenceRealFloorPass.lean}, and
\texttt{A5ReferenceRealLogBlock.lean}. Its scale definitions are in
\texttt{Tao/Section3/AmbientPartialTrace.lean}; its clock comparison is
in \texttt{Tao/Section3/TimeBudget.lean} and
\texttt{ND/LogTime/Clock.lean}. All refer here to the pinned source
commit specified in the bibliography. The underlying passage method
is Tao's~\cite{Tao2026v7}; the quantitative input remains the written
Tao--Rhin comparison above. The following is a written reconstruction
and refinement, not a claim to have replayed the whole Lean package.

Put
\[
 \alpha=\frac{1001}{1000},\qquad b=\frac1{17.232},\qquad
 a=\frac12-b,\qquad
 R(L)=L^{-b}(1+\log L)^a\quad(L\ge\log2).
\]
This positive rate is asymptotic to \(L^{-b}(\log L)^a\).
It only extends the displayed scalar rate to small positive logarithms;
it does not alter a measure, source window, or passage map.
For large \(L\), \(R\) is decreasing, since its logarithmic derivative
with respect to \(L\) is \(-b+a/(1+\log L)<0\).

\begin{proposition}[the exact effect of sending nonpassage to one]
\label{prop:Mazur-totalizer-exact}
For integer \(B\ge1\), let \(\mathcal Y_B\) and \(\Lambda_B\) be as in
Corollary~\ref{cor:Mazur-genuine-source}. The map
\[
 r_B:\mathcal Y_B\longrightarrow\{0,1,\ldots,B\},\qquad
 r_B(s)=s\ (s\ne\partial),\quad r_B(\partial)=1
\]
has image precisely the positive odd points of its codomain. Every
nonempty fibre is a singleton except \(r_B^{-1}(1)=\{1,\partial\}\).
The source's totalized map is exactly \(F_B=r_B\circ\Lambda_B\).

For probability measures \(\mu,\nu\) on \(\mathcal Y_B\), write
\(\Delta=\mu-\nu\), \(u=\Delta(\{1\})\), \(v=\Delta(\{\partial\})\).
With \(\|\Delta\|_1=\sum_s|\Delta(\{s\})|\),
\begin{equation}
 \|\Delta\|_1-\|(r_B)_*\Delta\|_1
 =|u|+|v|-|u+v|
 =\begin{cases}2\min(|u|,|v|),&uv<0,\\0,&uv\ge0.\end{cases}
 \label{eq:Mazur-totalizer-loss}
\end{equation}
Moreover \(\|\Delta\|_1=2\sup_A|\Delta(A)|\).
Consequently each natural-source pair in
Corollary~\ref{cor:Mazur-genuine-source} has totalized full \(\ell^1\)
distance \(O(R(\log B))\), and at most \(32(\log B)^{-c}\)
eventually for every \(0<c<b\).
\end{proposition}
\begin{proof}
First passage, if it exists, has one uniquely determined time and
landing point. This proves the map formula and fibres, including the
zero masses at even points and at zero. Pushforward adds only the two
coordinates \(u,v\); all other coordinates are unchanged, proving
the first equality in \eqref{eq:Mazur-totalizer-loss}. Splitting
according to the signs of \(u,v\) proves the second.

Since \(\sum_s\Delta(\{s\})=0\), the positive and negative parts each
have mass \(\|\Delta\|_1/2\). The event consisting of the positive
coordinates attains that value, and no event exceeds it in absolute
value. This proves the convention comparison and the rate.
The linear kernel of \((r_B)_*\) consists exactly of
\(t(\delta_1-\delta_\partial)\), \(t\in\mathbb R\).
Thus no inverse on arbitrary laws exists: \(\delta_1\) and
\(\delta_\partial\) have the same image. If the nonpassage mass is
provided separately, reconstruction is exact: retain every other
coordinate, set the \(\partial\)-coordinate to that mass, and subtract
it from the image mass at \(1\). An admissible mass lies between zero
and that image mass. This describes precisely what was lost.
\end{proof}

\begin{proposition}[exact overlap for unchanged finite source laws]
\label{prop:Mazur-source-overlap}
Let \(S,T\) be nonempty finite subsets of an ambient set and let \(w>0\)
on \(S\cup T\). Write \(W_A=\sum_{n\in A}w(n)\) and
\(p_A(n)=w(n)\mathbf1_A(n)/W_A\).
Then
\begin{equation}
 \sup_E|p_S(E)-p_T(E)|
 =1-\frac{W_{S\cap T}}{\max(W_S,W_T)}
 =\frac{\max(W_{S\setminus T},W_{T\setminus S})}
             {\max(W_S,W_T)}.
 \label{eq:Mazur-overlap}
\end{equation}
Their full \(\ell^1\) distance is twice this number.
For every common map \(f:S\cup T\to Z\), the output event distance
is at most \eqref{eq:Mazur-overlap}; its exact value is
\[
 \frac12\sum_{z\in f(S\cup T)}
 \left|\sum_{n:f(n)=z}(p_S(n)-p_T(n))\right|.
\]
\end{proposition}
\begin{proof}
For any two probability vectors \(p,q\),
\(\frac12\sum|p-q|=1-\sum\min(p,q)\), by
\(|p-q|=p+q-2\min(p,q)\). On the intersection here,
\(\min(p_S(n),p_T(n))=w(n)/\max(W_S,W_T)\);
elsewhere it is zero. This proves the first equality. Writing
\(W_S=W_{S\cap T}+W_{S\setminus T}\) and similarly for \(T\)
proves the second. Summing differences on each fibre gives the
displayed output formula. The triangle inequality on each fibre
proves contraction. This applies to the existing probabilities
without replacing either denominator. For an injective \(f\) equality
holds; cancellation of opposite signs within fibres describes every
strict loss.
\end{proof}

\begin{proposition}[real/floor transport with both endpoints retained]
\label{prop:Mazur-real-overlap}
For \(x\ge1\), set \(B=\lfloor x\rfloor\). For either
\(\beta=\alpha\) or \(\beta=\alpha^2\), put \(q=\alpha\beta\) and
\[
 S=[B^\beta,B^q]\cap(2\mathbb N+1),\qquad
 T_x=[x^\beta,x^q]\cap(2\mathbb N+1).
\]
Assume \(S,T_x\ne\varnothing\). Define \(\epsilon_F(x,\beta)\) and
\(\epsilon_H(x,\beta)\) by the exact right side of
\eqref{eq:Mazur-overlap}, using respectively \(w(n)=1\) and \(w(n)=1/n\).
For all sufficiently large \(B\), uniformly in \(B\le x<B+1\) and
the two branches,
\begin{align}
 \epsilon_F(x,\beta)
 &\le
 \frac{(x^q-B^q)/2+1}{(B^q-B^\beta)/2-1},
 \label{eq:Mazur-real-flat}\\
 \epsilon_H(x,\beta)
 &\le
 \frac{(q/2)\log(x/B)+B^{-\beta}}
      {((q-\beta)/2)\log B-B^{-\beta}}.
 \label{eq:Mazur-real-harmonic}
\end{align}
The displayed denominators are positive in this range.
In particular
\[
 \limsup_{B\to\infty}\ \sup_{B\le x<B+1} B\epsilon_F(x,\beta)\le q,
 \qquad
 \limsup_{B\to\infty}\ \sup_{B\le x<B+1}
 B\log B\,\epsilon_H(x,\beta)\le
 \frac{q}{q-\beta}=1001.
\]
The exact errors vanish at \(x=B\), even though these upper envelopes
need not vanish there.
\end{proposition}
\begin{proof}
On an interval of length \(v-u\), the number of odd integers differs
from \((v-u)/2\) by at most \(1\), for any choice of endpoint inclusion.
For completeness, if \(A(t)\) counts those odd points between \(u\)
and \(t\), then \(|A(t)-(t-u)/2|\le1\). Summation by parts gives
\[
 \sum_{n\in[u,v]\ {\rm odd}}\frac1n
 =\frac{A(v)}v+\int_u^v\frac{A(t)}{t^2}\,dt.
\]
The same formula uses the appropriate \(A\) for a deleted endpoint.
Replacing \(A(t)\) by \((t-u)/2\) gives
\(\tfrac12\log(v/u)\); the total absolute error is at most
\(1/v+\int_u^v t^{-2}dt=1/u\).
This also handles an atom at \(u\).

Eventually \(x^\beta\le B^q\), uniformly in \(x<B+1\), because
\((B+1)^\beta/B^q\to0\).
The two exclusive sets are then exactly
\([B^\beta,x^\beta)\cap(2\mathbb N+1)\) and
\((B^q,x^q]\cap(2\mathbb N+1)\).
The larger of their cardinalities is at most
\((x^q-B^q)/2+1\): indeed \(t\mapsto t^q-t^\beta\) is increasing
for \(t\ge1\), since \(q>\beta>1\).
The full mass of \(S\) is at least \((B^q-B^\beta)/2-1\).
Equation~\eqref{eq:Mazur-overlap} proves \eqref{eq:Mazur-real-flat}.

The exclusive reciprocal masses are at most
\(\frac{\beta}{2}\log(x/B)+B^{-\beta}\) and
\(\frac q2\log(x/B)+B^{-q}\).
Their maximum is bounded by the latter logarithmic term plus
\(B^{-\beta}\). The reciprocal mass of \(S\) is at least the
denominator of \eqref{eq:Mazur-real-harmonic}, proving that bound.

Finally \(x^q-B^q\le q(B+1)^{q-1}(x-B)\) and
\(\log(x/B)\le(x-B)/B\).
Divide numerator and denominator by the indicated powers.
The additive cardinal error contributes \(O(B^{1-q})\) after
multiplication by \(B\), and the harmonic error contributes
\(O(B^{1-\beta})\) after multiplication by \(B\log B\).
Both vanish. These estimates are uniform in \(x-B\in[0,1)\).
\end{proof}

\begin{theorem}[real-source passage and its actual scheduled failure]
\label{thm:Mazur-real-passage}
Whenever the unchanged closed odd block \([y,y^\alpha]\) is nonempty,
let \(H_y,F_y\) be its harmonic and uniform probability laws.
For real \(x\) let
\[
 n_x=\left\lfloor\frac{\log\lfloor x\rfloor}{10\log2}\right\rfloor,
 \qquad
 F_x(n)=
 \begin{cases}T^{\tau_x(n)}n,&\tau_x(n)<\infty,\\1,&\tau_x(n)=\infty.
 \end{cases}
\]
For either source branch \(y=x^\alpha,x^{\alpha^2}\) and either law,
the probability of \(\tau_x>n_x\), including nonpassage, is
\(O((\log x)^{-1/16})\). Every pair of the four totalized landing
laws has event distance \(O(R(\log x))\) and full \(\ell^1\)
distance \(O(R(\log x))\), uniformly as \(x\to\infty\).
No hypothesis about the existence of passage on every orbit is made.
\end{theorem}
\begin{proof}
Since every iterate is an integer,
\[
 T^j n\le x\ \Longleftrightarrow\ T^j n\le\lfloor x\rfloor
 \quad\hbox{for every }j.
\]
Thus first-hit times, nonpassage events, and landing values agree
exactly at \(x\) and \(B=\lfloor x\rfloor\).
The identity map on \(\{0,\ldots,B\}\) is the carrier equivalence;
there is no rounding error in the dynamics.

Let \(\mu_{\sigma,\beta}^{x}\) denote a real source law and
\(\mu_{\sigma,\beta}^{B}\) its floor source law.
Propositions~\ref{prop:Mazur-source-overlap} and
\ref{prop:Mazur-real-overlap} give event distance
\(\epsilon_\sigma(x,\beta)\) before the common map \(F_B\), hence
at most that distance after it. If \(D_B(i,j)\) is the exact
event distance between the natural-source outputs, the finite
pairwise bound is
\begin{equation}
 D_x(i,j)\le\epsilon_i+D_B(i,j)+\epsilon_j.
 \label{eq:Mazur-real-pair}
\end{equation}
Full \(\ell^1\) is twice this entire expression, not twice just its
middle term. In particular the natural bound \(32(\log B)^{-c}\)
from the totalizer proposition remains the middle full-\(\ell^1\)
term. The source's safe outer estimates were \(96x^{-99/100}\)
for uniform events and \(96000/(B\log B)\) for harmonic events.
Equations~\eqref{eq:Mazur-real-flat}--\eqref{eq:Mazur-real-harmonic}
replace those outer terms using the original finite sets.

The good-prefix argument in
Proposition~\ref{prop:Mazur-genuine-coverage} proves passage by
\(n_B=\lfloor\log B/(10\log2)\rfloor\), not just eventual passage.
Corollary~\ref{cor:Mazur-band-fullgood}, with \(C=1/2\), bounds
its complement by \(O((\log B)^{-1/16})\) on every band.
The exact closed-block mixture in
Corollary~\ref{cor:Mazur-genuine-source} adds at most the top
endpoint mass \(O(B^{-1})\). Applying the source overlap bound
to the single event \(\{\tau_B>n_B\}\) costs only
\(\epsilon_\sigma\), with no full-\(\ell^1\) factor.
This proves the scheduled failure estimate for real windows.
Finally \(\log B/\log x\to1\);
\(\epsilon_F=O(B^{-1})\), \(\epsilon_H=O((B\log B)^{-1})\),
and \(1/16>b\). Thus both the source perturbation and scheduled
failure are of smaller order than \(R(\log x)\).
Equation~\eqref{eq:Mazur-real-pair} proves the asserted rate.
\end{proof}

\begin{proposition}[the actual time-carrying scale iteration]
\label{prop:Mazur-critical-iteration}
There are constants \(C_0,B_0\), independent of \(B\) and \(y\),
such that for every integer \(B\ge B_0\) and every \(y\ge2\)
with a nonempty odd source block, the harmonic law \(H_y\) satisfies
\begin{equation}
 H_y\left\{n:\tau_B(n)>
 \frac{\log y}{10\alpha(\alpha-1)\log2}\right\}
 \le C_0 R(\log B).
 \label{eq:Mazur-harmonic-timed}
\end{equation}
In fact its integer time budget is the exact finite sum below.

For a seed \(y_0>1\), define
\(y_j=y_0^{\alpha^j}\), \(q_j=y_j^{1/\alpha}\), and
\(K_j=\sum_{i<j}n_{q_i}\). Let \(d_j\) be the event distance
between \((F_{q_j})_*H_{y_j}\) and
\((F_{q_j})_*H_{y_{j+1}}\), and put
\(s_j=H_{y_{j+1}}\{\tau_{q_j}>n_{q_j}\}\).
Whenever \(y_0^\alpha\le B\),
\begin{equation}
 H_{y_J}\{\tau_B>K_J\}\le\sum_{j<J}(d_j+s_j),\qquad
 K_J\le\frac{\log q_0(\alpha^J-1)}
                   {10(\alpha-1)\log2}.
 \label{eq:Mazur-timed-telescope}
\end{equation}
The finite sum of the critical comparison rates is bounded, uniformly
in \(J\), by
\begin{equation}
 \sum_{j<J}R(\alpha^j L)
 \le R(L)\left(\frac1{1-r}
       +\frac{a\log\alpha}{1+\log L}\frac r{(1-r)^2}\right),
 \qquad r=\alpha^{-b},\quad L\ge1.
 \label{eq:Mazur-critical-geometric}
\end{equation}
For the infinite sum, its ratio to \(R(L)\) tends to \(1/(1-r)\)
as \(L\to\infty\).
\end{proposition}
\begin{proof}
For \(K\in\mathbb N\), let \(A_{B,K}=\{n\text{ odd}:\tau_B(n)>K\}\).
For every \(q\ge1\) and positive odd \(n\),
\begin{align}
 \mathbf1_{A_{B,K}}(F_q(n))&\le\mathbf1_{A_{B,K}}(n),\nonumber\\
 \mathbf1_{A_{B,K+N}}(n)&\le
 \mathbf1_{A_{B,K}}(F_q(n))+\mathbf1_{\{\tau_q>N\}}(n)
 \qquad(N\in\mathbb N).
 \label{eq:Mazur-timed-pointwise}
\end{align}
For the first inequality, suppose \(\tau_B(n)\le K\). At nonpassage
\(F_q(n)=1\), which already lies below \(B\).
If first passage to \(q\) occurs before the hit at \(B\), the residual
time to \(B\) is at most \(K\). If it occurs later, then the earlier
hit at \(B\) had value above \(q\); hence \(q<B\), and the landing
at \(q\) is itself below \(B\). This proves the first inequality
without an assumed ordering between \(q\) and \(B\).
For the second, when \(\tau_q\le N\) and the landing hits \(B\)
within \(K\), concatenation supplies a hit within \(N+K\).
Otherwise its right side is at least one whenever necessary.

Integrating the second inequality under \(H_{y_{j+1}}\),
comparing its landing event with the one under \(H_{y_j}\),
and applying the first inequality gives
\[
 H_{y_{j+1}}(A_{B,K_j+n_{q_j}})
 \le H_{y_j}(A_{B,K_j})+d_j+s_j.
\]
These are the literal two source branches at threshold \(q_j\):
\(q_j^\alpha=y_j\) and \(q_j^{\alpha^2}=y_{j+1}\).
All points in the seed block are at most \(B\), so its bad
probability at \(K_0=0\) is zero. Induction proves the first
part of \eqref{eq:Mazur-timed-telescope}.
Since \(n_{q_j}\le\log q_j/(10\log2)\) and
\(\log q_j=\alpha^j\log q_0\), summing the finite geometric
series proves its clock bound, with no extra or unused step.

For \eqref{eq:Mazur-critical-geometric}, write \(h=\log\alpha\).
Then exactly
\[
 \frac{R(\alpha^jL)}{R(L)}
 =r^j\left(1+\frac{jh}{1+\log L}\right)^a.
\]
Here \(0<a<1\). Concavity, or integration of
\(a(1+t)^{a-1}\le a\), bounds the last factor by
\(1+ajh/(1+\log L)\). Summing
\(\sum r^j=(1-r)^{-1}\) and
\(\sum jr^j=r/(1-r)^2\) proves the upper bound.
The lower bound for the infinite sum is \((1-r)^{-1}\),
because the parenthesized factor is at least one; squeezing
proves the limit. All terms are nonnegative, so no interchange
of conditionally convergent sums occurs.

Now fix \(y\) with \(y^\alpha>B\).
Choose the least integer \(J\ge1\) such that
\(y_0=y^{\alpha^{-J}}\) satisfies \(y_0^\alpha\le B\).
Existence follows from \(y^{\alpha^{-J}}\to1\).
Minimality gives \(y_0^{\alpha^2}>B\), whence
\[
 B^{1/\alpha^2}<y_0\le B^{1/\alpha},\qquad
 B^{1/\alpha^3}<q_0\le B^{1/\alpha^2}.
\]
Thus one sufficiently large \(B_0\), chosen before \(y,J\), puts
all \(q_j\) in the range of Theorem~\ref{thm:Mazur-real-passage}.
That theorem gives \(d_j+s_j\le C R(\log q_j)\) with one \(C\).
The geometric bound is \(O(R(\log q_0))\).
Since \(\log q_0\in(\alpha^{-3}\log B,\alpha^{-2}\log B]\),
direct substitution gives
\(R(\log q_0)\le\alpha^{3b}R(\log B)\) for large \(B\).
The denominator \(1-r\) is fixed and strictly positive.
This proves the probability estimate, uniformly in the trace length.
Finally \(\alpha^J\log q_0=(\log y)/\alpha\), proving
\eqref{eq:Mazur-harmonic-timed}.
If \(y^\alpha\le B\), all starts already hit at time zero and the
same statement holds. Both this case and an empty trace have zero cost.
\end{proof}

\begin{theorem}[timed natural density from the original Mazur schedules]
\label{thm:Mazur-critical-density}
There is an absolute \(C\) such that for every integer \(B\ge2\) and
real \(X\ge2\),
\begin{align}
 \frac1X\#\left\{1\le n\le X:n\text{ odd},\
       \tau_B(n)>\frac{100}{\log2}\log n\right\}
 &\le C R(\log B),\label{eq:Mazur-odd-density}\\
 \frac1X\#\left\{1\le n\le X:\
       \tau_B^{\rm raw}(n)>\frac{301}{\log2}\log n\right\}
 &\le 2C R(\log B).\label{eq:Mazur-raw-density}
\end{align}
Here the ordinary Collatz map is \(n/2\) for even \(n\) and
\(3n+1\) for odd \(n\); \(\tau_B^{\rm raw}\) is its first hitting
time of \([1,B]\). These are not accelerated raw steps.

In particular every function \(f:\mathbb N_{>0}\to\mathbb R\)
tending to infinity is exceeded downwards by an iterate within
\((301/\log2)\log n\) ordinary steps for a set of natural density one:
the iterate is \emph{strictly} below \(f(n)\).
The corresponding odd-relative-density statement holds for functions
tending to infinity on the odd integers with the Syracuse clock
\((100/\log2)\log n\). For every \(0<c<b\), the right sides of the
fixed-target estimates are also \(O_c((\log B)^{-c})\).

These constants refine the source's
\(501501/(5000\log2)\) and \(1509503/(5000\log2)\)
within its own scheduled construction. They are not claimed to be
the best clocks in this workbench: the earlier Allikvere-based
reconstruction has a much smaller leading Syracuse coefficient.
\end{theorem}
\begin{proof}
First take \(B\) large enough for the preceding results.
When \(X\le B\) the counted sets are empty.
For \(X>B\), put \(y=X^{1/\alpha}\) and
\(q=X^{1/\alpha^2}\). The top odd block is exactly
\([y,X]\cap(2\mathbb N+1)\), including both integer endpoints.
If \(y^\alpha\le B\) there is nothing to prove; otherwise choose
the harmonic trace ending at \(H_y\) as in the preceding proposition,
with its integer budget \(K\). Integrating
\eqref{eq:Mazur-timed-pointwise} under \(F_y\), comparing its
landing law at \(q\) with that of \(H_y\), then using the first
inequality there, gives
\[
 F_y\{\tau_B>K+n_q\}
 \le H_y\{\tau_B>K\}
    +D((F_q)_*F_y,(F_q)_*H_y)
    +F_y\{\tau_q>n_q\}.
\]
Here \(D\) denotes event-supremum distance. All three terms are
\(O(R(\log B))\), uniformly in \(X>B\): the first by the trace,
and the last two by the real passage theorem and
\(\log q=\alpha^{-2}\log X\).
For the comparison to \(R(\log B)\), use eventual monotonicity at
\(\alpha^{-2}\log B\), followed by direct substitution.

Write \(k=10\log2\). The exact trace bound and the final schedule yield
\[
 K+n_q\le
 \frac{\log y}{k\alpha(\alpha-1)}+\frac{\log q}{k}
 =\frac{\log X}{k\alpha(\alpha-1)}.
\]
There is no use of \(\log q\le\log X\) in this equality.
Every point in the top block satisfies \(\log X\le\alpha\log n\).
It follows that \(K+n_q\le\log n/[k(\alpha-1)]
=(100/\log2)\log n\).
Thus the top-block bad fraction is \(O(R(\log B))\).
There are at most \(y+1\) positive integers below its lower boundary,
so the contribution of all omitted odd points divided by \(X\) is
at most \(X^{-(1-1/\alpha)}+X^{-1}\).
Uniformly for \(X>B\) this is \(O(R(\log B))\), since every negative
power of \(B\) is \(o(R(\log B))\).
The number of odd points in the top block is at most \(X\) for
\(X\ge2\), so changing its probability to count divided by \(X\)
does not increase its bound. This proves \eqref{eq:Mazur-odd-density}
for large \(B\).
For the finitely many integers \(2\le B<B_0\), the counted ratio
is at most one and \(\min_{2\le B<B_0}R(\log B)>0\);
increasing one \(C\) handles all of them. The same odd count bound
also holds for \(0<X<2\), since the only possible positive odd input,
\(1\), is good at time zero.

For the raw clock, write uniquely \(n=2^v u\), \(u\) positive odd.
Suppose \(u=u_0,\ldots,u_t\) is a Syracuse trajectory with \(u_t\le B\)
and \(t\le(100/\log2)\log u\).
If \(a_i=v_2(3u_i+1)\), then
\[
 2^{\sum_{i<t}a_i}u_t=\prod_{i<t}(3+1/u_i)\,u
 \le 4^t u.
\]
Taking base-two logarithms and using \(u_t\ge1\) gives
\(\sum a_i\le2t+\log_2u\). The ordinary trajectory reaches
\(u_t\) in exactly \(v+t+\sum a_i\) steps. Therefore its first
hit is no later than
\[
 v+3t+\log_2u
 \le \frac{300}{\log2}\log u+\frac{\log n}{\log2}
 \le \frac{301}{\log2}\log n.
\]
This is an exact map from a Syracuse path to the ordinary path:
each odd edge expands into one \(3u+1\) edge followed by its
\(a_i\) halving edges, after the \(v\) initial halvings.
There are no deleted labels or valuations.

Consequently every raw-bad \(n\) has odd part \(u\) in the odd-bad
set. The unique factorization gives, with no residual cutoff,
\[
 \#\{\text{raw-bad }n\le X\}
 \le\sum_{v\ge0}\#\{\text{odd-bad }u\le X/2^v\}
 \le C R(\log B)\sum_{v\ge0}X/2^v
 =2CX R(\log B).
\]
The actual count sum is finite. The bound also covers its terms with
\(X/2^v<2\), whose counts vanish. This proves the raw estimate.

For the strict growing-threshold conclusion, fix any integer \(B\ge2\).
Outside a finite set, \(f(n)>B\). Thus failure to have an iterate
strictly below \(f(n)\) within the stated time is contained in the
fixed-\(B\) bad set plus that finite set. Take the upper density in
\(X\) first, then let \(B\to\infty\); \(R(\log B)\to0\).
For odd relative density the denominator is
\(\#\{n\le X:n\text{ odd}\}\sim X/2\), changing only a factor two.
Finally \(R(L)=o(L^{-c})\) for every \(0<c<b\), which proves the
last rate assertion. No pointwise assertion that every orbit is
bounded follows from these density estimates.
\end{proof}

The probability recurrence, rather than a sequence of independent
random orbits, is what permits the iteration. Its test event changes
with the accumulated integer budget \(K_j\); uniformity over all
landing events in Theorem~\ref{thm:Mazur-real-passage} is therefore
essential. The construction does not presume stationarity or
independence at successive scales. The two identities
\(q_j^\alpha=y_j\), \(q_j^{\alpha^2}=y_{j+1}\) identify every source
law used, and the pointwise inequalities identify exactly how an
actual path passes between the tests.

\begin{center}
\fbox{\begin{minipage}{0.93\linewidth}
\textbf{The measures and the two clocks.}
The exact source overlap controls real/floor transport.
Actual landing with \(\partial\) maps to the source's landing-with-\(1\);
only its \(\{1,\partial\}\) fibre loses information.
At each threshold \(q_j\), compare the two harmonic landing laws,
charge actual failure to land by \(n_{q_j}\), and add precisely
\(n_{q_j}\) to the clock. Their finite sum, plus the one uniform
top-block passage, is at most
\(\log X/[10\alpha(\alpha-1)\log2]\).
The inequality \(\log X\le\alpha\log n\) gives \(100\log_2 n\)
Syracuse steps. Expanding the valuation-labelled odd edges gives
\(301\log_2 n\) ordinary steps.
Proofs: Propositions~\ref{prop:Mazur-totalizer-exact}--%
\ref{prop:Mazur-critical-iteration} and
Theorem~\ref{thm:Mazur-critical-density}.
\end{minipage}}
\end{center}
